% Generated by code/assemble_manuscript.py; edit source fragments.

\section{A. The conormal step in full generality}

\subsection{A.1. Canonical cotangent kernels and ordinary crossings}

% A.1.1 from part_A.tex

Throughout this section $i:S\hookrightarrow Y$ is the inclusion of a reduced effective Cartier divisor into a complex manifold of dimension $n+1$, $j:S_{\rm reg}\hookrightarrow S$ is the inclusion of its regular locus, and $A$ is an invertible analytic $\OO_S$-module. We write
\[
 N^*=\OO_Y(-S)|_S,\qquad V=\Omega_Y^1|_S,\qquad
 T=\ker(\Omega_S^1\longrightarrow j_*\Omega^1_{S_{\rm reg}}),\qquad Q=\Omega_S^1/T.
\]
All maps between sheaves below are $\OO_S$-linear maps of coherent analytic sheaves, except for the displayed restriction into $j_*\Omega^1_{S_{\rm reg}}$. The latter is also $\OO_S$-linear; its kernel and image have the coherent descriptions proved below.

\begin{statement}{A.1.1}{Proved}{The conormal filtration and its exact section criterion}
The conormal map $d:N^*\to V$, defined by $[F]\mapsto dF|_S$ for a local defining equation $F$, is injective. With $K=\ker(V\to Q)$, there are canonical exact sequences
\begin{align}
0&\longrightarrow N^*\xrightarrow{d}K\longrightarrow T\longrightarrow0,
\label{A:filtration1}\\
0&\longrightarrow K\longrightarrow V\longrightarrow Q\longrightarrow0.
\label{A:filtration2}
\end{align}
Let $\partial_A:H^0(S,Q\otimes A)\to H^1(S,K\otimes A)$ and $\delta_A:H^0(S,T\otimes A)\to H^1(S,N^*\otimes A)$ be their connecting maps. Then
\[
 H^0(S,V\otimes A)=0
 \quad\Longleftrightarrow\quad
 H^0(S,K\otimes A)=0\ \hbox{and}\ \partial_A\hbox{ is injective}.
\tag{A.1}
\]
Under the two vanishings in Lead 1, $H^0(S,V\otimes A)$ is canonically $\ker\delta_A$.
\end{statement}
\begin{proof}
Locally $R=\C\{x_0,\ldots,x_n\}/(F)$ is reduced. Its total quotient ring is the product of the fraction fields of its irreducible components. At the generic point of every component, $S$ is smooth and $dF$ is nonzero. Consequently $a\,dF=0$ in $R^{n+1}$ implies that $a$ vanishes in each of those fields. The map from the reduced ring $R$ into their product is injective, so $a=0$. This proves injectivity in the usual right-exact conormal sequence. Taking the inverse image of $T\subset\Omega_S^1$ gives \eqref{A:filtration1} and \eqref{A:filtration2}. Tensoring by $A$ is exact because $A$ is locally free. The cohomology sequence of \eqref{A:filtration2} identifies the image of $H^0(VA)$ in $H^0(QA)$ with $\ker\partial_A$ and its kernel with $H^0(KA)$, proving (A.1). When $H^0(QA)=0$, it identifies $H^0(VA)$ with $H^0(KA)$. The first sequence then identifies this group with $\ker\delta_A$ when $H^0(N^*A)=0$.
\end{proof}


% A.1.2 from part_A.tex

\begin{statement}{A.1.2}{Proved}{The canonical normal-crossings formula}
For a reduced normal-crossings divisor, the map in (A.3) is an isomorphism. In a chart with $F=x_1\cdots x_r$, it is explicitly
\[
 \bigoplus_{i=1}^r R/(x_i)\longrightarrow K,
 \qquad (a_i)_i\longmapsto
 \sum_{i=1}^r a_i\prod_{j\ne i}x_j\,dx_i.
\tag{A.6}
\]
The tensor factor $N^*$ is trivialized here by $[F]$. At a double crossing $T\simeq R/(x,y)$; at a triple crossing,
\[
 T\simeq R^3/\bigl(R(1,1,1)+R(x,0,0)+R(0,y,0)+R(0,0,z)\bigr).
\tag{A.7}
\]
\end{statement}
\begin{proof}
On the normalization branch $x_i=0$, an ambient form maps to zero in the cotangent bundle of that branch precisely when all coefficients except that of $dx_i$ vanish. Thus a coefficient of $dx_i$ must vanish on every branch $x_j=0$ for $j\ne i$. In the convergent power-series ring it is divisible by $\prod_{j\ne i}x_j$. Its remaining factor is defined modulo $x_i$, which gives exactly (A.6). The coefficients of different $dx_i$ show injectivity. Coefficients of the coordinates not among $x_1,\ldots,x_r$ vanish on every branch and hence vanish in $R$. This proves surjectivity in every dimension.

For $r=1$ the map is $a\mapsto a\,dx$. For $r=2$ it is $(a,b)\mapsto ay\,dx+bx\,dy$, and quotienting by $dF=y\,dx+x\,dy$ leaves the generator $y\,dx=-x\,dy$, annihilated by $(x,y)$. For $r=3$ the branch generators are $yz\,dx,xz\,dy,xy\,dz$. Their sum is $dF$ and their respective annihilators are $(x),(y),(z)$, proving (A.7). Multiplication of $F$ by a unit and changes of coordinates were checked in A.2.1. Formula (A.6) is the restriction of the canonical fractional-ideal map, so it glues without a chosen branch ordering, equation, or coordinate system.
\end{proof}


\subsection{A.2. Inverse Jacobian ideals and conductor calculations}

% A.2.1 from part_A.tex

\begin{statement}{A.2.1}{Proved}{The Jacobian fractional ideal for every reduced hypersurface}
Let $\mathcal J\subset\OO_S$ be the Jacobian ideal, generated locally by $F_{x_0},\ldots,F_{x_n}$ modulo $F$. In the total quotient sheaf $\mathcal M_S$ define
\[
 \mathcal J^{-1}=\{q\in\mathcal M_S:q\mathcal J\subset\OO_S\}.
\]
There are canonical identifications
\[
 K\simeq N^*\otimes\mathcal J^{-1},\qquad
 T\simeq N^*\otimes(\mathcal J^{-1}/\OO_S),
\tag{A.2}
\]
under which the first map in \eqref{A:filtration1} is induced by $\OO_S\hookrightarrow\mathcal J^{-1}$. Thus (A.1), with this explicitly defined fractional ideal, is an exact necessary and sufficient criterion for every reduced hypersurface.
\end{statement}
\begin{proof}
The ideal $J$ contains a nonzerodivisor: it is nonzero at every minimal prime, and prime avoidance supplies an element outside all minimal primes. Every $R$-linear map $\varphi:J\to R$ is multiplication by a unique element of the total quotient ring. Indeed, for a nonzerodivisor $a\in J$, the identity $a\varphi(b)=b\varphi(a)$ for $b\in J$ gives the multiplier $q=\varphi(a)/a$. Hence $J^{-1}=\Hom_R(J,R)$ is coherent.

At the generic points, an element $\alpha\in K$ is uniquely $q\,dF$ for a tuple of meromorphic functions $q$. The coefficients of $\alpha$ belong to $R$ exactly when $qF_{x_i}\in R$ for every $i$, which is the condition $q\in J^{-1}$. Conversely such a vector restricts to zero in the cotangent space of every regular component, hence belongs to $K$. The map $[F]\otimes q\mapsto q\,dF|_S$ proves the first identification. Its quotient by $R\,dF$ gives the second. Replacing $F$ by $uF$, with $u$ a holomorphic unit, replaces $dF|_S$ by $u\,dF|_S$ and replaces the conormal generator by $u[F]$. Coordinate changes transform the derivative vector by an invertible Jacobian matrix. These are exactly the transitions in (A.2), so the identifications are canonical.
\end{proof}


% A.2.2 from part_A.tex

\begin{statement}{A.2.2}{Proved}{Normalization, conductor, and the excess kernel}
Let $\eta:\widetilde S\to S$ be the finite analytic normalization, $\mathcal B=\eta_*\OO_{\widetilde S}$, and $\mathfrak c=\Ann_{\OO_S}(\mathcal B/\OO_S)$ its conductor. There are canonical inclusions
\[
 \mathcal J\subset\mathfrak c,\qquad
 \mathcal B\subset\mathcal J^{-1},\qquad
 \eta_*\eta^*N^*\hookrightarrow K.
\tag{A.3}
\]
In a local ring $R$ with normalization $B$, put $D=\Hom_B(JB,B)$, viewed as a fractional ideal. Then
\[
 J^{-1}=\ker\left[D\longrightarrow
 \Hom_R\bigl(J,(B/\mathfrak c)/(R/\mathfrak c)\bigr)\right],
\quad q\longmapsto\bigl(b\mapsto qb\bmod R\bigr).
\tag{A.4}
\]
For any locally free $\OO_S$-module $E$, the conductor gives a second, independent exact section criterion:
\begin{multline}
H^0(S,E)=
\ker\left[H^0(\widetilde S,\eta^*E)\oplus H^0(D_c,E|_{D_c})\right.\\
\left.\longrightarrow H^0(\widetilde D_c,\eta^*E|_{\widetilde D_c})\right],
\qquad (a,b)\longmapsto a|_{\widetilde D_c}-\eta^*b,
\tag{A.5}
\end{multline}
where $D_c$ and $\widetilde D_c$ are defined by $\mathfrak c$ in $S$ and $\widetilde S$. In particular, the ambient group vanishes exactly when the map in (A.5), for $E=V\otimes A$, is injective.
\end{statement}
\begin{proof}
We prove $JB\subset R$ rather than assume a conductor formula. Choose a generic linear projection for which Weierstrass preparation writes the equation as a unit times a monic polynomial $P(z)$ over $A_0=\C\{x_1,\ldots,x_n\}$. Thus $R=A_0[z]/(P)$ is finite free over $A_0$ and generically separable. The functional taking the coefficient of $z^{m-1}$ in a remainder modulo $P$, where $m=\deg P$, generates $\Hom_{A_0}(R,A_0)$ as an $R$-module: the matrix of the pairing on $1,z,\ldots,z^{m-1}$ has zeros below the reverse diagonal and ones on that diagonal. Over the fraction field, Lagrange interpolation identifies that functional with $r\mapsto\operatorname{Tr}(r/P'(z))$. Therefore the trace dual is $P'(z)^{-1}R$.

For $b\in B$ and $r\in R$, the element $br$ is integral over the integrally closed ring $A_0$. Its trace belongs to $A_0$. It follows that $b$ lies in this trace dual, so $P'(z)b\in R$. Modulo $F$, the derivative in the chosen projection direction differs from $P'(z)$ by a unit. The permissible projection directions form a nonempty open set; choosing a linear basis of such directions proves $F_{x_i}B\subset R$ for every $i$. This proves $J\subset\mathfrak c$ and hence $B\subset J^{-1}$. The last inclusion of (A.3) is (A.2) and the projection formula.

The module $D$ consists exactly of the meromorphic $q$ satisfying $qJ\subset B$. Requiring that these products descend to $R$ is precisely the vanishing of their classes in $B/R=(B/\mathfrak c)/(R/\mathfrak c)$, proving (A.4). Finally the ring map
\[
 R\longrightarrow B\oplus R/\mathfrak c,
 \qquad r\longmapsto(r,r\bmod\mathfrak c)
\]
has image the kernel of $(b,\bar r)\mapsto b-\bar r\in B/\mathfrak c$. Surjectivity onto $B/\mathfrak c$ comes from the first summand. Tensoring this exact sequence by the free module corresponding to $E$ and taking global sections gives (A.5). All conductor restrictions are scheme-theoretic, so this argument includes nonreduced conductors.
\end{proof}

The fractional ideal appearing in (A.2) has an established interpretation. Granger--Schulze, Proposition~3.4 and Corollary~3.5 of \cite{GS2014}, identify $J^{-1}$ with the module of logarithmic residues and prove the inclusions in (A.3). The preceding proof gives the maps needed here directly. Their Theorems~1.1 and~1.3 also give the following exact extension of the normal-crossings case.


% A.2.3 from part_A.tex

\begin{statement}{A.2.3}{Known, with derivation}{The precise range of the normalization formula}
The canonical inclusion $\eta_*\eta^*N^*\hookrightarrow K$ is an isomorphism exactly when $S$ has normal crossings outside an analytic subset of codimension at least two in $S$.
\end{statement}
\begin{proof}
Locally a logarithmic one-form can be written as $\alpha/F$, with $\alpha$ holomorphic and $\alpha\wedge dF$ divisible by $F$. Restriction to $S$ is uniquely $q\,dF$ at its generic points; the holomorphic coefficients of $\alpha$ give $q\in J^{-1}$. Conversely, for $q\in J^{-1}$ choose holomorphic lifts of the coefficients $qF_{x_i}$. Their exterior product with $dF$ vanishes modulo $F$, so they define a logarithmic form with residue $q$. The kernel consists exactly of holomorphic forms. Thus the logarithmic residue module in \cite[Proposition~3.4]{GS2014} is exactly our $J^{-1}$, with no change of fractional-ideal embedding. Weakly holomorphic functions are exactly $B$. Theorems~1.1 and~1.3 there say that all logarithmic residues lie in $B$ exactly when the stated normal-crossings condition holds. Combining this with $B\subset J^{-1}$ and (A.2) proves the result.
\end{proof}


% A.2.4 from part_A.tex

\begin{statement}{A.2.4}{Proved}{Pinch points and cuspidal edges}
At the pinch point $F=x^2-y^2z$, with normalization $x=uv,y=u,z=v^2$, write $B=\C\{u,v\}$. Then
\[
 K=B\,dF,\qquad T\simeq B/R\simeq R/(x,y),
\quad \tau=2yz\,dx-2xz\,dy-xy\,dz
\tag{A.8}
\]
is a torsion generator. At the cuspidal edge $F=x^2-y^3$, with normalization $x=t^3,y=t^2$ and unchanged coordinate $z$, write $B=\C\{t,z\}$. Then
\[
 K=t^{-1}B\,dF,\qquad
 T\simeq t^{-1}B/R\simeq R/(x,y^2),
 \quad\tau=2y\,dx-3x\,dy.
\tag{A.9}
\]
The conductor is $(x,y)$ for the pinch point and $(x,y)$ for the cusp; in the cusp normalization it is $t^2B$. An ordinary surface, meaning a surface whose singularities are double normal crossings, triple normal crossings, and pinch points, satisfies canonically $K=\eta_*\eta^*N^*$ and $T=N^*\otimes(\eta_*\OO_{\widetilde S}/\OO_S)$. For the cuspidal edge the additional quotient $J^{-1}/B$ is generated by $t^{-1}$ over $\C\{z\}$.
\end{statement}
\begin{proof}
For the pinch point,
\[
 R=\C\{u,v^2,uv\}=\C\{v^2\}+uB,
 \quad J=(uv,uv^2,u^2)=u(u,v)B.
\]
The equality of the two descriptions of $J$ follows by separating even and odd powers of $v$. An element $q\in J^{-1}$ satisfies $(qu)(u,v)\subset B$. Since $B$ is a regular local ring, $\Hom_B((u,v),B)=B$: a meromorphic function whose products with $u$ and $v$ are holomorphic has no pole along any divisor, and its possible isolated singularity is removable by the holomorphic extension theorem. Hence $qu=h\in B$. The requirements $hv,hv^2\in R$, restricted to $u=0$, say respectively that $h(0,v)v$ and $h(0,v)v^2$ are even power series. The first requires $h(0,v)$ to be odd and the second to be even. Therefore $h(0,v)=0$, so $h\in uB$ and $q\in B$. The reverse inclusion follows either directly or from (A.3). Also $B/R$ is generated by $v$ and its annihilator is $(u,uv)=(y,x)$. Multiplying $dF$ by $v=x/y$ gives the displayed $\tau$.

For the cusp, $R=\C\{t^2,t^3,z\}$ and $J=(t^3,t^4)$. A meromorphic $q$ with $qt^3,qt^4\in R$ has no poles away from $t=0$; writing its Laurent expansion in $t$ shows that the lowest allowed exponent is $-1$. Indeed an exponent $-3$ produces a forbidden $t$ term after multiplication by $t^4$, an exponent $-2$ produces one after multiplication by $t^3$, and smaller exponents produce a pole. Conversely every $t^{-1}B$ has both products in $R$. Thus $J^{-1}=t^{-1}B$. Modulo $R$, the quotient has the $\C\{z\}$-basis $t^{-1},t$; multiplication by $y$ sends the first basis vector to the second, while $x$ and $y^2$ kill both. This proves the cyclic presentation and (A.9). The conductor calculations follow directly from the displayed descriptions of $R\subset B$. Formula (A.2) and the three local computations prove the ordinary-surface assertion by gluing.
\end{proof}


% A.2.5 from part_A.tex

\begin{statement}{A.2.5}{Proved}{Vertical residue poles at a tacnode--cusp intersection}
For $F=x^2-y^4z^3$, normalization is $x=u^2v^3$, $y=u$, $z=v^2$. Put $A_0=\C\{u,v^2\}$ and $B=\C\{u,v\}$. Then
\[
 \begin{aligned}
 R&=A_0\oplus u^2v^3A_0,&\mathfrak c&=u^2v^2B,\\
 J^{-1}&=A_0\oplus(uv)^{-1}A_0,&
 J^{-1}/B&\simeq A_0/(uv^2).
 \end{aligned}
 \tag{A.9a}
\]
In a trivialization of $N^*$ the torsion module is
\[
 T\simeq R/(x,y^3z^2),\qquad
 \tau=2yz\,dx-4xz\,dy-3xy\,dz.
 \tag{A.9b}
\]
In particular $v^{-1}\in J^{-1}\setminus B$. For the ruling with base coordinate $v$, this is a residue constant on its generic fibres with a pole on the vertical divisor $v=0$, although the horizontal conductor $u=0$ has coefficient two.
\end{statement}
\begin{proof}
The element $v=x/(y^2z)$ belongs to the fraction field and satisfies $v^2=z$. Adjoining it gives the regular ring $B$, which is therefore the finite normalization. Separating even and odd powers of $v$ gives the displayed description of $R$. An even monomial $u^av^b$ belongs to the conductor exactly when $a\ge2$ and $b\ge2$; an odd monomial already in $R$ has $a\ge2$, $b\ge3$ and is then in the conductor. Thus $\mathfrak c=u^2v^2B$.

The normalized Jacobian generators, with irrelevant nonzero constants removed, are $u^2v^3,u^3v^6,u^4v^4$. If $qJ\subset R$, the first product shows that $q$ has no other polar divisor, with possible orders at most two along $u=0$ and three along $v=0$. A term of order $v^{-3}$ would produce a forbidden $v^1$ term in $qu^4v^4$; a term of order $v^{-2}$ would produce one in $qu^2v^3$. Both coefficients must vanish, so $q$ has order at least $-1$ in $v$. A coefficient $a(v)u^{-2}$ would require both $a(v)v^3$ and $a(v)v^6$ to be even Laurent series, from the $u^0$ and $u^1$ coefficients of the first two products. The first requirement makes $a$ odd and the second makes it even, so $a=0$. Thus $q=t/(uv)$ with $t\in B$.

Write $t=a(u,v^2)+v b(u,v^2)$. The first Jacobian product equals $uv^2t$ and belongs to $R$ exactly when $b$ is divisible by $u$. The other two products automatically belong to $R$ under this condition. This gives $J^{-1}=A_0\oplus(uv)^{-1}A_0$. The odd part of $B$ is $(uv)^{-1}(uv^2)A_0$, giving $J^{-1}/B$ in (A.9a). The odd part of $R$ is $(uv)^{-1}u^3v^4A_0$, so $J^{-1}/R$ is generated by $(uv)^{-1}$ with annihilator $(x,y^3z^2)$. Multiplication of $dF$ by $(uv)^{-1}$ gives (A.9b).

The particular vertical pole has the completely regular gradient products
\[
 v^{-1}dF=2y^2z\,dx-4xyz\,dy-3xy^2\,dz.
\]
Pulling this expression to $B$ gives zero. It is a nonzero ambient-kernel vector, with multiplier $v^{-1}$, rather than a section of the normalization submodule. The exact identities and finite monomial checks are saved in \path{code/A_vertical_pole.py}; the Laurent argument above establishes the full convergent module. This local example does not assert the existence of a compact CDP fibration realizing that configuration.
\end{proof}

For ordinary surfaces the exact ambient-vanishing criterion is (A.1) with $K\otimes A=\eta_*\eta^*(N^*\otimes A)$. For surfaces with cuspidal-edge charts it is (A.1) with the locally specified sheaf $N^*\otimes t^{-1}B$. These specify the first group and the target of the connecting homomorphism, including the extra polar term at the cusp. Under $H^0(QA)=0$ they reduce respectively to the vanishing of these explicitly computed kernel sections. The general criterion (A.4)--(A.5) retains the full conductor and applies without a singularity-type restriction. The exact polynomial identities and exponent computations are recorded in \path{code/A_local_algebra.py} and its output.


% A.2.6 from part_A.tex

\begin{statement}{A.2.6}{Proved}{A transverse cusp whose critical-curve equation is only a symbolic square}
Put
\[
 \begin{aligned}
 I&=(y^2-xz,\ yz^2-x^4,\ z^3-x^3y)\subset\C\{x,y,z\},\\
 F&=4x^7-15x^3yz^2+10x^2y^3z-3xy^5+4z^5.
 \end{aligned}                                                   \tag{A.NL.1}
\]
The reduced critical locus of the reduced hypersurface $F=0$ is the
irreducible curve $D$ parametrized by $(x,y,z)=(t^5,t^6,t^7)$.
Its transverse type at every point with $t\ne0$ is the ordinary cusp
$A_2$. The curve is not a local complete intersection at the origin,
and
\[
             F\in I^{(2)}\setminus I^2.                         \tag{A.NL.2}
\]
Thus the ordinary-square membership used to define the conormal
quadratic jet in A.5.5 does not extend to an arbitrary reduced
critical curve, even with generic transverse type $A_2$.
This is a local example, without an asserted compact CDP realization.
\end{statement}
\begin{proof}
Write the three displayed generators of $I$ as $g_1,g_2,g_3$.
The quotient is a finite $\C\{x\}$-module generated by
$1,y,z,yz,z^2$: their products by $y$ and $z$ reduce to this list
using $y^2=xz$, $yz^2=x^4$, and $z^3=x^3y$.
For example these identities also give $y^5=x^6$ and $z^5=x^7$,
which make the analytic quotient finite over $\C\{x\}$.
The parametrization sends the five generators to
$1,t^6,t^7,t^{13},t^{14}$. Their exponents have distinct residues
modulo five, so they are linearly independent over $\C\{t^5\}$.
Consequently this actual ring homomorphism is injective and has
image $\C\{t^5,t^6,t^7\}$; it proves the stated normalization and
the primeness of $I$. The three generators are minimal. Indeed their
initial terms have degrees two, three, three; after the degree-two
coefficient vanishes, the cubic terms $yz^2$ and $z^3$ are independent
modulo the span of $xg_1,yg_1,zg_1$. A height-two ideal with three
minimal generators is not a complete intersection.

The following polynomial identities determine the critical locus:
\[
 \begin{aligned}
 F_y&=-15xg_1^2,\\
 F_x&=(17xyz-3y^3)g_1-28x^2g_2,\\
 F_z&=10x^2yg_1+20zg_3.
 \end{aligned}                                                   \tag{A.NL.3}
\]
At $x=0$, the equations $F_x=F_z=0$ give $y=z=0$.
At $x\ne0$, $F_y=0$ gives $g_1=0$, and then $F_x=0$
gives $g_2=0$. The value $z=0$ is impossible there, so $F_z=0$
also gives $g_3=0$. Conversely all three partials vanish on $D$.
Weighted Euler's identity, with weights $(5,6,7)$ and degree $35$,
puts $F$ in its Jacobian ideal. The critical locus is therefore
exactly $D$ as a reduced set. A repeated factor of $F$ would put
a two-dimensional component in this critical locus, which proves
that $F$ is reduced.

The Hessian along the normalization has the exact factorization
\[
 \operatorname{Hess}(F)|_D
 =2t^{21}
 \begin{pmatrix}7t^2\\0\\-5\end{pmatrix}
 \begin{pmatrix}7t^2&0&-5\end{pmatrix},
 \qquad \frac{1}{6}F_{yyy}|_D=-20t^{17}.                       \tag{A.NL.4}
\]
For $t\ne0$, the plane $x=t^5$ is transverse to $D$.
The restricted Hessian has rank one and its radical is the
$y$ direction; its cubic coefficient in that direction is nonzero.
The holomorphic splitting lemma and a local cube root of a unit
give the transverse cusp equation. At the generic point of $D$,
its conormal is locally free and injects into ambient differentials.
Since $F$ and $dF$ vanish there, $F\in I^2\C\{x,y,z\}_I$,
which is precisely $F\in I^{(2)}$.

To decide ordinary-square membership, use the positive weights
$(5,6,7)$. The three generators have weights $12,20,21$.
Among products of two generators, only $g_1^2$ can contribute
to weight $35$: it needs a coefficient of weight $11$, whose
space is spanned by $xy$. Products $g_1g_2$ and $g_1g_3$ would
need weights $3$ and $2$, which do not occur; the other products
already have greater weight. Thus the weight-$35$ part of $I^2$
is exactly $\C\,xyg_1^2$. Its $x^7$ coefficient is zero, whereas
that of $F$ is four. Taking the weight-$35$ component of a
convergent series identity is legitimate because all weights are
positive. This proves $F\notin I^2$ in the analytic local ring.

For comparison of the curve invariants, the normalization conductor
is $t^{10}\C\{t\}$: the semigroup gaps are $1,2,3,4,8,9$.
The differential image in $\C\{t\}\,dt$ is $t^4\C\{t\}\,dt$.
The smallest valuation of a two-by-two minor of the displayed
generators' Jacobian matrix is $19$. Thus the normalized conormal
determinant correction is $19-4=15$, while the conductor exponent
is $10$; the complete-intersection adjunction comparison cannot
identify these two quantities here.

The symbolic-square correction is visible directly on the free
normalized conormal. Set $a=dg_1|_D$ and
$b=dg_2|_D+2t^8a$. Then
\[
 \begin{aligned}
 a&=(-t^7,2t^6,-t^5),& b&=(-6t^{15},5t^{14},0),\\
 dg_3|_D&=-3t^9a+tb,&
 q&=t^{-5}(5t^8a-2b)^2.
 \end{aligned}                                                   \tag{A.NL.5}
\]
The first two rows generate the entire normalized conormal image,
and $q$ is the quadratic term, namely half the Hessian.
If $A,B$ are the dual basis, its radical frame $2A+5t^8B$
is represented in the ambient normal quotient by $t^{-6}\partial_y$.
Its intrinsic cubic is therefore $-20t^{-1}$. These are actual
quadratic and cubic poles of orders five and one. The corresponding
combined invariant $q^3k^2$ has pole order $17$; no bound for such
poles on all non-lci curves is deduced from this one computation.
All substitutions, polynomial identities, weighted coefficient
spaces, conormal comparisons, and Jacobian minors are checked
exactly in \path{code/A_non_lci_cusp.py}.
\end{proof}


\subsection{A.3. Full fibres and their multiplicities}

% A.3.1 from part_A.tex

\begin{statement}{A.3.1}{Proved}{All thickenings and their lifting maps}
Let $S$ be any effective Cartier divisor in a complex manifold, let $E$ be any locally free sheaf on that manifold, and let $E_k=E|_{kS}$. For each $k\ge2$ there is a canonical exact sequence
\[
0\longrightarrow (E|_S)\otimes(N^*)^{k-1}
 \longrightarrow E_k\longrightarrow E_{k-1}\longrightarrow0.
\tag{A.10}
\]
Write $G_{k-1}=(E|_S)\otimes(N^*)^{k-1}$ and let
$\epsilon_k:H^0((k-1)S,E_{k-1})\to H^1(S,G_{k-1})$
be the connecting map. Then
\[
0\longrightarrow H^0(S,G_{k-1})\longrightarrow H^0(kS,E_k)
 \longrightarrow\ker\epsilon_k\longrightarrow0.
\tag{A.11}
\]
In particular $H^0(kS,E_k)=0$ exactly when $H^0(S,G_{k-1})=0$ and $\epsilon_k$ is injective. Starting with $H^0(S,E|_S)$, equations (A.10)--(A.11) compute the group on $mS$ through all its extension maps. The vanishing of every $H^0(S,G_j)$, $0\le j<m$, is a sufficient criterion. These statements apply to $E=\Omega_Y^1\otimes B$ for every line bundle $B$, and include smooth, normal-crossings, and all other reduced supports.
\end{statement}
\begin{proof}
With $I=\OO_Y(-S)$, the sequence is the tensor product of $E$ with
$0\to I^{k-1}/I^k\to\OO_Y/I^k\to\OO_Y/I^{k-1}\to0$.
The Cartier property identifies $I^{k-1}/I^k$ with $(N^*)^{k-1}$ on $S$; local freeness makes the tensor product exact. Taking global sections proves (A.11) and its equivalence. Induction on $k$ proves the sufficient criterion. No splitting of the thickening, cyclic base change, or assumption about the normal bundle has been used.
\end{proof}


% A.3.2 from part_A.tex

\begin{statement}{A.3.2}{Refuted, with computed replacement}{Vanishing on the support alone does not imply vanishing on a multiple fibre}
There is a compact smooth complex threefold $g:Y\to\PP^1$ with connected general fibre an abelian surface and a smooth fibre support $S$ of multiplicity $m$, for every $m\ge3$, and a line bundle $B$ on $Y$, such that
\[
H^0(S,(\Omega_Y^1\otimes B)|_S)=0,
\qquad h^0(mS,(\Omega_Y^1\otimes B)|_{mS})=3.
\tag{A.12}
\]
Its conormal has exact order $m$. The replacement is (A.10)--(A.11), with the graded twists retained.
\end{statement}
\begin{proof}
Take an abelian surface $F$ and a translation $a:F\to F$ of exact order $m$. Let $\zeta$ be a primitive $m$th root of unity. The action generated by $(p,z)\mapsto(a(p),\zeta z)$ on $F\times\PP^1$ is free, so the quotient $Y$ is a smooth compact threefold and the quotient map is finite etale. The invariant map $t=z^m$ gives $g:Y\to\PP^1$. Away from $0$ and $\infty$ its fibre is isomorphic to $F$: the $m$ lifts of a base point are permuted transitively and freely. At $0$ the fibre is $mS$, where $S=F/\langle a\rangle$. Set $B=\OO_Y(-S)$.

The connected cyclic cover $F\to S$ identifies $N^*=B|_S$ with the character line of weight one. Its order is exactly $m$: otherwise a nontrivial proper character power would have a nonzero invariant trivialization on the connected compact cover, whose holomorphic functions are constant. Pullback to $F$ splits the ambient cotangent bundle into two invariant tangent forms and the normal form $dz$ of weight one. Consequently
\[
 (\Omega_Y^1\otimes B)|_S\simeq (N^*)^{\oplus2}\oplus(N^*)^2.
\]
The nontrivial flat lines on the abelian surface $S$ have no holomorphic sections. For completeness, their pullbacks to $F$ are trivial; a section would be a constant function transforming by a nontrivial character, which is zero. Since $m\ge3$, both characters in the preceding splitting are nontrivial, proving the first equality.

The inverse image of $mS$ is the ringed space with structure sheaf $\OO_F\otimes\C[z]/(z^m)$. Write $e$ for the local generator of the pulled-back line $B$; it has weight one. All global sections on this inverse image are linear combinations of $z^k\alpha_1e,z^k\alpha_2e,z^kdz\,e$ for $0\le k<m$, where $\alpha_1,\alpha_2$ are a basis of invariant one-forms on $F$. The first two have weight $k+1$ and the third weight $k+2$. The invariant monomials are exactly
$z^{m-1}\alpha_1e,z^{m-1}\alpha_2e,z^{m-2}dz\,e$.
Taking invariants is exact over $\C$, by averaging over the finite group. These three monomials therefore give the claimed basis and prove the second equality. Translations act identically on $H^q(F,\OO_F)$, whose dimensions are $\binom{2}{q}$. The same calculation gives all fibre groups: $h^q(S,E|_S)=(0,0,0)$ and $h^q(mS,E|_{mS})=(3,6,3)$ for $q=0,1,2$. For completeness, on the compact threefold itself the two tangent summands pull back to $\OO_F\boxtimes\OO_{\PP^1}(-1)$ and have no cohomology. The normal summand is $\OO_F\boxtimes\Omega^1_{\PP^1}(-\{0\})$. Its first cohomology on $\PP^1$ has the Cech basis $dz/z,dz$, of weights zero and one, respectively. Indeed sections near zero are multiples of $z\,dz$, and sections near infinity are multiples of $z^{-2}dz$, leaving exactly those two Laurent classes. Its invariant part is one-dimensional. Kunneth and finite-cover descent therefore give $h^q(Y,E)=(0,1,2,1)$ for $q=0,1,2,3$.
\end{proof}

For this example the complete dimensions on $kS$, $1\le k\le m$, are $0$ for $k\le m-2$, $1$ for $k=m-1$, and $3$ for $k=m$; the same character calculation proves this. More generally the character count with $B=\OO_Y(-bS)$ is
\[
\begin{aligned}
h^0(kS,(\Omega_Y^1\otimes B)|_{kS})
&=2\#\{0\le j<k:j+b\equiv0\pmod m\}\\
&\quad+\#\{0\le j<k:j+b+1\equiv0\pmod m\}.
\end{aligned}
\tag{A.13}
\]
The finite counts for $3\le m\le12$ and all characters are checked in \path{code/A_multiple_fibres.py}.


% A.3.3 from part_A.tex

\begin{statement}{A.3.3}{Proved}{The precise finite-translation step in Remarks 7.2 and 7.4}
For a fixed multiple fibre $mS$, suppose that the set of line bundles $B$ for which $H^0(S,(\Omega_Y^1\otimes B)|_S)=0$ contains a dense open subset $U$ of the relevant Picard component. The intersection
\[
\bigcap_{j=0}^{m-1}\{B:B\otimes\OO_Y(-jS)\in U\}
\tag{A.14}
\]
is a dense open subset on which the full-fibre ambient group vanishes, whenever the displayed translations preserve that component. The same assertion for a finite union of components follows by using the appropriate dense open subset in each translated component.
\end{statement}
\begin{proof}
Tensoring by the fixed line bundle $\OO_Y(-jS)$ is a biholomorphic translation between Picard components. Its inverse image of a dense open set is dense and open. The finite intersection retains these properties. Its defining condition is exactly the vanishing of every graded group in A.3.1, which proves the result.
\end{proof}

This finite-translation argument supplies the thickening step intended in Remark~7.2 once ambient vanishing on the smooth support is established. On a smooth support, the conormal sequence has $Q=\Omega_S^1$ and $K=N^*$, so the two requisite support vanishings do establish that ambient vanishing. Remark~7.4 requires the same argument on a singular support; a statement solely about $Q=\widetilde\Omega_S^1$ leaves the kernel in A.1.1 to be checked. Thus the thickening operation itself has a complete exact treatment, while the cited singular-support assertion cannot supply its missing kernel hypothesis. Example A.3.2 disproves the inference from a single support vanishing even with torsion conormal; it does not refute a generic-vanishing assertion already established for all the finitely translated twists.


% A.3.4 from part_A.tex

\begin{statement}{A.3.4}{Proved}{The smooth-multiple assertion of Remark 7.2}
Under the hypotheses of CDP20 Proposition~7.3, let $mS$ be a full fibre whose reduced support $S$ is smooth. For every real irrational transgression parameter $r$ of D.1.9,
\[
 H^0\bigl(mS,(\Omega_Y^1\otimes L_r)|_{mS}\bigr)=0.
 \tag{A.15a}
\]
The same simultaneous vanishing on all such fibres holds on the complement of a proper complex analytic subset of $\Pic(Y)$. The assertion includes the entire scheme-theoretic fibre, with every nilpotent layer.
\end{statement}
\begin{proof}
Choose a disc $\Delta$ about the fibre value and write $U=g^{-1}(\Delta)$. Form the base change $t=s^m$ and normalize it, obtaining $p:U'\to U$ and $g':U'\to\Delta'$. This map $p$ is finite étale of degree $m$. Indeed, near a point of the smooth support write $g=u z^m$ with $u$ a holomorphic unit. The normalization is given by
\[
 w^m=u,\qquad s=wz.
\]
The derivative $m w^{m-1}$ is invertible, so these equations give an unramified cover of the original smooth chart. They cover the normalized base change, since they coincide away from $z=0$ and are normal. The local constructions glue by uniqueness of normalization. The divisor $S'=(g')^{-1}(0)$ is smooth and reduced, $g'$ is a submersion near it, and $p^{-1}(mS)=mS'$ as Cartier divisors.

Pull back the unitary local system (A.22) on $U$ to $U'$; its associated holomorphic line is $A'=p^*L_r$. Every nearby smooth fibre of $g'$ maps isomorphically onto the corresponding nearby fibre of $g$. Its character is therefore the irrational exponential of a nonzero integral class. The smooth-fibre calculation in A.4.9 gives zero cohomology in degrees zero and one for that local system. Proper smooth local triviality at $S'$ transports the same cohomology statement to $S'$. If a component presentation of the cover is used, the argument applies on each component and the groups add.

A holomorphic section of a unitary-flat line on a compact smooth complex manifold is parallel: its squared norm is plurisubharmonic and hence constant by the maximum principle. Thus $H^0(S',A')=0$ follows from the zero degree-zero local-system cohomology. A.4.7 gives $H^0(S',\Omega_{S'}^1\otimes A')=0$ from degree-one vanishing. The conormal of the reduced full fibre $S'$ is trivial, with global generator $dg'$. Its conormal sequence, tensored by $A'$, now gives
$H^0(S',(\Omega_{U'}^1\otimes A')|_{S'})=0$.

All the graded factors for $mS'$ in A.3.1 equal this same restricted bundle, because the conormal is trivial. Induction in (A.10) proves $H^0(mS',(\Omega_{U'}^1\otimes A')|_{mS'})=0$. Since $p$ is étale, $p^*\Omega_U^1=\Omega_{U'}^1$, and finite faithfully flat pullback injects sections on $mS$ into sections on $mS'$. One can verify that injection locally by tensoring the inclusion of a module into its faithfully flat scalar extension, or by the trace divided by $m$. This proves (A.15a).

There are finitely many non-submersion values. For each of these fixed full fibres the coherent ambient bundle varies in the holomorphic Picard family constructed in A.4.9, and proper semicontinuity makes its nonvanishing locus analytic. Each locus is proper because it excludes every $L_r$ with $r$ real irrational. Their finite union is proper, proving the simultaneous generic assertion. This proves the smooth-multiple conclusion of Remark~7.2. For a singular support the identical root cover retains its singularities and the kernel $N^*\otimes\mathcal J^{-1}$; the present proof does not erase that term in Remark~7.4.
\end{proof}


% A.3.5 from part_A.tex

\begin{statement}{A.3.5}{Proved}{Every reducible full fibre has generic ambient vanishing}
Let $Y$ be a smooth compact complex threefold with $H^2(Y,\Q)=0$, and let $F=\sum_i m_iS_i$ be an effective Cartier divisor with connected reducible support. Every connected component of $\Pic(Y)$ contains a dense open subset, with proper closed analytic complement, on which
\[
 H^0(F,(\Omega_Y^1\otimes L)|_F)=0.
 \tag{A.15b}
\]
The multiplicities may be unequal. For a proper connected fibration over a compact curve, the vanishing is simultaneous on all reducible full fibres after intersecting finitely many such open sets. Consequently a full fibre with ambient sections for every Picard twist has irreducible support.
\end{statement}
\begin{proof}
Every component $S_i$ meets a different component $S_j$, because the support is connected and reducible. The canonical section of $\OO_Y(S_j)$ restricts to a nonzero section on $S_i$. It is a nonzerodivisor in every local ring of $S_i$: otherwise one local branch of $S_i$ would be contained in $S_j$, and analytic continuation on the global irreducible component would give $S_i=S_j$. Its zero scheme is therefore a nonzero effective Cartier divisor $D$ on $S_i$.

For any fixed locally free sheaf $E$ on the reduced irreducible compact surface $S_i$, the injections
\[
 H^0(S_i,E(-nD))\hookrightarrow H^0(S_i,E)
\]
form a decreasing sequence of finite-dimensional subspaces. Their intersection is zero. At a point of $D$, a section in the intersection has its germ in every power of the Cartier ideal; the Krull intersection theorem for the Noetherian analytic local ring makes this germ zero. The identity theorem on the reduced irreducible surface, applied in local trivializations of $E$, then makes the global section zero. A descending sequence of finite-dimensional subspaces stabilizes, so $H^0(S_i,E(-nD))=0$ for all sufficiently large $n$.

Fix a Picard component and a representative $L_0$ in it. The integral class $c_1(\OO_Y(S_j))$ is torsion because $H^2(Y,\Q)=0$; choose a positive integer $q$ killing it. Then $\OO_Y(-qnS_j)$ belongs to $\Pic^0(Y)$ for every $n$. Apply the preceding argument to $E=(\Omega_Y^1\otimes L_0)|_{S_i}$. It supplies, in the chosen Picard component, a twist $L_0\otimes\OO_Y(-qnS_j)$ with zero ambient group on $S_i$. Local holomorphic Picard families and proper semicontinuity show that the nonvanishing locus on that component is a proper closed analytic subset. Its complement is dense and open. The same argument works in every Picard component.

Choose an ordered list $A_1,\ldots,A_N$ of the prime divisors $S_i$, with $S_i$ repeated $m_i$ times, and put $F_0=0$, $F_k=A_1+\cdots+A_k$. The exact Cartier sequences, tensored with the locally free ambient bundle, are
\[
\begin{split}
0\longrightarrow
 (\Omega_Y^1\otimes L\otimes\OO_Y(-F_{k-1}))|_{A_k}
 \longrightarrow (\Omega_Y^1\otimes L)|_{F_k}\\
 \longrightarrow(\Omega_Y^1\otimes L)|_{F_{k-1}}
 \longrightarrow0.
\end{split}
\]
This follows locally from multiplication by an equation of $F_{k-1}$, and remains exact when a component occurs repeatedly. For each $k$ take the dense open set just proved on $A_k$ and translate it by the fixed global line $\OO_Y(-F_{k-1})$, using the appropriate Picard component. Their finite intersection is dense open. Every displayed first term has no sections there, so induction from $F_0=0$ proves (A.15b) on the full scheme $F_N=F$.

A proper holomorphic map from a smooth compact manifold to a smooth curve has finitely many critical values, since its proper analytic critical-value set does not contain a general point. Its smooth connected fibres are irreducible. Thus there are finitely many reducible fibres, and a further finite intersection gives the asserted simultaneous vanishing. This proves the final consequence and resolves the unequal-multiplicity transfer in this case without any assertion about torsion-free differentials on individual components.
\end{proof}


\subsection{A.4. The torsion-free cotangent assertion}

% A.4.1 from part_A.tex

\begin{statement}{A.4.1}{Proved}{Both fibre groups on the specified $W$}
For the reduced normal-crossings fibre $W$ with normalization $\eta:\operatorname{dP}_6\to W$ in the specified threefold $X$, and every $L\in\Pic(X)$, put $A=(K_X\otimes L^{-1})|_W$. Then
\[
H^0(W,Q\otimes A)=0,
\qquad H^0(W,\Omega_X^1|_W\otimes A)\simeq\C.
\tag{A.15}
\]
The same two conclusions hold with $A=B|_W$ for every $B\in\Pic(X)$.
\end{statement}
\begin{proof}
The only global topological input is $H^2(X,\Z)=0$, and the only identification of the fibre used is its specified normalization as the blowup of $\PP^2$ at three points. A blowup at a smooth point has $\pi_*\OO=\OO$ and $R^1\pi_*\OO=0$. The latter follows from the theorem on formal functions: successive exceptional-neighborhood quotients are $\OO_{\PP^1}(k)$, $k\ge0$, whose first cohomology is zero. Thus $H^1(\operatorname{dP}_6,\OO)=H^1(\PP^2,\OO)=0$. The exponential sequence makes $c_1:\Pic(\operatorname{dP}_6)\to H^2(\operatorname{dP}_6,\Z)$ injective. The pullback of every line bundle on $X$ has zero first Chern class and is therefore trivial.

The differential of normalization gives $\Omega_W^1\to\eta_*\Omega^1_{\operatorname{dP}_6}$. Its target is torsion-free as an $\OO_W$-module, so it factors through $Q$. The resulting map is injective because it is an isomorphism on the dense regular locus and its source is torsion-free. A holomorphic one-form on the blowup restricts away from the exceptional curves to a one-form on $\PP^2$ minus three points. Its coefficients extend over these points by Hartogs' theorem, and $H^0(\PP^2,\Omega^1_{\PP^2})=0$ by the Euler sequence. Hence $H^0(\operatorname{dP}_6,\Omega^1)=0$. The projection formula and the triviality of $\eta^*A$ prove the first equality in (A.15).

Because $W$ is the full reduced fibre, its conormal is trivial. A.1.2 identifies $K\otimes A$ with $\eta_*\eta^*A=\eta_*\OO_{\operatorname{dP}_6}$. The normalization is connected and compact, so its global holomorphic functions are $\C$. The cohomology sequence of A.1.1 and the first equality in (A.15) identify this group with the ambient group. The same argument applies to every $B$.
\end{proof}


% A.4.2 from part_A.tex

\begin{statement}{A.4.2}{Refuted}{A specific implication in the printed proof of Proposition 7.3}
The implication that a ruled minimal surface $S_0$ with a one-parameter family of numerically trivial line bundles $M_t$ and nonzero sections of $\Omega^1_{S_0}\otimes M_t$ must be ruled over a curve of genus at least two is false when the family $M_t$ is allowed to be constant, as the printed proof explicitly allows.
\end{statement}
\begin{proof}
Take $S_0=E\times\PP^1$ for an elliptic curve $E$ and take $M_t=\OO_{S_0}$ for every parameter. The product cotangent decomposition gives
\[
 h^q(S_0,\Omega^1_{S_0})=(1,2,1),\qquad q=0,1,2.
\]
Indeed the $E$ cotangent summand contributes $(1,1,0)$, while the $\PP^1$ cotangent summand contributes $(0,1,1)$ by the Kunneth formula and $h^1(\PP^1,\OO(-2))=1$. A nonzero invariant form on $E$ provides the asserted sections for every parameter, whereas the ruling has genus one. This refutes the stated implication. It is a counterexample to that implication with its explicit data; it supplies no threefold satisfying all of Proposition~7.3's ambient hypotheses.
\end{proof}


% A.4.3 from part_A.tex

\begin{statement}{A.4.3}{Proved}{The omitted genus-one case for the standard rank-one fibre}
Let $E$ be an elliptic curve, let $P\in\Pic(E)$, and let $T=\PP_E(\OO_E\oplus P)$ with its zero and infinity sections $D_0,D_\infty$. Glue $D_0$ to $D_\infty$ by a translation of $E$ to form the compact normal-crossings surface $S$, and denote its normalization by $\eta:T\to S$. Every line bundle $A$ on $S$ whose pullback $\eta^*A$ is topologically trivial satisfies
\[
H^0(S,\widetilde\Omega_S^1\otimes A)=
\begin{cases}\C,&A\simeq\OO_S,\\0,&A\not\simeq\OO_S.\end{cases}
\tag{A.16}
\]
\end{statement}
\begin{proof}
For the ruling $p:T\to E$, every topologically trivial line bundle on $T$ is $p^*P_A$ for a unique $P_A\in\Pic^0(E)$. Indeed its degree on every ruling fibre is zero, so the adjunction map from the pullback of its rank-one direct image is an isomorphism, checked on each fibre by $H^0(\PP^1,\OO)=\C$ and $H^1(\PP^1,\OO)=0$. The line on $E$ has degree zero because its first Chern class pulls back to zero and $p$ has a section.

The relative cotangent sequence and restriction to ruling fibres give
\[
 H^0(T,\Omega_T^1\otimes p^*P_A)
 =p^*H^0(E,\Omega_E^1\otimes P_A).
\]
The relative quotient has no sections because its restriction to every fibre is $\OO_{\PP^1}(-2)$. Since $\Omega_E^1\simeq\OO_E$, the displayed group vanishes for nontrivial $P_A$. This proves (A.16) in that case by the injection into normalization differentials.

For $P_A=\OO_E$, choose a trivialization of $\eta^*A$. The descent datum for $A$ is a nonzero constant $\lambda$ identifying the two restrictions to $D_0,D_\infty$, because these are trivial lines on the compact connected curve $E$. Multiplying the chosen trivialization by a constant changes neither $\lambda$ nor its meaning. The local calculation for $xy=0$ identifies the image of $\widetilde\Omega_S^1$ in $\eta_*\Omega_T^1$ with pairs of forms whose tangential restrictions to the double curve agree. The unique normalization form is a scalar multiple of $p^*\alpha$, where $\alpha$ is a nonzero invariant form on $E$. Translation pulls $\alpha$ back to itself. Twisted descent is therefore the equation $\alpha=\lambda\alpha$, so a nonzero section exists exactly when $\lambda=1$. That equality also says precisely that the trivialization descends to $A\simeq\OO_S$. This proves both dimensions.
\end{proof}


% A.4.4 from part_A.tex

\begin{statement}{A.4.4}{Proved}{Canonical triviality on an irreducible reduced torus fibre}
Let $g:Y\to C$ be a proper holomorphic map from a connected smooth compact threefold onto a smooth compact curve, with connected fibres and general fibre a complex two-torus. Every irreducible reduced full fibre $S=g^{-1}(c)$ has $\omega_S\simeq\OO_S$. When $H^2(Y,\Q)=0$, it therefore satisfies $h^q(S,\OO_S)=(1,2,1)$.
\end{statement}
\begin{proof}
The coherent sheaf $g_*\omega_Y$ is torsion-free: a nonzero function on an open subset of the base cannot annihilate a holomorphic section of a line bundle on the dominating smooth total space. It has generic rank one because the smooth torus fibre has one holomorphic top form; after restricting to a base open set where Grauert base change holds, this identifies its rank. On the smooth curve it is consequently a line bundle $M$. The evaluation morphism $g^*M\to\omega_Y$ is generically nonzero, so its zero scheme is an effective Cartier divisor $D$ and $\omega_Y\simeq g^*M\otimes\OO_Y(D)$. No component of $D$ dominates the base: on a general smooth torus fibre its defining section is a nonzero holomorphic top form and is nowhere zero. Thus $D$ is vertical.

All vertical components meeting the specified irreducible full fibre are that fibre itself. Consequently $\OO_Y(D)|_S$ is a power of $\OO_Y(S)|_S$, and the latter is trivial because a coordinate on $C$ defines the full reduced fibre. The pullback $g^*M|_S$ is also trivial. Adjunction gives $\omega_S=(\omega_Y\otimes\OO_Y(S))|_S\simeq\OO_S$. The fibre is reduced and connected, so $h^0(\OO_S)=1$; it is a Cartier divisor, hence Cohen--Macaulay and Gorenstein, so Serre duality gives $h^2(\OO_S)=1$. The divisor Riemann--Roch calculation in A.5.2 gives $\chi(\OO_S)=0$ under the stated rational cohomology hypothesis, yielding $h^1=2$.
\end{proof}


% A.4.5 from part_A.tex

\begin{statement}{A.4.5}{Proved}{The full torsion-free cotangent assertion for irreducible normal-crossings torus fibres}
Let $g:Y\to C$ have the hypotheses of A.4.4 and let $S$ be a singular irreducible reduced full fibre with normal crossings. For every nontrivial unitary-flat line bundle $A$ on $S$,
\[
H^0(S,\widetilde\Omega_S^1\otimes A)=0,
\qquad
H^0(S,\Omega_Y^1|_S\otimes A)
 \simeq H^0(\widetilde S,\eta^*A).
\tag{A.17}
\]
The latter group has dimension one when $\eta^*A$ is holomorphically trivial, and zero otherwise. The normalization is exactly one of the following: a diagonal Hopf surface with two disjoint elliptic curves identified; a ruled surface over an elliptic curve with two disjoint sections identified by translation; or a rational surface with an anticanonical hexagon as in Schr\"oer--Siebert Proposition~2.6.
\end{statement}
\begin{proof}
We check the hypotheses of the classification, including its deformation sheaf. A.4.4 gives $\omega_S\simeq\OO_S$. Dualizing the locally free conormal resolution of $\Omega_S^1$ yields
\[
 \mathcal T^1_S=\mathcal Ext^1_{\OO_S}(\Omega_S^1,\OO_S)
 \simeq\OO_S/\mathcal J=\OO_{S_{\rm sing,red}}.
\]
Here the conormal is trivial because $S$ is the full reduced fibre; the local identities $\mathcal J=(x,y)$ at a double point and $\mathcal J=(xy,xz,yz)$ at a triple point show that this is the reduced singular locus. Hence $S$ is d-semistable in the exact sense of \cite[Definition~2.1]{SS2002}. That definition imposes irreducibility, normal crossings, d-semistability and trivial dualizing sheaf; it does not require a prior smoothing to Kodaira surfaces. Propositions~2.4, 2.5 and 2.6 there give the three normalization cases. In the ruled case their Proposition~4.1 says precisely that trivial $\omega_S$ forces the section identification to be a translation. The comparison with their notation is $X_{\mathrm{SS}}=S$ and $S_{\mathrm{SS}}=\widetilde S$.

In the rational case $\eta^*A$ is trivial: its first Chern class is torsion because the line is unitary-flat, whereas the integral second cohomology of a smooth rational surface is torsion-free and its Picard group injects into that group by the exponential sequence. Such a surface has no holomorphic one-forms, by birational invariance as proved in A.4.1. The injection $\widetilde\Omega_S^1\to\eta_*\Omega^1_{\widetilde S}$ therefore proves the first vanishing.

In the elliptic ruled case the proof of A.4.3 applies to a decomposable ruled surface of any degree: the two disjoint sections express it as $\PP_E(\OO_E\oplus P)$, with $P\in\Pic(E)$. The integral second cohomology of this ruled surface is torsion-free, by the projective-bundle formula. A unitary-flat line has torsion first Chern class, so its pullback is topologically trivial and comes from $\Pic^0(E)$. No degree assumption on $P$ occurs in the cotangent calculation, and A.4.3 proves the required vanishing for every nontrivial $A$.

In the Hopf case, pull back to a finite primary diagonal Hopf cover. Its universal covering is $\C^2\setminus\{0\}$ with contraction $(z_1,z_2)\mapsto(a_1z_1,a_2z_2)$, where $|a_i|<1$. A holomorphic one-form valued in a unitary-flat line lifts to $f_1(z)dz_1+f_2(z)dz_2$, with each $f_i$ holomorphic on the punctured space and hence on $\C^2$ by Hartogs' theorem. For a contraction character $\rho$ of modulus one, equivariance says that the coefficient of $z_1^{k_1}z_2^{k_2}dz_i$ can be nonzero only when
\[
 a_1^{k_1}a_2^{k_2}a_i=\rho
\]
(or its inverse, according to the convention for the associated line). The modulus on the left is strictly less than one and the character has modulus one. Thus every coefficient is zero. Finite etale pullback is injective on sections, so the twisted one-form group on the Hopf normalization also vanishes, and again the normalization injection proves the first assertion on $S$.

Finally A.1.2 and the first vanishing identify the ambient group with $H^0(\widetilde S,\eta^*A)$. A section of a unitary-flat line on a compact connected complex manifold is either zero or parallel and nowhere zero: its squared norm is locally a sum of squares of holomorphic functions, hence plurisubharmonic, and the maximum principle makes it constant. It therefore trivializes the line, and in the trivial case the section space is $\C$. This proves the dimension statement.
\end{proof}


% A.4.6 from part_A.tex

\begin{statement}{A.4.6}{Proved}{Euler numbers of these irreducible normal-crossings fibres}
For $S$ in A.4.5,
\[
e(S)=\begin{cases}
2,&\widetilde S\text{ is rational},\\
0,&\widetilde S\text{ is elliptic ruled or Hopf}.
\end{cases}
\tag{A.18}
\]
Thus an arbitrary compact torus fibration whose fibres are irreducible reduced normal-crossings surfaces or multiples of smooth surfaces of Euler number zero has
\[
c_3(Y)=e(Y)=2r,
\tag{A.19}
\]
where $r$ is the number of rational-normalization fibres. For $b_1(Y)=b_2(Y)=0$, one has $r\in\{0,1\}$ and $b_3(Y)=2-2r$.
\end{statement}
\begin{proof}
Write $D\subset\widetilde S$ for the conductor support and $N\subset S$ for the singular curve. Normalization is a homeomorphism off these sets, so additivity of the topological Euler characteristic gives $e(S)=e(\widetilde S)-e(D)+e(N)$. In the Hopf case $e(\widetilde S)=0$, $D$ is a disjoint union of two elliptic curves, and $N$ is one elliptic curve; all three Euler numbers vanish. The ruled case has the same conductor curves and $e(\widetilde S)=e(\PP^1)e(E)=0$.

In the rational case the normalization is a Hirzebruch surface blown up twice, so its Euler number is $4+2=6$. The conductor is a cycle of six rational curves, with six nodes, so $e(D)=12-6=6$. Its image consists of three rational curves with exactly two triple points; counting the three preimages of each triple point once gives $e(N)=6-2(3-1)=2$. Hence $e(S)=2$.

Away from finitely many critical values the proper holomorphic map is a differentiable fibre bundle, and a two-torus has Euler number zero. Additivity for a stratification by this open base and the exceptional points therefore gives $e(Y)=\sum_c e((Y_c)_{\rm red})$. Nilpotents do not change the underlying topological space. The assumed fibre classes and (A.18) give (A.19). The equality with $c_3$ is the complex-oriented Gauss--Bonnet theorem. Poincare duality for the connected oriented six-manifold gives $e(Y)=2-2b_1+2b_2-b_3$. Under the two Betti vanishings this is $2-b_3$, so $b_3=2-2r\ge0$ and $r$ is zero or one.
\end{proof}


% A.4.7 from part_A.tex

\begin{statement}{A.4.7}{Proved}{Unitary-flat differentials on every compact reduced complex surface}
Let $S$ be a compact reduced complex surface, let $\mathbb A$ be a rank-one unitary complex local system on its underlying topological space, and let $A=\mathbb A\otimes_{\C}\OO_S$ be its associated holomorphic line bundle. There is a canonical complex-linear injection
\[
 H^0(S,\widetilde\Omega_S^1\otimes A)
       \longrightarrow H^1(S,\mathbb A).
 \tag{A.20}
\]
The construction uses continuous local primitives that are holomorphic on the normalization; it does not assert that those primitives are holomorphic on a non-seminormal surface.
\end{statement}
\begin{proof}
Let $\eta:\widetilde S\to S$ be normalization and choose a resolution $\pi:Z\to\widetilde S$ such that the exceptional curves have normal crossings; the components of $Z$ are compact smooth surfaces. Pullback of Kähler differentials along $h=\eta\circ\pi:Z\to S$ kills the torsion of $\Omega_S^1$, because the resulting form vanishes on the dense inverse image of $S_{\rm reg}$ and the target is locally free. Thus a section $\alpha$ on the left of (A.20) gives a holomorphic one-form $\alpha_Z$ with values in the unitary-flat line $h^*A$. This map on sections is injective for the same dense-open reason.

Write $\nabla$ for the flat unitary connection. Holomorphicity gives $\nabla\alpha_Z$ type $(2,0)$. The Hermitian pairing identifies the tensor product of the flat line with its conjugate with the trivial local system. Stokes' theorem gives, on each connected component of $Z$,
\[
 \int_Z \nabla\alpha_Z\wedge\overline{\nabla\alpha_Z}
 =\int_Z d\bigl(\alpha_Z\wedge\overline{\nabla\alpha_Z}\bigr)=0.
\]
After multiplication by the fixed dimensional sign, the integrand is the nonnegative squared norm of a $(2,0)$-form times a positive volume form. Hence $\nabla\alpha_Z=0$.

We prove the required local primitive assertion at a point $x\in S$. Trivialize $\mathbb A$ on a sufficiently small connected neighbourhood $V$ of $x$; complex analytic spaces are locally contractible, so such a flat trivialization exists. The section $\alpha$ has a local lift to $\Omega_S^1\otimes A$, since the latter surjects onto the torsion-free quotient. Its pullback to each reduced fibre of $h$ over $x$ is zero: the morphism of that fibre to $S$ is constant, so the induced morphism on Kähler differentials is zero.

For each point $\widetilde x\in\eta^{-1}(x)$, a sufficiently small regular neighbourhood of the exceptional fibre $\pi^{-1}(\widetilde x)$ in the good resolution retracts onto that fibre. This neighbourhood can be constructed by plumbing the normal disc bundles of the exceptional curves at their transverse intersection points; it is a neighbourhood of the entire fibre. A sufficiently small inverse image of a neighbourhood of $\widetilde x$ lies inside it by properness. In the smooth case one instead uses a coordinate ball. Every period of $\alpha_Z$ on the regular neighbourhood is zero, because the form is closed and its restriction to the retracting fibre is zero. Integration from a fixed point therefore produces a single-valued holomorphic primitive: a smooth primitive has differential of type $(1,0)$, so its $\bar\partial$ is zero. Normality and $\pi_*\OO_Z=\OO_{\widetilde S}$ descend that primitive to a holomorphic function $b$ on a neighbourhood of $\widetilde x$. Choose its constant so that $b(\widetilde x)=0$. Perform this independently at every point over $x$ and shrink $V$ so that all its inverse image is covered by these neighbourhoods.

These primitives have equal values on every fibre of $\eta$ over $V$. To see this, form the finite analytic fibre product
$W=\widetilde S\times_S\widetilde S$ with projections $p_1,p_2$. On the normalization of every positive-dimensional irreducible component of $W_{\rm red}$, the differential of $p_1^*b-p_2^*b$ is zero on a dense open subset: both differentials are pullbacks of the same local lift of $\alpha$ from $S$. The isolated singularities of $\widetilde S$ can be omitted for this comparison; holomorphicity then extends the equality. Thus the difference is constant on each such component. The same conclusion is automatic on a zero-dimensional component. Shrink $V$ once more so that every component of the finite space $W|_V$ meets the fibre over $x$; components disjoint from that fibre can be removed by properness. Our normalization $b(\widetilde x)=0$ makes all these constants zero. Hence $p_1^*b=p_2^*b$ as continuous functions, proving equality on the fibres of $\eta$.

A finite proper surjection is a topological quotient map. The functions $b$ therefore descend to continuous functions on $V$, holomorphic after pullback to normalization, whose differential on the regular locus is $\alpha$. On overlaps, any two such primitives differ by a locally constant function: their difference is constant on each local normalized branch, and continuity identifies the constants whenever branches meet. With the flat transition functions of $\mathbb A$, these differences define a Cech class in $H^1(S,\mathbb A)$. Changing the chosen primitives changes it by a locally constant coboundary. This defines the canonical linear map (A.20).

If the resulting class is zero, adjust the local primitives by locally constant sections so that they glue. Their pullback is then a global holomorphic section of the unitary-flat line on each compact component of $Z$. Its squared norm is plurisubharmonic, so the maximum principle makes the section parallel; its covariant differential is zero. Consequently $\alpha_Z=0$, and injectivity of pullback gives $\alpha=0$. This proves (A.20).
\end{proof}


% A.4.8 from part_A.tex

\begin{statement}{A.4.8}{Proved}{The torsion-free cotangent assertion for arbitrary singular supports in a torus fibration}
Let $g:Y\to\PP^1$ be a proper surjective holomorphic map with connected fibres from a smooth compact complex threefold with $b_1(Y)=b_2(Y)=0$, and suppose its general fibre is a complex two-torus. Let $L_r=\exp(2\pi\mathrm i r\beta)\in\Pic^0(Y)$ be the transgression family constructed in D.1.9. For every real irrational $r$ and every reduced fibre support $S=(Y_c)_{\rm red}$,
\[
 H^1(S,\mathbb A_{r,S})=0,
 \qquad
 H^0(S,\widetilde\Omega_S^1\otimes L_r|_S)=0,
 \tag{A.21}
\]
where $\mathbb A_{r,S}$ is the unitary local system whose associated holomorphic line is $L_r|_S$. In particular this proves Equation~(11) of CDP20 for all singular reduced fibre supports in the torus case, with no assumption on their singularities or their scheme-theoretic multiplicities.
\end{statement}
\begin{proof}
We specify all comparison maps used from D.1.9. On a sufficiently small disc $\Delta$ about $c$, put $U=g^{-1}(\Delta)$ and $U^*=U\setminus S$. The integral section $\gamma$ of $R^1g_*\Z$ has a lift $\gamma_U\in H^1(U,\Z)$ by Leray over the contractible disc. Its image in $H^1(U,\OO_U)$ equals $\beta|_U$, by the coherent Leray sequence. Thus the restriction of the holomorphic line $L_r$ to $U$ is the line associated to the flat character
\[
 \pi_1(U)\longrightarrow S^1,\qquad
 a\longmapsto\exp\bigl(2\pi\mathrm i r\langle\gamma_U,a\rangle\bigr).
 \tag{A.22}
\]
Denote this local system by $\mathbb A_U$. D.1.8--D.1.9 prove that $\gamma_U$ restricts nontrivially to every nearby smooth fibre. Irrationality of $r$ therefore makes (A.22) nontrivial on such a fibre $F$.

A nontrivial rank-one local system on the real torus underlying $F$ has zero cohomology in every degree. Explicitly, use the standard cellular cochain complex of the product of four circles with coefficients in the character. It is the Koszul complex for the four scalars $\chi(e_j)-1$. At least one scalar is nonzero; contraction in that factor, divided by this scalar, is a contracting homotopy. Therefore $H^q(F,\mathbb A_U|_F)=0$ for every $q$. The map $U^*\to\Delta^*$ is a differentiable fibre bundle. Its local-system Leray spectral sequence has zero fibre-cohomology sheaves, hence
$H^1(U^*,\mathbb A_U|_{U^*})=0$.

Inclusion induces a surjection $\pi_1(U^*)\to\pi_1(U)$: a loop in the smooth real six-manifold $U$ can be perturbed off the divisor, whose real codimension is at least two, relative to a base point in $U^*$. The induced map on first cohomology with these local coefficients is injective. For an explicit verification, identify first cohomology with crossed homomorphisms modulo coboundaries. If a crossed homomorphism $a:G\to M$ becomes $a(q(h))=(\chi(q(h))-1)v$ after pullback by a surjection $q:H\to G$, surjectivity gives $a(g)=(\chi(g)-1)v$ for every $g\in G$. Thus its original class was already zero. Applied to the displayed fundamental-group map, this proves $H^1(U,\mathbb A_U)=0$.

The argument applies after every sufficiently small shrinking of $\Delta$, with the restrictions of the same local system. Taking the direct limit and using proper base change identifies the resulting stalk with $H^1(S,\mathbb A_U|_S)$; scheme-theoretic nilpotents do not affect this topological local system. This proves the first vanishing in (A.21). The second is exactly A.4.7 applied to the restriction of (A.22). All singular values are treated by the same global parameter $r$, so the assertion is simultaneous.
\end{proof}


% A.4.9 from part_A.tex

\begin{statement}{A.4.9}{Proved}{CDP20 Proposition 7.3 with a complete replacement proof}
Let $Y$ be a connected smooth compact complex threefold with $H^1(Y,\Z)=H^2(Y,\Z)=0$, algebraic dimension one and holomorphic algebraic reduction $g:Y\to C$. There is a dense open subset of $\Pic(Y)$, whose complement is a proper complex analytic subset, such that
\[
 H^0\bigl((Y_c)_{\rm red},
       \widetilde\Omega^1_{(Y_c)_{\rm red}}\otimes L|_{(Y_c)_{\rm red}}\bigr)=0
 \tag{A.23}
\]
for every singular reduced fibre support and every $L$ in that subset. In fact the transgression twists $L_r$ of D.1.9 satisfy (A.23) on every reduced fibre support, including smooth and reducible supports, for every real irrational $r$.
\end{statement}
\begin{proof}
The independently rederived smooth-fibre classification C.3.1 leaves complex tori, bielliptic surfaces, primary or secondary Kodaira surfaces, and primary or secondary Hopf surfaces or Inoue surfaces. We prove that for each such smooth fibre $F$ and each nonzero $\gamma_F\in H^1(F,\Z)$ the local system
\[
 \mathbb A_{r,F}:\quad a\longmapsto
       \exp(2\pi\mathrm i r\langle\gamma_F,a\rangle)
 \quad(r\in\mathbb R\setminus\Q)
\]
satisfies $H^0(F,\mathbb A_{r,F})=H^1(F,\mathbb A_{r,F})=0$. Its character has infinite image, so the assertion in degree zero is immediate.

For a torus the Koszul argument in A.4.8 applies. A bielliptic surface has a finite étale torus cover. Pullback of the integral class is nonzero: transfer composed with pullback is multiplication by the degree. The character stays nontrivial after pullback, and transfer with local coefficients makes pullback on cohomology injective over $\C$. Thus degree-one vanishing descends from the torus cover.

A primary Kodaira surface is a principal elliptic bundle over an elliptic curve, as described in \cite[\S1 and the central-extension diagram in the proof of Lemma~1.1]{SS2002}. Its fundamental group has a central extension
\[
 1\longrightarrow\Z^2\longrightarrow\Gamma\longrightarrow\Z^2
       \longrightarrow1.
 \tag{A.24}
\]
For any nontrivial one-dimensional character $\chi$ of $\Gamma$, if its restriction to the central subgroup is nontrivial, the latter's cohomology with these coefficients is zero by the two-variable Koszul contraction. The Hochschild--Serre spectral sequence then gives $H^1(\Gamma,\C_\chi)=0$. If its central restriction is trivial, $\chi$ descends to a nontrivial character of the quotient $\Z^2$. Centrality makes its action on the cohomology of the central subgroup equal to that scalar character, and the same Koszul contraction kills every term on the second page. Again $H^1(\Gamma,\C_\chi)=0$. First cohomology of a connected space with local coefficients is first cohomology of its fundamental group, so this proves the required assertion. A secondary Kodaira surface has a finite primary cover; the same transfer argument as above applies. Hopf surfaces have finite primary covers, and a primary Hopf surface is diffeomorphic to $S^1\times S^3$; the cover statement and topology are recorded in \cite[\S2.1.1 and Theorem~2.2]{Albanese2020}. The nontrivial circle character has an invertible cellular differential, giving the required vanishing and its finite-cover descent.

For completeness we also give the Inoue calculation, rather than assuming a generic-vanishing theorem for its cotangent bundle. The explicit constructions in \cite[\S2.1.2 and Proposition~2.3]{Albanese2020} identify every Inoue surface with a mapping torus $N_\varphi$ of a compact three-manifold $N$. In the $S_M$ cases $N$ is a three-torus and the eigenvalues of $\varphi^*$ on $H^1(N,\C)$ are $\alpha,\beta,\bar\beta$, with $\alpha>1$ and $|\beta|=\alpha^{-1/2}<1$, after reversing the circle orientation when necessary. In the $S^+$ and $S^-$ cases $N$ is the nontrivial circle bundle over a two-torus; the nonzero Euler class makes $H^1(N,\C)$ the pullback of $H^1(T^2,\C)$ by its Gysin sequence. The monodromy on it has eigenvalues $\alpha,\alpha^{-1}$ or $\alpha,-\alpha^{-1}$, respectively. These assertions follow as well directly from the conjugation relations displayed in that construction: on the abelian rank-three kernel the matrix is $M$, and on the quotient of the Heisenberg kernel by its centre the matrix is $N\in\mathrm{GL}_2(\Z)$ with precisely those eigenvalues.

None of these degree-one eigenvalues has modulus one. In particular there are no invariant degree-one classes on $N$. The ordinary mapping-torus sequence therefore identifies $H^1(F,\Z)$ with the copy of $\Z$ from the circle, since the invariant subgroup of the torsion-free group $H^1(N,\Z)$ is zero. Thus $\mathbb A_{r,F}$ is pulled back from the circle, with monodromy $\lambda\in S^1$ of infinite order. The local-coefficient Leray sequence for $N\to N_\varphi\to S^1$ gives
\[
\begin{split}
0\longrightarrow
 \coker(\lambda-1:H^0(N,\C)\to H^0(N,\C))
 \longrightarrow H^1(F,\mathbb A_{r,F})\\
 \longrightarrow
 \ker(\lambda\varphi^*-1:H^1(N,\C)\to H^1(N,\C))
 \longrightarrow0.
\end{split}
\]
The first term is zero because $\lambda\ne1$, and the last term is zero by the eigenvalue moduli. Using the reciprocal convention for the circle local system replaces $\lambda$ by $\lambda^{-1}$ and gives the same conclusion. This completes the degree-one calculation for every surface type in C.3.1.

Now repeat the neighbourhood argument of A.4.8 with these two fibre-cohomology vanishings. D.1.8--D.1.9 supply the same nonzero integral class on every nearby smooth fibre and the unitary local system on every fibre neighbourhood. The only terms that can contribute to first cohomology of the punctured neighbourhood are its degree-zero and degree-one fibre cohomology, just proved zero. The fundamental-group surjection and proper base change therefore give $H^1((Y_c)_{\rm red},\mathbb A_{r,c})=0$ for every $c$. Applying A.4.7 proves (A.23) for the single global twist $L_r$, simultaneously on all reduced supports.

Finally, the exponential sequence and the integral vanishings identify $\Pic(Y)$ with the complex vector space $H^1(Y,\OO_Y)$. Choose holomorphic Cech cocycles representing a basis; exponentiating their linear combinations gives a holomorphic family of line bundles on $Y\times H^1(Y,\OO_Y)$ whose fibre is the indicated line bundle. There are finitely many singular fibre supports. For each of them, tensor its fixed coherent torsion-free cotangent sheaf by this family and apply proper semicontinuity: the locus where its space of sections is nonzero is a closed complex analytic subset of the parameter space. Each is proper because it excludes $L_r$. Their finite union is proper and has empty interior, giving the asserted dense open set. This proves the statement of Proposition~7.3. The invalid genus-two deduction in its printed proof remains refuted by A.4.2; the proof above replaces that argument through the explicit maps (A.20)--(A.24).
\end{proof}


\subsection{A.5. Normalizations and persistent ambient sections}

% A.5.1 from part_A.tex

\begin{statement}{A.5.1}{Proved}{What trivial normalization pullback forces}
Let $S$ be a connected reduced full fibre of a proper map from a complex manifold to a smooth curve. For every line bundle $A$ on $S$ with $\eta^*A\simeq\OO_{\widetilde S}$, the ambient group contains the canonical image of $H^0(\widetilde S,\OO_{\widetilde S})$. In particular it is nonzero. This applies to the subcase of Proposition~7.3 in which all $\eta^*(L_t|_S)$ are trivial. It is compatible with the vanishing of $H^0(S,Q\otimes A)$ through the exact maps of A.1.1, as A.4.1 demonstrates.
\end{statement}
\begin{proof}
A local coordinate of the base trivializes $N^*$. Tensoring the injection in (A.3) by $A$ and applying the projection formula gives
$\eta_*\eta^*A\hookrightarrow K\otimes A\hookrightarrow V\otimes A$.
Taking global sections preserves injectivity. Every connected component of the compact normalization has the nonzero constant function $1$, proving nonvanishing. No normal-crossings assumption is needed for this inclusion.
\end{proof}

The line in the printed proof asserting $H^2(\widetilde S,\omega_{\widetilde S})=0$ also cannot hold for a connected compact normal Cohen--Macaulay surface with its dualizing sheaf: Serre duality identifies its dual with $H^0(\widetilde S,\OO)=\C$. Replacing $\omega_{\widetilde S}$ there by $\OO_{\widetilde S}$ matches the vanishing stated at the beginning of that proof, but does not settle its omitted genus-one case. We have established (11) for $W$ and the exact obstruction to deducing ambient vanishing from it. The full statement of Proposition~7.3 is proved below in A.4.9 by a different argument through unitary local systems. The following argument establishes Proposition~4.11 independently of the monodromy-dependent restrictions.


% A.5.2 from part_A.tex

\begin{statement}{A.5.2}{Proved}{Normal Gorenstein surfaces with vanishing first Chern class and Euler characteristic}
A compact connected normal Gorenstein complex surface $S$ with $c_1(\omega_S)=0$ in $H^2(S,\Q)$ and $\chi(S,\OO_S)=0$ is smooth. Consequently every singular irreducible reduced divisor in a compact complex threefold $Y$ with $H^2(Y,\Q)=0$ is nonnormal. In particular the assertion of CDP20 Proposition~4.11 holds without its monodromy-dependent restrictions on the smooth fibre.
\end{statement}
\begin{proof}
Let $\pi:Z\to S$ be the minimal resolution and set $\ell=h^0(S,R^1\pi_*\OO_Z)$. Normality gives $\pi_*\OO_Z=\OO_S$, and the fibres have dimension at most one, so Leray gives $\chi(Z,\OO_Z)=-\ell$. We first prove $\ell=0$. The classification table in \cite[Theorem~2.2]{Peters2021}, together with Noether's formula and invariance of $\chi(\OO)$ under blowing up, shows that a smooth compact surface of negative holomorphic Euler characteristic is a blown-up ruled surface over a smooth curve $B$ of genus $g\ge2$. Indeed the sums $c_1^2+c_2$ in every other row of that table are nonnegative. Thus, when $\ell>0$, there is a ruling $p:Z\to B$ and $\chi(Z,\OO_Z)=1-g$, giving $\ell=g-1$.

An exceptional irreducible curve $E$ of $\pi$ dominating $B$ would have $h^1(E,\OO_E)\ge g$: normalize $E$, apply Riemann--Hurwitz to its finite map to $B$, and use the normalization exact sequence on $E$. The exact sequence $0\to\mathcal I_E\to\OO_Z\to\OO_E\to0$, pushed forward by $\pi$, gives a surjection $R^1\pi_*\OO_Z\to R^1\pi_*\OO_E$ because $R^2\pi_*\mathcal I_E=0$. It follows that $\ell\ge h^1(E,\OO_E)\ge g$, a contradiction. Hence every exceptional curve is vertical for $p$. The ruling is constant on each connected fibre of $\pi$ and descends to a holomorphic map $p_S:S\to B$: local holomorphic coordinate functions descend by $\pi_*\OO_Z=\OO_S$, and properness supplies the required neighborhoods. The map $p$ is the composite of blowups and a $\PP^1$-bundle, so $R^1p_*\OO_Z=0$. The Leray sequence for $p=p_S\circ\pi$ then gives
\[
0=R^1p_*\OO_Z\longrightarrow (p_S)_*R^1\pi_*\OO_Z
 \longrightarrow R^2(p_S)_*\OO_S=0.
\]
A nonzero sheaf supported at finitely many points cannot have zero direct image, so again $\ell=0$, a contradiction. This proves rationality of all singularities.

We give the crepancy argument to specify the comparison of canonical bundles. Proper duality applied to $R\pi_*\OO_Z=\OO_S$ identifies $\pi_*\omega_Z$ with the invertible dualizing sheaf $\omega_S$. A local generator of $\omega_S$ therefore pulls back to a holomorphic two-form on $Z$, agreeing with the usual pullback off the exceptional set. Its zero divisor is effective and exceptional. These divisors glue to $D\ge0$ with $K_Z=\pi^*K_S+D$. Every exceptional irreducible curve has arithmetic genus zero by the same surjection to its first cohomology, hence is a smooth rational curve. Minimality of the resolution gives $E^2\le-2$, and adjunction gives $D\cdot E=K_Z\cdot E=-2-E^2\ge0$. The exceptional intersection matrix is negative definite. Thus a nonzero $D$ would satisfy both $D^2<0$ and $D^2=\sum_E d_E(D\cdot E)\ge0$. Therefore $D=0$.

Consequently $c_1(K_Z)=0$ in rational cohomology. It follows by adjunction that $Z$ contains no $(-1)$-curve, so $Z$ is a minimal compact surface; Noether's formula gives $c_2(Z)=0$. We exclude rational curves on each surviving row of the classification. A ruled surface is excluded by $K_Z\cdot F=-2$ on a ruling fibre. General-type surfaces have $K_Z^2>0$ and K3 and Enriques surfaces have positive $c_2$, so none occurs. A torus has no rational curve because a map from $\PP^1$ lifts to its vector-space universal cover and is constant; the same holds for a finite etale quotient of a torus. A primary Kodaira surface is an elliptic bundle over an elliptic curve, so a rational curve would map constantly to the base and then constantly to a fibre; finite etale quotients have the same property. A class VII surface here has $b_2=c_2=0$, and the Hopf--Inoue classification gives no rational curves. Equivalently one can use their standard universal covers: $\C^2\setminus\{0\}$ for a primary Hopf cover and $\mathbb H\times\C$ for an Inoue surface. A map from $\PP^1$ lifts and has constant coordinate functions.

It remains to check the properly elliptic row. For a minimal elliptic surface, the Kodaira fibre list expresses $c_2$ as a sum of nonnegative fibre Euler numbers, with zero contributions exactly at multiples of smooth elliptic curves. Since $c_2(Z)=0$, all fibre reductions are smooth elliptic curves. The canonical bundle formula on a Kähler elliptic surface with $\chi(\OO_Z)=0$ has rational first Chern class
\[
 c_1(K_Z)=\left(2g(B)-2+\sum_i(1-m_i^{-1})\right)[F].
\]
In Kodaira dimension one the coefficient is positive, and the Kähler class has positive integral over $F$, contradicting $c_1(K_Z)=0$. In the non-Kähler case, a horizontal rational curve would be a multisection $C$. For $D=C+nF$ with $n$ sufficiently large, $D^2>0$ and $(K_Z-mD)\cdot F<0$ for $m>0$. Riemann--Roch therefore gives quadratic growth of $h^0(Z,\OO_Z(mD))$, making $Z$ Moishezon and hence projective as a smooth compact surface. This contradicts the non-Kähler case. Vertical rational curves have already been excluded by the fibre list. Thus $Z$ has no rational curves in any case. The exceptional curves of $\pi$ were rational, so there are none and $S$ is smooth.

Finally take an irreducible reduced divisor $S\subset Y$ with $H^2(Y,\Q)=0$. It is a Cartier divisor in the smooth manifold $Y$, hence Gorenstein, and adjunction gives $\omega_S=(K_Y\otimes\OO_Y(S))|_S$, with rational first Chern class zero. The divisor sequence and Hirzebruch--Riemann--Roch give
$\chi(\OO_S)=\chi(\OO_Y)-\chi(\OO_Y(-S))=0$, since $c_1(\OO_Y(S))=0$ rationally. A normal $S$ is therefore smooth by the result just proved.
\end{proof}


% A.5.3 from part_A.tex

\begin{statement}{A.5.3}{Proved}{The effective conductor divisor on a minimal resolution}
Let $S$ be a compact reduced hypersurface in a smooth threefold, let $\eta:\coprod_i T_i\to S$ be its normalization, and let $\pi_i:Z_i\to T_i$ be the minimal resolutions. Put $h_i=\eta\pi_i$. There are canonical effective integral divisors $D_i$ on $Z_i$ with
\[
 \omega_{Z_i}(D_i)\simeq h_i^*\omega_S.
 \tag{A.25}
\]
The nonexceptional part of $D_i$ is exactly the divisorial conductor, with its scheme-theoretic multiplicities. When $S\subset Y$ and $H^2(Y,\Q)=0$, every $K_{Z_i}+D_i$ has rational first Chern class zero. A ruled $Z_i\to B$ with $g(B)\ge2$ then has
\[
 D_i=2C+V,
 \tag{A.26}
\]
where $C$ is a section and $V$ is vertical. In the case where $C$ is $\pi_i$-exceptional, $\eta^*(L|_S)|_{T_i}$ is trivial for every $L\in\Pic(Y)$.
\end{statement}
\begin{proof}
Finite duality gives $\omega_{T_i}=\mathfrak c_i\otimes\eta^*\omega_S$ as rank-one reflexive sheaves, where $\mathfrak c_i$ is the conductor viewed as an ideal on $T_i$. Locally this is the equality
$\Hom_R(B,R)=\{q:qB\subset R\}=\mathfrak c$ for a reduced Gorenstein ring $R$ and its finite normalization $B$, tensored with the invertible dualizing module. The equality uses the meromorphic multiplication identification as in A.2.1. The conductor is a maximal Cohen--Macaulay module because it is this dual of the maximal Cohen--Macaulay normalization over the Gorenstein ring. On a normal surface a rank-one maximal Cohen--Macaulay module is reflexive, so the displayed equality is determined by its codimension-one orders. Those orders define an effective conductor divisor $D_i^0$ on $T_i$.

Off the exceptional set the canonical identification pulls back to $\omega_{Z_i}(D_i^{\rm str})\simeq h_i^*\omega_S$, where $D_i^{\rm str}$ denotes strict transform. The meromorphic extension of this line-bundle identification gives integral exceptional coefficients $a_j$ with
$K_{Z_i}+D_i^{\rm str}+\sum_j a_jE_j=h_i^*K_S$.
We verify their signs. For each exceptional irreducible curve $E_j$, minimality excludes a smooth rational $(-1)$-curve. Adjunction therefore gives
$K_{Z_i}\cdot E_j=2p_a(E_j)-2-E_j^2\ge0$;
also $D_i^{\rm str}\cdot E_j\ge0$. The exceptional intersection matrix $M=(E_j\cdot E_k)$ is negative definite with nonnegative off-diagonal entries, and
\[
 M a=-v,\qquad v_j=(K_{Z_i}+D_i^{\rm str})\cdot E_j\ge0.
\]
To see that $a\ge0$ without an unproved sign convention, write $a=a_+-a_-$ with disjoint nonnegative coordinate supports. If $a_-\ne0$, then
$a_-^tMa=a_-^tMa_+-a_-^tMa_->0$,
whereas $Ma=-v$ gives $a_-^tMa\le0$. This contradiction proves $a_j\ge0$. Thus $D_i=D_i^{\rm str}+\sum_j a_jE_j$ proves (A.25). The conductor duality and effective exceptional comparison also appear, for projective surfaces with trivial dualizing sheaf, in \cite[Lemma~1.1 and Corollary~1.4]{Furushima1993}; the argument here supplies the maps and does not use that paper's additional K3-type hypothesis $h^1(\OO_S)=0$.

For a divisor in $Y$, adjunction gives $\omega_S=(K_Y\otimes\OO_Y(S))|_S$, proving the numerical assertion. On a general ruling fibre $f$, $K_{Z_i}\cdot f=-2$, so $D_i\cdot f=2$. Every horizontal irreducible curve has normalization of genus at least $g(B)\ge2$. Such a curve $C_0$ cannot occur in $D_i$ with coefficient one: otherwise
\[
 0=(K_{Z_i}+D_i)\cdot C_0
   =2p_a(C_0)-2+(D_i-C_0)\cdot C_0>0.
\]
Every horizontal coefficient is consequently at least two. The degree-two intersection with $f$ forces exactly one horizontal component, coefficient two and degree one over $B$, giving (A.26). The degree-one finite morphism $C\to B$ is an isomorphism because $B$ is normal; thus $C$ is a smooth section.

Finally, integral $H^2$ of a smooth ruled surface and of its blowups is torsion-free, and its topologically trivial line bundles are precisely the pullbacks of $\Pic^0(B)$. These facts follow from the projective-bundle and blowup cohomology formulas and from $p_*\OO_{Z_i}=\OO_B$, $R^1p_*\OO_{Z_i}=0$. Restriction to the section $C$ is inverse to this Picard pullback. The line $h_i^*L$ has torsion, hence zero, integral first Chern class. When $C$ is exceptional, its restriction to $C$ is trivial because $h_i$ maps $C$ to a point. Therefore $h_i^*L$ is trivial on $Z_i$. Normality gives $(\pi_i)_*\OO_{Z_i}=\OO_{T_i}$, and the projection formula descends this triviality to $T_i$. This proves the last assertion.
\end{proof}


% A.5.4 from part_A.tex

\begin{statement}{A.5.4}{Proved}{The remaining normalization class for a persistent ambient obstruction}
Let $Y$ satisfy the CDP hypotheses, and let $S$ be a connected nonnormal reduced fibre support for which
\[
 H^0(S,\Omega_Y^1|_S\otimes L|_S)\ne0
 \quad\hbox{for every }L\in\Pic(Y).
 \tag{A.27}
\]
There is a normalization component $T_i$ whose minimal resolution $Z_i$ has Kodaira dimension $-\infty$. Every line from $\Pic(Y)$ pulls back to a unitary-flat line on that $Z_i$. A rational or class VII $Z_i$ has trivial Picard pullback from $Y$, and then
\[
 H^2\!\left(T_i,\omega_{T_i}\otimes
             \eta^*((K_Y^{-1}\otimes L)|_S)|_{T_i}\right)\simeq\C
 \quad\hbox{for every }L\in\Pic(Y).
 \tag{A.28}
\]
The same conclusion holds for the exceptional-section case of A.5.3. Thus the only normalization type left by this argument, when (A.28) has not been obtained, is an irrational ruled surface with nonconstant Picard pullback; in base genus at least two its divisor (A.26) has its horizontal section in the strict conductor.
\end{statement}
\begin{proof}
A.4.9 gives $H^0(S,\widetilde\Omega_S^1\otimes L|_S)=0$ on a dense open subset of $\Pic(Y)$. By A.1.1 and (A.27), $H^0(S,K\otimes L|_S)$ is nonzero there. Its nonvanishing locus is closed analytic in the holomorphic Picard family, hence is all of $\Pic(Y)$.

Pull $K$ to each $Z_i$, quotient the resulting sheaf by its torsion, and take its double dual. The result is a fixed holomorphic line bundle $M_i$, since its generic rank on each component is one and $Z_i$ is a smooth surface. The induced maps give an injection
\[
 H^0(S,K\otimes L|_S)\hookrightarrow
       \bigoplus_i H^0(Z_i,M_i\otimes h_i^*L).
 \tag{A.29}
\]
Indeed a section in the kernel vanishes generically on every component, and $K\subset\Omega_Y^1|_S$ is torsion-free over the reduced space. Each summand's nonvanishing locus is closed analytic in the connected complex vector space $\Pic(Y)=H^1(Y,\OO_Y)$. A finite union of proper analytic subsets cannot cover that space. Consequently one component $i$ has
$H^0(Z_i,M_i\otimes h_i^*L)\ne0$ for every $L$.

Fix a Gauduchon metric on $Z_i$, with fundamental form $\omega$ satisfying $\partial\bar\partial\omega=0$. The degree of a line bundle is the integral of its real Chern curvature against $\omega$; it is independent of the chosen Hermitian metric and additive under tensor product. A nonzero section has an effective zero divisor, so its line has nonnegative degree. For any $G\in\Pic(Y)$ the preceding nonvanishing applied to $G^{\otimes n}$ therefore gives
\[
 \deg_\omega(M_i)+n\deg_\omega(h_i^*G)\ge0
 \quad\hbox{for every }n\in\Z.
\]
It follows that $\deg_\omega(h_i^*G)=0$ for every $G$ and every Gauduchon metric.

We record the precise flatness implication. For a topologically trivial line, the real Chern curvature $F$ is exact. If $\int F\wedge\omega=0$, the Gauduchon scalar elliptic equation supplies a real $u$ with $\Lambda_\omega(F-i\partial\bar\partial u)=0$. The remaining form is closed, exact, real and primitive of type $(1,1)$, hence anti-selfdual. Its square integrates both to zero by exactness and to minus its squared norm. It is therefore zero. Changing the line metric by $e^{-u}$ gives zero Chern curvature, hence a flat unitary connection. This is exactly the curvature argument of \cite[Lemma~2.3(1)--(2)]{Teleman2006}. Since $H^2(Y,\Z)=0$, it applies to every $h_i^*G$.

The divisor $D_i$ of A.5.3 is nonzero. For an irreducible nonnormal support this follows from its nonempty conductor; for a connected reducible support every component meets another component, again giving a nonempty conductor. If $Z_i$ had nonnegative Kodaira dimension, some positive pluricanonical power would have a nonzero section, making $\deg_\omega(K_{Z_i})\ge0$. But $D_i$ is a nonzero effective divisor and therefore has strictly positive degree. Equation (A.25), with $\omega_S=(K_Y\otimes\OO_Y(S))|_S$, would give
\[
 0=\deg_\omega h_i^*(K_Y\otimes\OO_Y(S))
   =\deg_\omega(K_{Z_i})+\deg_\omega(D_i)>0,
\]
a contradiction. Thus $\kappa(Z_i)=-\infty$, whose classification consists of rational, irrational ruled, and class VII surfaces.

For a rational $Z_i$, topologically trivial lines are trivial. For a class VII $Z_i$, $h^1(\OO_{Z_i})=1$ and $b_1(Z_i)=1$. With our exponential convention $u\mapsto\exp(2\pi\mathrm i u)$, the unitary-flat subgroup of $\Pic^0(Z_i)$ has real tangent space equal to the injective image of $H^1(Z_i,\mathbb R)$ in $H^1(Z_i,\OO_{Z_i})$. This is \cite[Lemma~2.3(2)]{Teleman2006} after multiplication by the fixed scalar $\mathrm i$, or follows from the pluriharmonic maximum principle. That real tangent space has dimension one. The holomorphic Picard pullback has complex-linear differential and, by the proved flatness, takes values in this real one-dimensional subspace. Its differential is zero, so the homomorphism is constant and equals zero at the trivial line. Every $h_i^*L$ is therefore trivial. Projection formula and $(\pi_i)_*\OO_{Z_i}=\OO_{T_i}$ again descend triviality to $T_i$. The exceptional-section case has the same conclusion by A.5.3.

In each of these cases, the line multiplying $\omega_{T_i}$ in (A.28) is trivial, since it too comes from $Y$. A normal complex surface is Cohen--Macaulay, so its dualizing sheaf and Serre duality give $H^2(T_i,\omega_{T_i})=H^0(T_i,\OO_{T_i})^\vee=\C$. This proves (A.28) and the stated exact reduction of the remaining normalization problem.
\end{proof}


% A.5.5 from part_A6_critical_curves.tex

\begin{statement}{A.5.5}{Proved}{Conductor-cleared jets on a complete-intersection critical curve}
Let $Y$ be a smooth compact complex threefold with
$H^2(Y,\mathbb Q)=0$, let $S\subset Y$ be a reduced divisor,
and let $D\subset\operatorname{Sing}(S)$ be a reduced
irreducible compact curve which is a local complete intersection
in $Y$. Assume that the generic plane-curve transversal to $D$
in $S$ is an ordinary cusp, an ordinary tacnode, or an ordinary
triple point.
If the normalization $C\to D$ has positive genus, then $D$ is
a smooth elliptic curve. The transversal type is the indicated
cusp, tacnode, or triple point at every point of $D$;
in particular no degeneration of that type occurs along $D$.

In the cusp case the proof constructs, on $C$, the actual
meromorphic cubic on the radical of the Hessian and bounds its
pole divisor by twice the conductor. In the tacnode case the
corresponding fourth-order invariant has poles bounded by four
times the conductor. These intermediate jets need not extend
holomorphically before the conductor correction.
\end{statement}
\begin{proof}
Let $\nu:C\to D$ be the normalization and $j:C\to Y$
its composite with the inclusion,
$I=\mathcal I_{D/Y}$, and
\[
 N=\bigl(\nu^*(I/I^2)\bigr)^*,\qquad
 M=j^*\mathcal O_Y(S).
\]
These are a rank-two vector bundle and a line bundle on $C$.
Let $\Xi$ be the conductor divisor on $C$, so that
$\nu_*\mathcal O_C(-\Xi)$ is the conductor ideal in
$\mathcal O_D$, and put $c=\deg\Xi\geq0$.
Adjunction for the complete-intersection curve and finite
duality give exact line-bundle comparisons
\[
 \det N=\nu^*\omega_D\otimes j^*K_Y^{-1},\qquad
 \omega_C=\nu^*\omega_D\otimes\mathcal O_C(-\Xi).
                                                        \tag{A.J.1}
\]
For the second comparison one can work directly in the total
quotient ring: $\operatorname{Hom}_{\mathcal O_D}
(\nu_*\mathcal O_C,\mathcal O_D)$ consists precisely of
elements multiplying the normalization into $\mathcal O_D$,
which is its conductor. Since $D$ is Gorenstein, tensoring
this identity by its invertible dualizing sheaf is finite
duality. Both $M$ and $j^*K_Y$ have degree zero, by the
stated rational cohomology hypothesis. Consequently
\[
 \deg\det N=2g(C)-2+c.                              \tag{A.J.2}
\]
This is the normal bundle of the complete-intersection ideal,
with its specified normalization pullback. It is the bundle
needed in the following jet calculation.

First treat an ordinary cusp. Choose a local defining section
$f$ of $S$, using a local frame of $\mathcal O_Y(S)$.
Since $D\subset\operatorname{Sing}(S)$, its conormal class
maps to zero under
$I/I^2\to\Omega_Y^1|_D$. This homomorphism is injective:
its source is locally free on the reduced curve and the map
is injective at every generic point. Therefore $f\in I^2$.
The associated-graded identity for a regular sequence gives
$I^2/I^3=\operatorname{Sym}^2(I/I^2)$. The class of $f$
is thus a quadratic form
\[
 q\in H^0\bigl(C,\operatorname{Sym}^2N^*\otimes M\bigr).
\]
It has rank one generically. Let $R\subset N$ be its
saturated radical and $Q=N/R$. Both are line bundles.
The quadratic form annihilates $R$ everywhere by holomorphic
continuation, and gives a nonzero section of $M\otimes Q^{-2}$.
In particular
\[
 \deg Q\leq0.                                      \tag{A.J.3}
\]

On $D_{\rm reg}$, the third-order term on the radical is
intrinsic. Indeed, after writing the quadratic form as
$a\ell^2$, subtract a holomorphic lift of $a\ell^2$ from
$f$ and restrict its cubic class to $R^3$. Changing that lift
adds terms divisible by $\ell$, which vanish on $R^3$.
Changing normal generators has the same effect. Its nonzero
value characterizes the ordinary-cusp transversal among the
rank-one quadratic cases. Denote this cubic by $k$.

We now extend precisely this section across $C\setminus
\nu^{-1}(D_{\rm reg})$. Work at a point of $D$ and choose
two generators $X,Y$ of $I$. After a constant invertible linear
change of these generators, the normalized linear factor of
$q$ has the form $X+hY$ on every normalization branch over
that point. Such a choice exists: the saturated quotient
$N\to Q$ is surjective at each of the finitely many points,
so a constant vector can be chosen outside their finitely
many kernels. Thus $h$ is holomorphic on the normalization,
and
\[
 q=a(X+hY)^2,
\]
where $a=q_{XX}$ is the pullback of an actual coefficient
in $\mathcal O_D$. Choose a local conductor generator $d$
on the normalization. Both $d$ and $dh$ belong to
$\mathcal O_D$, since the conductor is an ideal in the
normalization contained in that ring. Choose lifts
$\widehat d,\widehat{dh},\widehat a$ to $\mathcal O_Y$,
and set
\[
 G=\widehat d^{\,2} f-
    \widehat a\bigl(\widehat d X+\widehat{dh}Y\bigr)^2.
                                                        \tag{A.J.4}
\]
The quadratic class of $G$ is zero after normalization.
It is already zero on $D$: the map from the locally free
sheaf $\operatorname{Sym}^2(I/I^2)$ to its normalization
pullback is injective because $D$ is reduced. Hence
$G\in I^3$.

Its class in $I^3/I^4=\operatorname{Sym}^3(I/I^2)$,
pulled back to $C$ and evaluated on $R^3$, is holomorphic.
Dividing that class by $d^2$ gives $k$ on the dense regular
locus. To check this equality, there $d$ is a unit and
(A.J.4) is exactly the subtraction of the square used in the
definition of the radical cubic. Different choices of lifts
add cubic terms containing its linear factor, and therefore
have zero restriction to $R^3$. Thus these local expressions
are the same meromorphic section wherever they overlap.
Since $d$ has divisor $\Xi$, the resulting bound is
\[
 0\ne k\in H^0\bigl(C,M\otimes R^{-3}
                               \otimes\mathcal O_C(2\Xi)\bigr).
                                                        \tag{A.J.5}
\]
This proves the claimed pole control without a splitting of
the formal neighbourhood of the singular curve.

Taking degrees in (A.J.5) gives $\deg R\leq2c/3$.
Together with (A.J.2)--(A.J.3), this yields
\[
 2g(C)-2+c=\deg R+\deg Q\leq\frac{2c}{3},
 \qquad 2g(C)-2+\frac c3\leq0.                     \tag{A.J.6}
\]
For $g(C)\geq1$, this forces $g(C)=1$ and $c=0$.
Conductor zero means $\mathcal O_D=\nu_*\mathcal O_C$;
a normal complex curve is smooth. Hence $D=C$ is elliptic.
Now (A.J.2)--(A.J.5) give $\deg R=\deg Q=0$.
The Hessian coefficient and radical cubic are nonzero
sections of degree-zero line bundles on this elliptic curve,
so neither has a zero. Locally the holomorphic Morse lemma
in the nondegenerate normal direction, with the curve
coordinate as parameter, writes $f=y^2+g(x,z)$.
The conditions just proved give $g=x^3u(x,z)$ with $u$ a
holomorphic unit. Replacing $x$ by a local cube root of $u$
times $x$ gives the actual local equation $y^2+x^3=0$.
This proves the asserted constant transversal type.

Next treat an ordinary tacnode. Retain $N,M,R,Q,q,a,h,d,P$
from the cusp calculation, with
$P=\widehat d X+\widehat{dh}Y$ and
$G=\widehat d^{\,2}f-\widehat aP^2\in I^3$.
The generic radical cubic is now zero, so its continuation
$[G]_3|_{R^3}$ is zero everywhere on the normalization.
Since $\ell=X+hY$ is monic in $X$, polynomial division gives
an actual holomorphic quadratic $C_2$ on the normalization
such that
\[
 [G]_3=\ell C_2.
\]
Each coefficient of $dC_2$ is in the conductor and hence in
$\mathcal O_D$. Lift those coefficients and let $H$ be the
resulting homogeneous quadratic in $X,Y$. Set
\[
 K=\widehat d^{\,2}G-PH.
\]
Its cubic class pulls back to
$d^2\ell C_2-d\ell\cdot dC_2=0$.
The same locally free graded-piece injection used above gives
$K\in I^4$. There is an exact identity in the local ring of
the ambient smooth threefold:
\[
 4\widehat a\widehat d^{\,6}f
 -(2\widehat a\widehat d^{\,2}P+H)^2
 =4\widehat a\widehat d^{\,2}K-H^2.               \tag{A.J.7}
\]
On the dense smooth locus, complete the square in the
nondegenerate normal direction. If the relevant terms there
are $a\ell^2+b\ell r^2+er^4$, the fourth-order invariant is
$\Delta=ae-b^2/4$. A substitution
$\ell\mapsto\ell+ur^2$ leaves this expression unchanged;
terms divisible by $\ell^2$ and higher changes in the radical
coordinate do not alter it. Thus it is the intrinsic section
of $M^2\otimes Q^{-2}\otimes R^{-4}$ on that dense locus.
It is nonzero precisely for the ordinary-tacnode transversal
among the rank-one Hessian cases with zero radical cubic.

Divide (A.J.7) by $4\widehat d^{\,6}$ and take the
fourth-order class on $R^4$. Since
$[H]_2=dC_2$, this expresses exactly that same invariant as
\[
 \Delta=
 \frac{a[K]_4|_{R^4}-(C_2|_{R^2})^2/4}{d^4}.
                                                        \tag{A.J.8}
\]
Its numerator is holomorphic on the normalization.
Different choices agree on the dense regular locus and hence
give one meromorphic section. Therefore
\[
 0\ne\Delta\in H^0\bigl(C,M^2\otimes Q^{-2}
                    \otimes R^{-4}\otimes\mathcal O_C(4\Xi)\bigr).
                                                        \tag{A.J.9}
\]
Write $q_0=\deg Q\leq0$. Using
$\deg R=2g(C)-2+c-q_0$, the degree of the line bundle
in (A.J.9) is
\[
 -2q_0-4\bigl(2g(C)-2+c-q_0\bigr)+4c
       =-8g(C)+8+2q_0.
\]
Its nonzero section and $g(C)\geq1$ force
$g(C)=1$ and $q_0=0$.

The section $q$ of $M\otimes Q^{-2}$ is now nowhere zero.
In the local presentation above, $a=q_{XX}$ is therefore a
unit on every normalization branch. Since it belongs to
$\mathcal O_D$, it is a unit in that local ring as well.
The coefficient $h=q_{XY}/a$ consequently belongs to
$\mathcal O_D$; here $q_{XY}$ denotes the coefficient in
the symmetric matrix, so the cross term is $2q_{XY}XY$.
We can thus repeat the entire construction with $d=1$.
Polynomial division by the monic $X+hY$ then takes place
over $\mathcal O_D$ itself; the zero radical cubic descends
by the same injectivity argument. Formula (A.J.8) now has
no poles. The resulting nonzero section of
$M^2\otimes Q^{-2}\otimes R^{-4}$ has degree $-4c$,
which forces $c=0$.

Thus again $D=C$ is smooth elliptic. Both the nonzero
quadratic coefficient and $\Delta$ are sections of
degree-zero line bundles and are nowhere zero. The parametric
Morse lemma first gives $f=y^2+g(x,z)$, and the vanished
radical cubic and nonzero fourth invariant give
$g=x^4u(x,z)$ with $u$ a unit. Taking a local fourth root
produces the actual equation $y^2+x^4=0$ at every point of
$D$. This proves the tacnode assertion.

For an ordinary triple transversal, the quadratic class
vanishes generically, hence vanishes in the locally free
sheaf $I^2/I^3$. Thus $f\in I^3$, and its initial cubic
is an honest global section
$b\in H^0(C,\operatorname{Sym}^3N^*\otimes M)$.
The discriminant of a binary cubic is homogeneous of degree
four in its coefficients and transforms by the sixth power
of the determinant of a change of variables. Its value is
therefore a section
\[
 0\ne\operatorname{disc}(b)\in
 H^0\bigl(C,M^4\otimes(\det N)^{-6}\bigr).
                                                        \tag{A.J.10}
\]
It is nonzero because the three generic tangent lines are
distinct. Equations (A.J.2) and (A.J.10) give
$2g(C)-2+c\leq0$, again forcing $g(C)=1$ and $c=0$
under the genus assumption. The discriminant has degree
zero and is nowhere zero. The three simple tangent roots
extend holomorphically on each small curve disc. The implicit
function theorem gives three smooth local branches of $S$
with those pairwise distinct normal tangents. Choose the
first two branches as the coordinate axes; the third is
$y=xu(x,z)$ with $u$ a unit, and replacing $y$ by $y/u(x,z)$
straightens it to $y=x$. This gives the local equation
$xy(x-y)=0$ up to a unit at every point of $D$.
\end{proof}


The pole allowance in (A.J.5) has a concrete local use. Put
$h=z^2-x^3$, $D=(h,w)=0$ in $\mathbb C^3$, and
\[
 f=xw^2+2zwh+x^2h^2+h^3.
\]
On the normalization $x=s^2,z=s^3$ of $D$, its quadratic
class is $s^2(w+sh)^2$. Completing the square on $D_{\rm reg}$
gives
\[
 f=x\left(w+\frac zx h\right)^2+
                  \left(1-\frac1x\right)h^3.
\]
The radical cubic has a pole of order two at $s=0$.
The conductor generator is $d=x=s^2$, and the exact
conductor-cleared identity is
\[
 x^2f-x(xw+zh)^2=x(x-1)h^3\in I_D^3.
\]
Thus the corrected jet comparison applies even in this local
model where holomorphic extension of the uncorrected cubic
would fail. The equalities are checked exactly in
\texttt{code/E\_singular\_curve.py}.

The local calculation below is related to the conductor
derivations and characteristic-zero pole calculation in
Reid~\cite[\S3.9, \S4.4--4.8]{Reid1994Nonnormal}.
We give the canonical-module comparison directly, without
using an anticanonical polarization or a classification of
the surface.


% A.5.6 from part_A.tex

\begin{statement}{A.5.6}{Proved}{Uniform negative-degree vanishing for a fixed vector bundle}
Let $Z$ be a compact smooth complex surface, $\omega$ a Gauduchon metric, and $E$ a fixed holomorphic vector bundle. There is a real constant $b(E,\omega)$ such that
\[
 \deg_\omega P<b(E,\omega)
 \quad\Longrightarrow\quad H^0(Z,E\otimes P)=0
 \tag{A.30}
\]
for every holomorphic line bundle $P$. Thus persistence of sections of $E\otimes h^*L$ for every $L\in\Pic(Y)$ forces $\deg_\omega h^*L=0$ for every $L$, for any holomorphic map $h:Z\to Y$. Persistence on one Picard component already forces degree zero on the pullback of $\Pic^0(Y)$. Under $H^2(Y,\Q)=0$, it forces degree zero and unitary flatness for every ambient Picard pullback. On a class VII surface $Z$ with $H^2(Y,\Z)=0$, the entire holomorphic Picard pullback is the zero map.
\end{statement}
\begin{proof}
Choose a Hermitian metric on $E$ and let $C_E$ bound the largest eigenvalue of the Hermitian curvature endomorphism $\mathrm i\Lambda_\omega F_E$ on the compact surface. Write $V_\omega=\int_Z\omega^2/2$ and normalize degree by
$\deg_\omega P=(2\pi)^{-1}\int_Z\mathrm iF_P\wedge\omega$.
The scalar Gauduchon elliptic equation changes a Hermitian line metric so that its contracted curvature is the constant
\[
 \mathrm i\Lambda_\omega F_P
       =c_P:=2\pi\deg_\omega(P)/V_\omega.
\]
Indeed the difference from this constant has integral zero, and the image of $u\mapsto\mathrm i\Lambda_\omega\partial\bar\partial u$ is precisely the mean-zero functions for a Gauduchon metric; this is the scalar equation used in \cite[Lemma~2.3]{Teleman2006}. The Hermitian line metric may therefore be rescaled by the corresponding $e^{-u}$.

For a holomorphic section $s$ of $E\otimes P$, the Chern-connection identity is
\[
 \mathrm i\Lambda_\omega\partial\bar\partial|s|^2
 =|\nabla^{1,0}s|^2
  -\langle\mathrm i\Lambda_\omega F_{E\otimes P}s,s\rangle.
\]
It follows directly in a holomorphic frame by differentiating the Hermitian norm; thus no Kähler identity is being assumed. At a positive maximum of $|s|^2$, the left side is nonpositive, whereas the right side is at least
$-(C_E+c_P)|s|^2$.
If $c_P<-C_E$, this is strictly positive, a contradiction. We may take $b(E,\omega)=-C_EV_\omega/(2\pi)$, proving (A.30).

Apply this bound to $P=h^*(L^{\otimes n})$ for every positive and negative integer $n$. Nonzero sections for all $n$ are incompatible with a nonzero degree, proving the persistence assertion. For persistence on a single component, absorb one representative line into the fixed bundle $E$ and apply the same argument to all powers of every line in $\Pic^0(Y)$. If $H^2(Y,\Q)=0$, a positive power of every ambient line lies in $\Pic^0(Y)$, so additivity of degree gives degree zero for every pullback. Its real first Chern class is zero, even if its integral class is torsion. The Chern curvature is therefore exact, and the exact-curvature argument in A.5.4 again produces a flat unitary connection. On a class VII surface the differential of the holomorphic Picard pullback then lies in a real one-dimensional unitary tangent space. Complex linearity makes that differential zero. The exponential sequence makes $\Pic(Y)$ a connected quotient of $H^1(Y,\OO_Y)$ when $H^2(Y,\Z)=0$, so under that integral hypothesis the entire map is identically zero.
\end{proof}


% A.5.7 from part_A6_critical_curves.tex

\begin{statement}{A.5.7}{Proved}{A pure double conductor forces a smooth cusp image}
Let $S$ be a reduced Gorenstein complex surface and
$\pi:\overline S\to S$ its finite normalization.
Suppose that, over a neighbourhood of a point $p\in S$,
$\overline S$ is smooth and its conductor ideal is exactly
$\mathcal O_{\overline S}(-2C)$ for a smooth reduced
divisor $C$; its local branches over $p$ are disjoint.
Suppose that the generic transversal of $S$ along the
reduced image $D=\pi(C)$ is an ordinary cusp.
Then $D$ is smooth at $p$, $C\to D$ is an isomorphism
near the preimage of $p$, and the germ of $S$ at $p$ is
analytically a product of an ordinary cusp and a smooth disc.
No complete-intersection hypothesis on $D$ is required.
\end{statement}
\begin{proof}
Write $R=\mathcal O_{S,p}$ and let $\overline R$ be its
finite semilocal normalization. It is a product of regular
two-dimensional local analytic rings. Write $\mathfrak c$
for the conductor, and set
\[
 A=R/\mathfrak c,\qquad
 \overline A=\overline R/\mathfrak c,\qquad
 Q=\overline R/R=\overline A/A.
\]
We first give the exact duality comparison needed below.
The $R$-module $\overline R$ is maximal Cohen--Macaulay.
Since $R$ is Gorenstein,
$\operatorname{Ext}_R^j(\overline R,R)=0$ for $j>0$.
Its dual $\operatorname{Hom}_R(\overline R,R)$ is
precisely the conductor, viewed as a fractional ideal:
a homomorphism is multiplication in the total quotient ring,
and multiplication takes $\overline R$ into $R$ exactly
for conductor elements. Applying $\operatorname{Hom}_R(-,R)$
to $0\to R\to\overline R\to Q\to0$ therefore gives
\[
 \operatorname{Ext}_R^1(Q,R)=A.                    \tag{A.J.11}
\]
The depth lemma shows that $Q$ is a Cohen--Macaulay module
of dimension one. Gorenstein duality on such modules is the
involutive functor $M\mapsto\operatorname{Ext}_R^1(M,R)$.
Consequently $A$ is Cohen--Macaulay and
\[
 Q\simeq\operatorname{Ext}_R^1(A,R)=\omega_A,
 \qquad \operatorname{End}_A(Q)=A.                \tag{A.J.12}
\]
The final equality is the canonical multiplication map:
duality identifies its endomorphism ring with
$\operatorname{End}_A(A)=A$.

The ordinary-cusp hypothesis says that $A$ is generically
reduced. Its nilradical is therefore supported at finitely
many points; a one-dimensional Cohen--Macaulay ring has
no nonzero submodule with that support. Hence $A$ is reduced.
Choose local coordinates $(s_i,z_i)$ on the normalization
branches with $C=(z_i=0)$. Put
\[
 B=\prod_i\mathbb C\{s_i\},\qquad
 K=\prod_i\operatorname{Frac}\bigl(\mathbb C\{s_i\}\bigr).
\]
The specified conductor gives the actual ring comparison
\[
 \overline A=B[\epsilon]/(\epsilon^2),             \tag{A.J.13}
\]
component by component. The projection
$\overline A\to B$ is injective on $A$, since its kernel
is nilpotent and $A$ is reduced. It identifies $B$ with
the finite normalization of the curve ring $A$.

Under (A.J.13) the inclusion of $A$ takes the form
\[
 a\longmapsto a+\delta(a)\epsilon,
\]
where $\delta:A\to B$ is a complex-linear analytic
derivation. It extends to the total quotient ring, and on
each normalization branch has the form
$\delta(a)=u_i\,da/ds_i$ for a meromorphic germ $u_i$.
Let $u=(u_i)$ and write derivatives componentwise.
The map
\[
 \Phi:B[\epsilon]/(\epsilon^2)\longrightarrow K,
 \qquad b+c\epsilon\longmapsto c-u b'
                                                        \tag{A.J.14}
\]
is $A$-linear for this specified action: multiplication by
$a+u a'\epsilon$ changes its value to
$ac+u a'b-u(ab)'=a(c-u b')$. Its kernel at every generic
point is exactly the graph of the extended derivation on
$K$. It therefore induces a generic isomorphism from $Q$
to its image in $K$. The module $Q\simeq\omega_A$ is
torsion-free on the reduced curve, so this induced map is
injective.

Differentiation of convergent power-series germs onto
convergent power-series germs is surjective, by termwise
integration. Hence the image in (A.J.14) is exactly
\[
 I=B+uB\subset K.                                \tag{A.J.15}
\]
It is an invertible fractional ideal on the semilocal
normalization $B$. Every $A$-linear endomorphism of this
rank-one module is multiplication by an element of $K$;
preserving $I$ is equivalent, branch by branch, to belonging
to $B$. Thus $\operatorname{End}_A(I)=B$ as subrings of
$K$. Combining this with the actual $A$-linear isomorphism
$Q\xrightarrow{\Phi}I$ and (A.J.12) gives $A=B$.
Because $A$ is local, there is a single branch over $p$,
and its ring is regular. This proves that $D$ is smooth
and that $C\to D$ is an isomorphism there.

Finally $A=B=\mathbb C\{s\}$ makes $u=\delta(s)$
holomorphic. The analytic coordinate change
$s'=s+u(s)z$ has invertible differential at $z=0$ and
identifies the image of $A$ modulo $z^2$ with the functions
of $s'$ alone. Since $R$ is the inverse image of $A$ under
$\overline R\to\overline A$, it follows that
\[
 R=\mathbb C\{s'\}+z^2\mathbb C\{s',z\}
   =\mathbb C\{s',z^2,z^3\}
   \simeq\mathbb C\{s',x,y\}/(y^2-x^3).
                                                        \tag{A.J.16}
\]
This is the asserted product germ.
\end{proof}


% A.5.8 from part_A6_critical_curves.tex

\begin{statement}{A.5.8}{Proved}{Rational singularities for a vertically resolved ruled normalization}
Let $\pi:Z\to T$ be a resolution of a compact connected
normal complex surface, where $Z$ is an irrational ruled
surface with ruling $p:Z\to B$. Suppose that every
exceptional curve of $\pi$ is vertical for $p$.
Then there is a unique holomorphic map
$\widetilde p:T\to B$ satisfying
$p=\widetilde p\circ\pi$, and
\[
 R\pi_*\mathcal O_Z=\mathcal O_T,
 \qquad R^1\widetilde p_*\mathcal O_T=0.           \tag{A.J.17}
\]
In particular all singularities of $T$ are rational. For
$g=g(B)$, its complete structure-sheaf cohomology dimensions
are
\[
 (h^0(T,\mathcal O_T),h^1(T,\mathcal O_T),
                         h^2(T,\mathcal O_T))=(1,g,0).
                                                        \tag{A.J.18}
\]
\end{statement}
\begin{proof}
The exceptional fibre of the proper birational map $\pi$
over a point is connected, because $T$ is normal and
$\pi_*\mathcal O_Z=\mathcal O_T$. Its curve components
are vertical by assumption, so their intersection
connectivity makes $p$ constant on that whole fibre.
It is also constant on the one-point fibres. We verify
holomorphic descent, rather than only descending the
underlying continuous map.

Fix $t\in T$ and let $U\subset B$ be a coordinate disc
containing $p(\pi^{-1}(t))$. Properness makes
$\pi(Z\setminus p^{-1}(U))$ closed and disjoint from
$t$. There is therefore an open neighbourhood $W$ of $t$
with $\pi^{-1}(W)\subset p^{-1}(U)$. The coordinate
function of $p$ on $\pi^{-1}(W)$ is a section of
$\pi_*\mathcal O_Z(W)=\mathcal O_T(W)$, so descends
to a holomorphic function on $W$. Its values lie in the
coordinate image of $U$; the inverse coordinate gives a
holomorphic map $W\to U$. These local maps agree after
pullback to the surjective map $\pi$, hence agree on
overlaps. They construct the required unique
$\widetilde p$.

Contracting vertical ruling $(-1)$-curves factors $p$ as
\[
 Z\xrightarrow{b} Z_0\xrightarrow{p_0}B,
\]
where $Z_0\to B$ is a $\mathbb P^1$-bundle and $b$
is a finite sequence of point blowups. Every such blowup
has direct image $\mathcal O$ and no higher direct image
of $\mathcal O$: the successive layers on its
infinitesimal exceptional curves are
$\mathcal O_{\mathbb P^1}(k)$ for $k\geq0$, whose
degree-one cohomology is zero, and formal functions gives
the assertion. The bundle $p_0$ has the same direct-image
property by the cohomology of $\mathbb P^1$ and local
triviality. Composition consequently gives
\[
 p_*\mathcal O_Z=\mathcal O_B,
 \qquad R^j p_*\mathcal O_Z=0\quad(j>0).          \tag{A.J.19}
\]

The sheaf $\mathcal T=R^1\pi_*\mathcal O_Z$ has
finite support at the singular points of $T$. The
low-degree Leray sequence for $p=\widetilde p\circ\pi$
contains the actual arrows
\[
 0\longrightarrow R^1\widetilde p_*\mathcal O_T
 \longrightarrow R^1p_*\mathcal O_Z
 \longrightarrow\widetilde p_*\mathcal T
 \longrightarrow R^2\widetilde p_*\mathcal O_T.    \tag{A.J.20}
\]
The last term is zero: the nonconstant proper map
$\widetilde p$ has fibres of complex dimension one.
The middle term is zero by (A.J.19). Thus
$R^1\widetilde p_*\mathcal O_T=0$ and
$\widetilde p_*\mathcal T=0$. Pushforward is faithful
on a sheaf with finite support, since its stalk sections
at finitely many points cannot disappear under a proper
map. Hence $\mathcal T=0$. Higher $R^j\pi_*$ vanish
for $j>1$ by the dimension of its fibres. This proves
(A.J.17) and the defining rational-singularity property.

Finally (A.J.17) identifies every
$H^q(T,\mathcal O_T)$ with $H^q(Z,\mathcal O_Z)$.
Equation (A.J.19) identifies the latter with
$H^q(B,\mathcal O_B)$. A compact connected smooth curve
has dimensions $(1,g)$ in degrees zero and one and no
higher cohomology, proving (A.J.18).
\end{proof}


% A.5.9 from part_A.tex

\begin{statement}{A.5.9}{Proved}{The ruled obstruction with a reduced conductor image}
Let $S\subset Y$ be an irreducible nonnormal reduced Cartier
divisor in a compact smooth complex threefold with
$H^2(Y,\Z)=0$ and $\chi(\OO_S)=0$.
Let $\eta:T\to S$ be its normalization and
$\pi:Z\to T$ its minimal resolution, and put $h=\eta\pi$.
Suppose that $Z$ is irrational ruled, with ruling $p:Z\to B$,
that every exceptional curve of $\pi$ is vertical, and that
the effective divisor of A.5.3 has the form
\[
 K_Z+D_Z=h^*\omega_S,\qquad D_Z=2C+V,                 \tag{A.RC.1}
\]
where $C$ is a section in the strict conductor and $V$ is
vertical. Suppose also that the conductor image scheme
$\Gamma=\operatorname{Spec}_S(\OO_S/\mathfrak c)$ is reduced,
and that $H^0(S,K\otimes L|_S)\ne0$ for every $L\in\Pic(Y)$,
with $K$ as in A.1.1.
Then every ambient line bundle pulls back trivially to $T$.
In particular
\[
 H^2\bigl(T,\omega_T\otimes
       \eta^*((K_Y^{-1}\otimes L)|_S)\bigr)\simeq\C
       \quad\hbox{for every }L\in\Pic(Y).            \tag{A.RC.2}
\]
The proof forces $B$ and $\Gamma$ to be smooth elliptic curves,
$Z=T$ to be minimal ruled, and $D_Z=2C$ with $C^2=0$.
Thus no local-complete-intersection assumption on the original
conductor image curve is needed in this reduced-scheme case.
\end{statement}
\begin{proof}
Write $\mathcal B=\eta_*\OO_T$ and
$\mathcal Q=\mathcal B/\OO_S$. We first specify conductor
duality, including its twists. Since $S$ is Gorenstein of
dimension two and $\mathcal B$ is maximal Cohen--Macaulay,
its dual is $\mathfrak c=\mathcal Hom_S(\mathcal B,\OO_S)$,
and maximal Cohen--Macaulay biduality gives
$\mathcal Hom_S(\mathfrak c,\OO_S)=\mathcal B$.
The quotient $\OO_\Gamma$ is Cohen--Macaulay of dimension one:
this follows from the depth lemma in
$0\to\mathfrak c\to\OO_S\to\OO_\Gamma\to0$.
Applying $\mathcal Hom_S(-,\OO_S)$ to that sequence, its
initial terms identify the actual normalization inclusion and
its cokernel as
\[
 0\longrightarrow\OO_S\longrightarrow\mathcal B
 \longrightarrow
 \mathcal Ext^1_S(\OO_\Gamma,\OO_S)\longrightarrow0,
 \qquad
 \mathcal Q\simeq\omega_\Gamma\otimes\omega_S^{-1}|_\Gamma.
                                                               \tag{A.RC.3}
\]
Adjunction makes $\omega_S$ the restriction of
$K_Y\otimes\OO_Y(S)$. Its degree is zero on every compact
curve, because $H^2(Y,\Z)=0$. Curve Riemann--Roch and Serre
duality therefore give
$\chi(\mathcal Q)=\chi(\omega_\Gamma)=-\chi(\OO_\Gamma)$.
On the other hand, A.5.8 gives rational singularities of $T$
and $\chi(\OO_T)=1-g(B)$. The normalization sequence and
$\chi(\OO_S)=0$ now give the exact numerical identity
\[
                   \chi(\OO_\Gamma)=g(B)-1.          \tag{A.RC.4}
\]

We next prove that $\Gamma$ is connected. Contract ruling
$(-1)$-curves from $Z$ to a minimal ruled bundle $Z_0\to B$.
Under a single contraction $\sigma:Z_1\to Z_2$ with exceptional
curve $E$, let $D_2=\sigma_*D_1$, let $a$ be the coefficient
of $E$ in $D_1$, and let $m$ be the multiplicity of $D_2$ at
the contraction point. The canonical divisor formula gives
\[
 K_{Z_1}+D_1=\sigma^*(K_{Z_2}+D_2)+(1+a-m)E.
                                                               \tag{A.RC.5}
\]
Numerical triviality and intersection with $E$ force $a=m-1$.
The same numerical triviality descends to $Z_2$. On $Z_0$
the pushed divisor is $2C_0+V_0$, with $C_0$ a section and
$V_0$ a nonnegative sum of fibres, so its support is connected.
Reverse the contractions. A centre outside the existing support
would produce coefficient $-1$ on its exceptional curve; that
coefficient persists on strict transforms and contradicts the
effectiveness of the final $D_Z$. A centre of multiplicity one
is a smooth point of one coefficient-one component and its
exceptional coefficient is zero, so it does not disconnect the
support. At a centre of multiplicity at least two the exceptional
curve has positive coefficient and connects all strict-transform
branches through that point. Induction proves that
$\operatorname{Supp}D_Z$ is connected.

Every positive exceptional component of $D_Z$ maps into
$\Gamma$. Indeed away from $\Gamma$ the surface $T$ equals
the Gorenstein surface $S$ and has rational singularities by
A.5.8. Such a singularity is crepant on its minimal resolution:
a generator of its dualizing sheaf pulls back holomorphically
by rational duality; its effective exceptional zero divisor
has nonnegative intersection with each exceptional curve by
adjunction and minimality, so negative definiteness makes it
zero, as in A.5.2. Thus (A.RC.1) has no positive exceptional
coefficient there. The strict conductor maps onto the whole
conductor support. It follows that
$h(\operatorname{Supp}D_Z)=\operatorname{Supp}\Gamma$,
which proves connectedness.

Let $D\subset\Gamma$ be the image of $C$. Localize (A.RC.3)
at the generic point of $D$. Because $\Gamma$ is reduced,
$\OO_\Gamma$ there is the function field $k(D)$, and its
dualizing module has dimension one over that field. The
elementary conductor quotient sequence is
\[
 0\longrightarrow\OO_\Gamma\longrightarrow
          \mathcal B/\mathfrak c
   \longrightarrow\mathcal Q\longrightarrow0.
\]
It follows that $\mathcal B/\mathfrak c$ has dimension two
over $k(D)$. The conductor coefficient two on $C$ contributes
$2\deg(C\to D)$ to this dimension. Therefore this degree is
one and there is no other normalization branch over $D$.
The generic transverse ring has conductor $t^2\C\{t\}$ and
is exactly $\C+t^2\C\{t\}=\C\{t^2,t^3\}$, the cusp ring.
Thus $D^\nu\simeq C\simeq B$. This deduction also rules out
any additional vertical branch over the same image curve.

A connected compact reduced curve has only constant global
holomorphic functions. From (A.RC.4) we therefore obtain
$h^1(\Gamma,\OO_\Gamma)=2-g(B)$.
Restriction to the reduced component $D$ gives a surjection
on $H^1$, because its kernel is a coherent sheaf on a curve
and has no $H^2$. Normalization of $D$ gives
$h^1(D,\OO_D)\ge g(B)$. Hence
\[
             2-g(B)\ge h^1(D,\OO_D)\ge g(B).         \tag{A.RC.6}
\]
As $g(B)\ge1$, all quantities equal one. The normalization
quotient of $D$ has length zero, so $D$ itself is smooth
elliptic. Statement C2d applies to this actual smooth curve
in $Y$ and makes the transversal cusp constant along it.
No other singular curve meets $D$. Connectedness of $\Gamma$
therefore gives $\Gamma=D$. Its normalization inverse image
is the single smooth curve $C$, the normalization is smooth
near it, and C2d gives $C^2=0$.

Adjunction in (A.RC.1) yields $V\cdot C=0$. A ruling
$(-1)$-curve outside $V$ cannot meet $C$, since its
intersection with $K_Z+2C+V$ would be at least $-1+2$.
A component of $V$ cannot meet $C$ either. Thus all
contractions to $Z_0$ avoid $C$ and preserve its square zero.
On the minimal bundle, $K_{Z_0}+2C_0$ is numerically trivial,
so its nonnegative vertical divisor $V_0$ is zero.
The first reverse blowup would be away from $C_0$ and have
coefficient $-1$ by (A.RC.5). There are no such blowups.
Thus $Z$ is minimal ruled and $D_Z=2C$.
An exceptional curve for $Z\to T$ would have negative square;
adjunction with $K_Z\equiv-2C$ would make it rational and
therefore vertical. A minimal ruled bundle has only fibres
of square zero, so no such exceptional curve exists and $Z=T$.

The constant-cusp computation A.2.4 now gives the injection
$K\hookrightarrow\eta_*(h^*\OO_Y(-S)\otimes\OO_Z(C))$.
Every ambient line is topologically trivial on $Z$ and is
the pullback of a degree-zero line on the elliptic curve $B$.
The section exact sequence, together with $R^1p_*\OO_Z=0$,
gives
\[
 0\longrightarrow\OO_B\longrightarrow p_*\OO_Z(C)
 \longrightarrow N_{C/Z}\longrightarrow0,
 \qquad\deg N_{C/Z}=0.                              \tag{A.RC.7}
\]
After a degree-zero twist $P$, both end terms have zero
sections unless $P=\OO_B$ or $P=N_{C/Z}^{-1}$.
The image of $\Pic(Y)$ in $\Pic^0(B)$ is connected:
$H^2(Y,\Z)=0$ makes $\Pic(Y)$ the connected image of
$H^1(Y,\OO_Y)$ under the exponential map.
All-twist persistence of $K$, and translation by the ambient
line $\OO_Y(-S)$, put this entire image in those at most two
points. It is therefore the identity. This proves triviality
of every ambient pullback. Serre duality on $T$ proves
(A.RC.2).
\end{proof}


% A.5.10 from part_A6_cusp_chern.tex

\begin{statement}{A.5.10}{Proved}{Eliminating vertical conductor on a smooth ruled cusp normalization}
Let $Y$ be a smooth compact complex threefold with
$H^2(Y,\mathbb Q)=0$, let $S\subset Y$ be a reduced
irreducible divisor, and suppose its normalization
$h:Z\to S\subset Y$ is a smooth irrational
ruled surface $p:Z\to B$.
Suppose its conductor divisor is $D_Z=2C+V$, where
$C$ is a ruling section and $V$ is vertical, and
the generic transversal along $h(C)$ is an ordinary cusp.
Then $B$ is elliptic, $V=0$, $C^2=0$, and $Z$ is
already minimal.
The image $h(C)$ is smooth elliptic and $S$ is locally
an ordinary-cusp product along it. No local
complete-intersection hypothesis on the image curve is
required. The conductor-image scheme along the vertical
components may initially be nonreduced.
\end{statement}
\begin{proof}
Let $g=g(B)\geq1$, let $F$ be the full ruling fibre
class, and write $V=\sum_i c_i\Gamma_i$ for its
irreducible vertical components. Put
$v=V\cdot C$, $n=C^2$ and $k=2g-2$.
Conductor duality and Cartier adjunction give the exact
line-bundle comparison
\[
 K_Z\otimes\mathcal O_Z(2C+V)
       =h^*(K_Y\otimes\mathcal O_Y(S)).
                                                        \tag{A.J.21}
\]
Its right side has rational first Chern class zero.
Intersecting with $C$ and using adjunction on the section
gives
\[
 n=2-2g-v=-k-v.                                  \tag{A.J.22}
\]
Let $D_p$ be the divisorial common zero of
$dp:T_Z\to p^*T_B$. Its remaining zero scheme is
zero-dimensional, say $\Xi$, so the image is
$I_\Xi\otimes p^*T_B(-D_p)$.
The kernel $T_p$ is an invertible sheaf, by the depth
lemma. Taking determinants gives the actual comparison
\[
 T_p=K_Z^{-1}\otimes p^*K_B\otimes\mathcal O_Z(D_p).
\]
The section makes $dp$ surjective at every point of $C$.
Consequently $D_p$ and $\Xi$ are disjoint from $C$,
and $T_p|_C=N_{C/Z}$.

Put $\mathcal F=h^*T_Y$. Poincare duality on the compact
oriented six-manifold $Y$ also gives
$H^4(Y,\mathbb Q)=0$. Therefore
\[
 c_1(\mathcal F)=0,\qquad c_2(\mathcal F)=0
        \quad\hbox{in rational cohomology}.       \tag{A.J.23}
\]
The differential is generically injective. Let
$L\subset\mathcal F$ be the saturation of the image
of $T_p\xrightarrow{dh}\mathcal F$.
It is invertible: a saturated rank-one subsheaf of a
locally free sheaf on a smooth surface is reflexive.
The map $T_p\to L$ has no vertical divisorial zero.
Indeed $h$ is finite, so no ruling fibre is contracted;
in characteristic zero its derivative along that fibre
is nonzero generically.

At the generic point of $C$, an ordinary cusp has
normalization differential equivalent to that of
$(u,s)\mapsto(u^2,u^3,s)$.
The ruling tangent either lies in its kernel along $C$
or does not. The vanishing order in the first case is
exactly one, and in the second case it is zero.
Write $\epsilon=1$ or $\epsilon=0$ for these respective
alternatives. There are no other horizontal conductor
components, and the normalization is an immersion away
from the conductor. Hence the actual line comparison is
\[
 L=T_p\otimes\mathcal O_Z(\epsilon C).
                                                        \tag{A.J.24}
\]
This checks both positions of the derivative kernel
relative to the ruling tangent.

The quotient $E=\mathcal F/L$ is torsion-free of rank
two. The differential induces a generically nonzero map
\[
 T_Z/T_p=I_\Xi\otimes p^*T_B(-D_p)
                    \longrightarrow E\longrightarrow E^{**}.
\]
It extends uniquely to a map from $p^*T_B(-D_p)$ into
$E^{**}$, since both are locally free and the omitted
set has codimension two; locally this is Hartogs
extension of its coefficient functions.
Let $M$ be its saturated image line in the rank-two
bundle $E^{**}$. At the generic point of $C$, the
two-by-two differential minors have common vanishing
order one. Removing the order $\epsilon$ already
assigned to $L$ leaves order $1-\epsilon$ for this map.
At the generic point of $\Gamma_i$, take a plane transversal
to its image and a normalization parameter $u$ for the
corresponding plane-curve branch. If $m_i$ is its
multiplicity, its derivative modulo the nonzero tangent
along $\Gamma_i$ vanishes to order $r_i=m_i-1$.
These are all codimension-one zeros. Put
$R=\sum_i r_i\Gamma_i$ and $r=R\cdot C$. We obtain
\[
 M=p^*T_B(-D_p)\otimes
           \mathcal O_Z\bigl((1-\epsilon)C+R\bigr).
                                                        \tag{A.J.25}
\]

The plane-curve conductor gives the inequalities
\[
 2R\leq V,\qquad 2r\leq v.                      \tag{A.J.26}
\]
To spell out their local length content, on a single
plane branch of multiplicity $m_i$ the classes
$u,u^2,\ldots,u^{m_i-1}$ are linearly independent in
the normalization quotient. Its delta invariant is thus
at least $m_i-1$. The branch is Gorenstein. Finite
duality identifies the length of its ring modulo its
conductor with the length of its normalization quotient;
the intrinsic conductor exponent is consequently twice
that delta invariant. Other branches add their
nonnegative intersection multiplicities to this exponent.
For the last comparison, writing the plane equation
as $\prod_j f_j$ multiplies the residue denominator on
branch $i$ by $\prod_{j\ne i}f_j$, whose valuation is
exactly the sum of those intersection multiplicities.
This proves (A.J.26) for every branch, including cases
with nonreduced conductor-image scheme.

Let $\ell_1$ be the length of $E^{**}/E$. Saturation
of $M$ gives a zero-dimensional subscheme $\Xi_M$ and a
line bundle $N$ with
\[
 0\longrightarrow M\longrightarrow E^{**}
       \longrightarrow I_{\Xi_M}\otimes N\longrightarrow0.
\]
Write $\ell_2=\operatorname{length}\Xi_M$.
First Chern classes give $c_1(N)=-c_1(L)-c_1(M)$.
The second Chern classes in these typed quotient
sequences and (A.J.23) give the intersection identity
\[
 L^2+L\cdot M+M^2=\ell_1+\ell_2\geq0.            \tag{A.J.27}
\]
Indeed $c_2(E)=L^2$,
$c_2(E^{**})=L^2-\ell_1$, and
$c_2(E^{**})=-L\cdot M-M^2+\ell_2$.

Set $W=V+D_p$ and $B_v=R-D_p$, both vertical
divisors, the first effective. The exact comparisons
above and (A.J.21) give the numerical classes
\[
 L\equiv(2+\epsilon)C+W+kF,\qquad
 M\equiv(1-\epsilon)C+B_v-kF.
\]
The intersection form on vertical divisors is negative
semidefinite. One can check this directly by contracting
the ruling blowups: every vertical class is a multiple
of the pullback fibre class plus exceptional classes,
on which the form is negative definite. Hence
\[
 q_v=W^2+W\cdot B_v+B_v^2
      =(W+B_v/2)^2+3B_v^2/4\leq0.
\]
Its vanishing implies $W^2=B_v^2=0$; these classes
are then numerical multiples of the fibre class.

For $\epsilon=1$, using $D_p\cdot C=0$ in the
left side of (A.J.27) gives
\[
 \ell_1+\ell_2=-6k-3v+3r+q_v
                  \leq-6k-3v/2.
\]
Nonnegativity, with $k\geq0$ and $v\geq0$,
forces $g=1$, $v=0$, $r=0$ and $q_v=0$.

For $\epsilon=0$, the same intersection calculation reads
\[
 0\leq\ell_1+\ell_2
       =-6k-2v+4r+q_v\leq-6k.                   \tag{A.J.28}
\]
It forces $g=1$, $\ell_1=\ell_2=0$, $2r=v$,
and $q_v=0$.
Thus $E=E^{**}$ is locally free, and
$L\to\mathcal F$ is a subbundle inclusion.
Along $C$, the composite
$T_C\to\mathcal F|_C\to E|_C$ vanishes generically,
because $dh$ has rank one there and its nonzero ruling
tangent image is $L|_C$. It therefore vanishes
everywhere. The exact subbundle sequence gives an
actual nonzero holomorphic line-bundle map
\[
 T_C\longrightarrow L|_C=T_p|_C=N_{C/Z}.
                                                        \tag{A.J.29}
\]
It is nonzero since the finite map $C\to h(C)$ has
nonzero derivative generically. Its Hom line has degree
$n-(2-2g)=-v$ by (A.J.22). A nonzero section on a
compact curve has nonnegative degree, so $v=0$ in this
alternative too.

In both alternatives, $v=r=0$ and $q_v=0$.
Thus $W$ is numerically a fibre multiple and
$W\cdot C=V\cdot C+D_p\cdot C=0$, so it is
numerically zero. An effective nonzero divisor on the
projective ruled surface has positive intersection
with an ample class. It follows that $W=0$ and hence
$V=D_p=0$. Now $B_v=R$ is effective, has square zero
and has zero intersection with $C$, so the same
argument gives $R=0$. Equation (A.J.22) gives $C^2=0$.
Finally $K_Z+2C$ is numerically trivial. A ruling
$(-1)$-curve $E_0$ would satisfy
$0=(K_Z+2C)\cdot E_0=-1+2C\cdot E_0$, impossible
because $C\cdot E_0$ is an integer. There are therefore
no ruling blowups and $Z$ is minimal.

Statement A.5.7 now applies everywhere along
the conductor: the normalization is smooth and its
conductor is exactly $2C$. It makes the image smooth
and the surface locally the cusp product. The image is
elliptic because $C\to h(C)$ is an isomorphism.
\end{proof}


% A.5.11 from part_A.tex

\begin{statement}{A.5.11}{Proved}{The genus of a horizontal conductor image from a local torus group}
Let $f:Y\to C_0$ be a proper holomorphic map from a smooth
compact complex threefold to a smooth curve, with connected
fibres and complex two-torus general fibre. Let $S$ be an
irreducible reduced fibre support, $\eta:T\to S$ its
normalization, and $\pi:Z\to T$ a resolution which is
irrational ruled, with ruling $p:Z\to B$.
Suppose every exceptional curve of $\pi$ is vertical, and
the strict conductor has just one horizontal irreducible
component $C$, which is a section of $p$.
Let $D=\eta\pi(C)$ and let $\Sigma=D^\nu$ be its normalization.
Then $g(\Sigma)\le1$.
In particular, when $C\to D$ has degree one, $B$ is elliptic.
This applies to the horizontal-cusp case of $D_Z=2C+V$;
it allows arbitrary multiplicity of the full fibre and
arbitrary singularities or nonreduced conductor image on
the finitely many vertical curves.
\end{statement}
\begin{proof}
There is nothing to prove when $\Sigma$ is rational, so
assume $g(\Sigma)\ge1$.
Statement A.5.8 gives a holomorphic ruling
$\widetilde p:T\to B$ with $p=\widetilde p\pi$.
Identifying $C$ with $B$ by its section map, normalization
of its nonconstant map to $D$ gives a finite surjective
holomorphic map $\phi:B\to\Sigma$.
Put $\psi=\phi\widetilde p:T\to\Sigma$.

We describe precisely the finite identifications needed to
make this map descend to $S$. Every vertical conductor
component of $T$ lies over one of finitely many points of
$B$, so $\psi$ is constant on it. Its image under $\eta$
has rational normalization: a vertical component is the
image of a rational curve on $Z$, and a nonconstant map
from that curve to its image normalization is finite.
Such an image cannot equal $D$, because $g(\Sigma)\ge1$.
Consequently these vertical images meet $D$ at finitely
many points. Away from those points and from the finitely
many singular points of $D$, every normalization preimage
of a point of $D$ lies on $C$, and its $\psi$ value is the
unique corresponding point of $\Sigma$.
At a point outside $D$ with more than one normalization
preimage, all those preimages lie on vertical conductor
components. There are no other horizontal conductor
components by hypothesis. Thus there is a finite set
$A\subset\Sigma$ such that, for every $s\in S$, either
$\psi(\eta^{-1}(s))$ is a singleton or it is contained
in $A$. One can take $A$ to contain the constant values
on all vertical conductor components and all points of
$\Sigma$ above the singularities of $D$ and its
intersections with the vertical image curves.

Choose an embedded finite tree $J$ in the underlying
oriented surface $\Sigma$ containing $A$, and let
$q:\Sigma\to\Sigma/J$ collapse that tree to one point.
Such a tree is obtained by joining the finite set by
arcs and retaining a spanning tree after subdividing
their finitely many crossings. Triangulate the surface
with $J$ a subcomplex. A contractible subcomplex is a
cofibration; contracting it gives a homotopy equivalence
of this CW pair. Concretely the homotopy extension
property extends a contraction of $J$ and supplies a
homotopy inverse to the quotient map. In particular
\[
              q_*:\pi_1(\Sigma)\xrightarrow{\sim}
                         \pi_1(\Sigma/J).          \tag{A.Pi.1}
\]
The map $q\psi$ is constant on every fibre of the finite
proper surjection $\eta$. A proper continuous surjection
between these compact Hausdorff spaces is a quotient
map, so there is a unique continuous map
\[
 a:S\longrightarrow\Sigma/J,
 \qquad a\eta=q\psi.                              \tag{A.Pi.2}
\]
Let $\delta:\Sigma\to D\subset S$ be the curve normalization
map. On its dense smooth locus, lifting through $C$ gives
$a\delta=q$. Continuity gives this equality everywhere.
Equation (A.Pi.1) therefore makes $a_*$ surjective on
fundamental groups, without any assertion that a finite
point-identification quotient has the same fundamental
group as $\Sigma$.

We next extend the actual map (A.Pi.2) to a neighbourhood
of $S$ in $Y$. The simultaneous analytic triangulation
theorem of \cite[\S3, Theorem~2, p.~464]{Lojasiewicz1964}
triangulates the pair $(Y,S)$, since a complex analytic
subset is real semianalytic. Compactness makes the
subcomplex representing $S$ finite. After barycentric
subdivision it is full: vertices lying in it span a
face in that same subcomplex. Write $\lambda_v$ for
barycentric coordinates and let $L$ be its vertex set.
On the open neighbourhood where $\sum_{v\in L}\lambda_v>0$,
the formula
\[
 r(x)=\frac{\sum_{v\in L}\lambda_v(x)v}
                  {\sum_{v\in L}\lambda_v(x)}       \tag{A.Pi.3}
\]
defines a continuous retraction to $S$ in triangulation
coordinates. Its value lies in the face spanned by the
$L$ vertices of the same simplex, so these formulas
agree across faces. This supplies an actual open
neighbourhood $N$ and a map $r:N\to S$ fixed on $S$.

Let $c=f(S)$. Properness allows a small disc $\Delta$
about $c$ with $U=f^{-1}(\Delta)\subset N$ and no other
critical value. The map $(ar)|_U:U\to\Sigma/J$ is
surjective on fundamental groups: its restriction to
$\delta(\Sigma)\subset S\subset U$ is $q$, whose
fundamental-group map is an isomorphism.
Set $U^*=U\setminus S=f^{-1}(\Delta^*)$.
The inclusion induces a surjection
$\pi_1(U^*)\twoheadrightarrow\pi_1(U)$.
Indeed a loop in the smooth six-manifold $U$ can be
perturbed relative to its base point to be transverse
to an analytic stratification of $S$; every stratum
has real codimension at least two, so the resulting
one-dimensional loop misses every stratum.

Over $\Delta^*$, the proper submersion is a smooth
torus bundle. The fundamental-group exact sequence
and $\pi_2(\Delta^*)=0$ give
\[
 1\longrightarrow\Z^4\longrightarrow\pi_1(U^*)
       \longrightarrow\Z\longrightarrow1.
\]
A lift of the generator of $\Z$ splits this group
extension, hence
\[
 \pi_1(U^*)\simeq\Z^4\rtimes_T\Z
 \twoheadrightarrow\pi_1(U)
 \twoheadrightarrow\pi_1(\Sigma/J)
 \simeq\pi_1(\Sigma).                             \tag{A.Pi.4}
\]
The first group is metabelian: its commutator subgroup
is contained in the abelian normal subgroup $\Z^4$.
Every quotient in (A.Pi.4) is therefore metabelian.
A closed orientable surface of genus at least two
has a fundamental group surjecting onto the free
group on two generators: in its usual surface
presentation, retain two $a$ generators and send all
$b$ generators and the other $a$ generators to the
identity. That free group surjects onto $S_4$ by
sending its generators to $(1234)$ and $(12)$;
their conjugates supply the adjacent transpositions.
The commutator subgroup of $S_4$ is the nonabelian
group $A_4$, so this quotient is not metabelian.
This contradicts (A.Pi.4). Thus $g(\Sigma)=1$ in the
positive-genus case. When $C\to D$ has degree one,
$\Sigma\simeq C\simeq B$, which proves the final
assertion.
\end{proof}


% A.5.12 from part_A.tex

\begin{statement}{A.5.12}{Proved}{The arithmetic genus of conductor images from coherent duality}
Let $f:Y\to C_0$ be a proper connected-fibre map from a
smooth compact complex threefold with $H^2(Y,\Z)=0$ to
a smooth curve, and let its smooth general fibre have
trivial canonical bundle. This includes complex
two-tori and primary Kodaira surfaces.
Let $mS$ be a full fibre with irreducible
reduced nonnormal support $S$. Write $\eta:T\to S$ for
normalization and $\pi:Z\to T$ for a resolution.
If the pullback homomorphism
$\Pic(Y)\to\Pic(Z)$ is nonconstant, then every reduced
irreducible component $D$ of the conductor image scheme
has arithmetic genus at most one.
In particular every such component whose normalization
has positive genus is a smooth elliptic curve.
The conductor image scheme itself is allowed to be
nonreduced, and no local-complete-intersection assumption
is imposed on its reduced components.
\end{statement}
\begin{proof}
First we determine exactly the needed canonical line on
the reduced support, without a multiplicity formula for
the direct image. The sheaf $f_*K_Y$ is torsion-free of
generic rank one on the smooth curve and hence is a
line bundle. Choose a frame near $c=f(S)$ and evaluate
it in $K_Y$. A coordinate $t$ on the base disc supplies
the conormal frame $dt$, and adjunction identifies
$K_Y|_F$ with $K_F$ after dividing by this frame.
Under that line-bundle isomorphism, the evaluated
section restricts on a general smooth fibre to a
nonzero canonical two-form, which is nowhere zero.
Its zero divisor is therefore
vertical. Shrinking the base disc to remove the other
vertical zero divisors, irreducibility of the support
gives an effective divisor $kS$, with $k\ge0$, and an
actual line-bundle comparison on $U=f^{-1}(\Delta)$:
\[
 K_Y|_U\simeq\OO_U(kS),\qquad
 \omega_S\simeq\OO_Y((k+1)S)|_S,\qquad
 \omega_S^{\otimes m}\simeq\OO_S.                 \tag{A.CD.1}
\]
The last isomorphism uses the local base equation whose
zero divisor is exactly the full fibre $mS$.
In particular no assertion $k=m-1$ is used.

Hirzebruch--Riemann--Roch in
$0\to\OO_Y(-S)\to\OO_Y\to\OO_S\to0$ gives
$\chi(\OO_S)=0$: the rational Chern character of
$\OO_Y(-S)$ equals that of $\OO_Y$, since its first
Chern class is zero. A nonzero section of the torsion
line $\omega_S$ trivializes it. Indeed its $m$th power
is a nonzero holomorphic function on the compact
connected reduced space and hence is a nonzero constant.
Thus, with $e=h^0(S,\omega_S)$, we have
\[
 e\in\{0,1\},\qquad
 e=1\Longleftrightarrow\omega_S\simeq\OO_S,
 \qquad h^1(S,\OO_S)=1+e.                         \tag{A.CD.2}
\]
Here $h^0(\OO_S)=1$ and Serre duality identifies
$h^2(\OO_S)$ with $e$.

Let $\mathcal Q=\eta_*\OO_T/\OO_S$ and let
$\Gamma\subset S$ be the conductor image scheme.
Both $S$ and $T$ are connected and compact, so their
global holomorphic functions are the same constants.
The normalization sequence consequently gives an
injection and the indicated kernel:
\[
 H^0(S,\mathcal Q)\hookrightarrow H^1(S,\OO_S),
 \qquad
 \operatorname{im}H^0(S,\mathcal Q)
 =\ker\bigl(H^1(S,\OO_S)\to H^1(T,\OO_T)\bigr).
                                                               \tag{A.CD.3}
\]
Because $H^2(Y,\Z)=0$, $\Pic(Y)$ is the connected
exponential image of $H^1(Y,\OO_Y)$. A nonconstant
holomorphic group homomorphism out of it has nonzero
differential. The differential of the stated Picard
pullback is the actual composite
\[
 H^1(Y,\OO_Y)\longrightarrow H^1(S,\OO_S)
 \longrightarrow H^1(T,\OO_T)
 \longrightarrow H^1(Z,\OO_Z).
\]
Thus its middle arrow has positive rank. Equations
(A.CD.2)--(A.CD.3) imply $h^0(\mathcal Q)\le e$.

The conductor is Cohen--Macaulay of dimension one,
and the exact duality comparison (A.RC.3) applies
without a reducedness assumption on $\Gamma$:
$\mathcal Q=\omega_\Gamma\otimes\omega_S^{-1}|_\Gamma$.
Serre duality on that curve therefore gives
\[
 h^1(\Gamma,\omega_S|_\Gamma)
       =h^0(S,\mathcal Q)\le e.                   \tag{A.CD.4}
\]
For any reduced irreducible component $D$ of $\Gamma$,
restriction of the line gives an exact quotient
$\omega_S|_\Gamma\to\omega_S|_D$. Its kernel is a
coherent sheaf on a one-dimensional compact space and
has no $H^2$, so the corresponding map on $H^1$ is
surjective. Hence
$h^1(D,\omega_S|_D)\le e$.

For $e=1$, (A.CD.2) makes this precisely
$p_a(D)=h^1(D,\OO_D)\le1$.
For $e=0$, the same inequality makes
$h^1(D,\omega_S|_D)=0$.
The restriction $\omega_S|_D$ is torsion by (A.CD.1)
and has degree zero. Curve Riemann--Roch now gives
\[
 0\le h^0(D,\omega_S|_D)
       =\chi(D,\omega_S|_D)=1-p_a(D),
\]
again proving $p_a(D)\le1$.
Finally, for a reduced irreducible curve the
normalization exact sequence gives
$p_a(D)=g(D^\nu)+\operatorname{length}
(\nu_*\OO_{D^\nu}/\OO_D)$.
Positive normalization genus therefore forces genus
one and zero normalization quotient. A normal complex
curve is smooth, which proves the final assertion.
\end{proof}


% A.5.13 from part_A.tex

\begin{statement}{A.5.13}{Proved}{The conductor bound when the general fibre has $q=1$ and $p_g=0$}
Let $f:Y\to\PP^1$ be a proper connected-fibre map from
a smooth compact complex threefold with
$H^1(Y,\Z)=H^2(Y,\Z)=0$. Suppose its smooth general
fibre has $h^1(\OO)=1$ and $h^2(\OO)=0$.
Then
\[
 R^1f_*\OO_Y=\OO_{\PP^1},\qquad R^2f_*\OO_Y=0,
 \qquad\Pic(Y)=\C\beta,                            \tag{A.Q1.1}
\]
where the last identification uses the exponential
coordinates of D.1.9 and its nonzero transgression
class $\beta$. Every full fibre has structure-sheaf
cohomology dimensions $(1,1,0)$.
For an irreducible reduced support $S$, every
nonconstant ambient Picard pullback to a normalization
resolution forces every reduced conductor component
to have arithmetic genus at most one. In particular
each positive-genus component is smooth elliptic,
without any lci or conductor-reducedness hypothesis.
\end{statement}
\begin{proof}
By E4, $E=R^1f_*\OO_Y$ is a line bundle. The section
induced by D.1.9's integral transgression class is
nowhere zero: its image under the natural base-change
map at each point is the nonzero coefficient class on
that full fibre. A zero of the section in $E$ would
have zero image under that map. Thus $E\simeq\OO$.
Coherent Leray now gives $H^1(Y,\OO_Y)=\C\beta$.
The exponential sequence and the two stated integral
cohomology vanishings give the Picard identification
in (A.Q1.1).

The coherent sheaf $G=R^2f_*\OO_Y$ has generic rank
zero, hence is a torsion sheaf. Riemann--Roch on $Y$
gives $\chi(\OO_Y)=\int_Yc_1c_2/24=0$, since
$c_1(Y)$ is zero in rational cohomology. The Euler
characteristic of coherent Leray is therefore
\[
 0=\chi(\OO_{\PP^1})-\chi(E)+\chi(G)
      =\operatorname{length}G.
\]
Thus $G=0$. Top-degree base change gives
$H^2(F,\OO_F)=0$ for every full Cartier fibre $F$.
The multiplication-by-parameter sequence then gives
$H^1(F,\OO_F)=E\otimes\C(c)$, of dimension one;
E4 gives $h^0(F,\OO_F)=1$.

For a full fibre $F=mS$ with irreducible reduced
support, the quotient $\OO_F\to\OO_S$ induces a
surjection on $H^2$, because its kernel has support
of dimension two and has no $H^3$.
Hence $h^2(S,\OO_S)=0$. Riemann--Roch in the Cartier
sequence on $Y$ gives $\chi(\OO_S)=0$ as in A.5.12,
so $h^1(S,\OO_S)=1$.

Put $\mathcal Q=\eta_*\OO_T/\OO_S$ for the
normalization $T$ and write $\Gamma$ for the conductor
image scheme. Nonconstant Picard pullback has nonzero
differential factoring through
$H^1(S,\OO_S)\to H^1(T,\OO_T)$.
The source has dimension one, so this map has zero
kernel. The normalization sequence (A.CD.3) gives
$H^0(\mathcal Q)=0$. Conductor duality and Serre
duality give
\[
 H^1(\Gamma,\omega_S|_\Gamma)=0,
 \qquad H^1(D,\omega_S|_D)=0                      \tag{A.Q1.2}
\]
for each reduced irreducible component $D$, where
the second map is the restriction surjection in the
proof of A.5.12. Adjunction makes $\omega_S$ ambient,
so $\deg(\omega_S|_D)=0$ by $H^2(Y,\Z)=0$.
Curve Riemann--Roch yields
$0\le h^0(D,\omega_S|_D)=1-p_a(D)$.
The normalization genus and its nonnegative delta
length then give the asserted smooth elliptic
conclusion whenever that genus is positive.
\end{proof}


% A.5.14 from part_A6_critical_curves.tex

\begin{statement}{A.5.14}{Proved}{Arithmetic genus two and the embedded delta-one comparison}
Let $f:Y\to C_0$ be a proper connected-fibre map from a
smooth compact complex threefold with $H^2(Y,\Q)=0$ to
a smooth curve. Suppose its smooth general fibre has
trivial canonical bundle. Let $mS$ be a full fibre with
irreducible reduced nonnormal support, and let $D$ be
a reduced irreducible component of its conductor image.
Then
\[
                         p_a(D)\le2.                \tag{A.J.30}
\]
If the normalization of $D$ has positive genus, then
$D$ is a local complete intersection in $Y$; every
singular germ of $D$ is an ordinary node or an ordinary
cusp, and the sum of its delta invariants is at most
one. In particular, when the generic transverse type
along such a positive-genus $D$ is $A_2$, $A_3$, or an
ordinary triple point, $D$ is smooth elliptic and the
transverse analytic type is constant along all of $D$.
These conclusions do not require a nonconstant ambient
Picard pullback.
\end{statement}
\begin{proof}
The evaluation, adjunction and fibre-equation comparisons
in (A.CD.1) use no Picard assumption and apply with
rational second-cohomology vanishing. They give a torsion
dualizing line $\omega_S$. The divisor Riemann--Roch
calculation gives $\chi(\OO_S)=0$. Thus, for
$e=h^0(S,\omega_S)\in\{0,1\}$, the same Serre-duality
calculation as in (A.CD.2) gives
$h^1(\OO_S)=1+e$; if $e=1$, then
$\omega_S\simeq\OO_S$.

Write $\eta:T\to S$ for normalization,
$\mathcal Q=\eta_*\OO_T/\OO_S$, and $\Gamma$ for
the full conductor image scheme. Equality of the
constant global functions on $S$ and $T$ gives the
actual injection $H^0(\mathcal Q)\hookrightarrow
H^1(\OO_S)$. Conductor duality and curve Serre duality
therefore give
\[
 h^1(\Gamma,\omega_S|_\Gamma)
       =h^0(S,\mathcal Q)\le1+e.                  \tag{A.J.31}
\]
The restriction quotient
$\omega_S|_\Gamma\twoheadrightarrow\omega_S|_D$
has a kernel supported in dimension at most one, so
its $H^1$ map is surjective. For $e=1$ it follows that
$p_a(D)=h^1(\OO_D)\le2$. For $e=0$, put
$L=\omega_S|_D$. This torsion line has degree zero,
and (A.J.31) gives $h^1(D,L)\le1$. Curve
Riemann--Roch gives
\[
 p_a(D)=1+h^1(D,L)-h^0(D,L)\le2,
\]
proving (A.J.30) in both cases.

Let $\nu:\widehat D\to D$ be normalization. The
normalization sequence gives the exact numerical
comparison
\[
 p_a(D)=g(\widehat D)+
        \sum_{x\in D}\dim_{\C}
          (\nu_*\OO_{\widehat D}/\OO_D)_x.       \tag{A.J.32}
\]
If $g(\widehat D)>0$, every local delta is consequently
at most one and their sum is at most one. We give the
analytic local rings at a point with delta one.
Let $A=\OO_{D,x}$ and let
$B=\bigoplus_{i=1}^r\C\{t_i\}$ be its finite
normalization. Here $\dim_{\C}(B/A)=1$. The residue map
of the local ring $A$ has the same value on every
normalization branch. The resulting quotient of
$\C^r$ by its diagonal has dimension $r-1$, so
$r\le2$.

For $r=2$, the subring
\[
 A_0=\{(b_1,b_2)\in\C\{t_1\}\oplus\C\{t_2\}:
                                  b_1(0)=b_2(0)\}
\]
contains $A$ and has colength one in $B$. Hence
$A=A_0$. The explicit homomorphism
$\C\{u,v\}/(uv)\to A_0$ sends
$u$ to $(t_1,0)$, $v$ to $(0,t_2)$, and a constant
to the corresponding equal pair. Expanding a
convergent series shows that it is an isomorphism.
This is the ordinary node germ.

For $r=1$, an element of $A$ with valuation one would
be a local coordinate on the normalization disc by
the holomorphic inverse function theorem. Convergent
composition in the analytic algebra $A$ would then
give $\C\{t\}\subset A$, contrary to its colength
one. Therefore every nonconstant element of $A$ has
valuation at least two. We have
$A\subset\C+t^2\C\{t\}$; both rings have
colength one in $\C\{t\}$, so equality holds.
The map
\[
 \C\{u,v\}/(v^2-u^3)\longrightarrow A,
                 \qquad (u,v)\longmapsto(t^2,t^3)
\]
is an isomorphism: every series with no linear term
is uniquely the sum of an even series in $t^2$ and
$t^3$ times a series in $t^2$. This is the ordinary
cusp germ. A delta-zero germ is its smooth
normalization disc.

We finally check the embedding into the specified
threefold. These analytic rings have embedding
dimension at most two. In local coordinates of the
smooth threefold, the tangent space of $D$ therefore
has codimension at least one. There is an element
$z$ of its ideal with $dz(x)\ne0$. The inverse
function theorem makes $z$ an ambient coordinate,
and the ideal of $D$ is the inverse image of a
height-one ideal in the regular analytic ring
$\C\{u,v\}=\OO_{Y,x}/(z)$. This ideal is principal:
a pure reduced curve in a smooth surface is a Cartier
divisor, or equivalently its reduced height-one ideal
is a product of the prime factors in this factorial
ring. Thus $I_{D,x}=(z,h)$, a regular sequence.
This proves the asserted lci embedding, including
at the node and cusp. Statement A.5.5 now applies
to the three specified generic transverse types
and proves the final smooth elliptic and constant-type
conclusion.
\end{proof}


% A.5.15 from part_A6_milnor.tex

\begin{statement}{A.5.15}{Proved}{Hirzebruch germs cannot occur in a torus degeneration}
Let $f:Y\to C_0$ be a proper holomorphic map from a
smooth complex threefold to a smooth curve, with smooth
nearby fibres complex two-tori. There is no point of
a fibre at which, in holomorphic ambient coordinates
$(x,y,z)$ and a coordinate on the base, the map has
the equation
\[
 f(x,y,z)=u(x,y,z)\bigl(x^a-y^bz^c\bigr)^m,
 \quad a\ge2,\quad b,c,m\ge1,                                    \tag{A.J.33}
\]
where $u$ is a holomorphic unit. No Betti-number
assumption on $Y$ is needed. This includes the pinch
point, every germ $x^2-y^{2k}z^{2r+1}$ with
$k\ge1$, $r\ge0$, and the reducible germ
$x^2-y^2z^2$.
\end{statement}
\begin{proof}
Put $h=x^a-y^bz^c$. We first compare the actual map
with its reduced equation. This polynomial is weighted
homogeneous of degree $a(b+c)$ for weights
$(b+c,a,a)$. On a small simply connected coordinate
neighbourhood choose a holomorphic logarithm of $u$.
Define the holomorphic coordinate change
\[
 \Phi(x,y,z)=\left(
 x u^{1/(ma)},\;
 y u^{1/(m(b+c))},\;
 z u^{1/(m(b+c))}\right).
                                                        \tag{A.J.34}
\]
Its derivative at the origin is an invertible diagonal
matrix. The holomorphic inverse function theorem and
weighted homogeneity give
$h\circ\Phi=u^{1/m}h$ and therefore
$f=(h\circ\Phi)^m$ exactly. For a sufficiently small
nonzero regular value $t$, a chosen Milnor fibre of
$h$ at one of the $m$ roots of $t$ is consequently a
compact oriented submanifold with boundary of the
smooth global fibre $Y_t$. Denote this Milnor fibre
by $M$.

We first give a direct spherical cycle, without an
irreducibility assumption. Put $d=\gcd(b,c)$,
$b'=b/d$ and $c'=c/d$, and consider the smooth affine
curve
\[
 B=\{(x,v)\in\C^2:x^a-v^d=1\}.
\]
The affine surface $H=\{h=1\}$ maps holomorphically
to $B$ by $v=y^{b'}z^{c'}$. The curve $B$ is
connected: projection to the $v$-line has $a$ sheets,
and a small loop about any root of $v^d+1$ cyclically
permutes them. It is smooth because its two partial
derivatives cannot vanish on $x^a-v^d=1$.
The function $v$ has $a$ distinct simple zeros
$P_j=(\zeta_a^j,0)$ on $B$.

Choose two of these zeros, $P$ and $Q$, and an
embedded arc $\gamma:[0,1]\to B$ joining them and
avoiding the other zeros. Choose smooth complex
functions $A(t),B_0(t)$ on its interior such that
\[
 A(t)^{b'}B_0(t)^{c'}=v(\gamma(t)),\qquad
 |A(t)|=|v(\gamma(t))|^{1/(2b')},\quad
 |B_0(t)|=|v(\gamma(t))|^{1/(2c')}.
\]
They exist by choosing a continuous argument along
the arc. The continuous map
\[
 (t,\theta)\longmapsto
 \bigl(x(\gamma(t)),A(t)e^{\mathrm i c'\theta},
                 B_0(t)e^{-\mathrm i b'\theta}\bigr)
                                                        \tag{A.J.33a}
\]
from $[0,1]\times S^1$ to $H$ collapses each boundary
circle to the corresponding point $P$ or $Q$.
It therefore defines an oriented spherical cycle
$\sigma_\gamma:S^2\to H$. Coprimality of $b'$ and
$c'$ makes each interior circle embedded. The whole
map is injective away from the collapsed ends,
because the base arc is embedded.

Let $\gamma'$ be a nearby arc with the same endpoints
and disjoint interior, obtained by an isotopy in a
thin disc neighbourhood of $\gamma$. The isotopy
keeps all other zeros of $v$ outside that disc.
Choose the arguments in (A.J.33a) continuously during
the isotopy; both coordinates tend uniformly to zero
at the ends. This gives a continuous homotopy of
maps from $S^2$, so the resulting spherical cycles
are homologous. Their supports meet only at $P$ and
$Q$. Near either endpoint, $(y,z)$ are actual
holomorphic coordinates on $H$, since $a x^{a-1}$
is nonzero there. Arrange the two arcs to be radial
with different arguments in the local $v$-coordinate.
On a sufficiently small sphere in these coordinates,
the links of the two cycles are disjoint parallel
curves of slope $(c',-b')$ on a Clifford torus.
Their oriented linking number is $-b'c'$.
Indeed the torus framing of a $(p,q)$ curve has
linking number $pq$ with its parallel copy: push the
two curves to the opposite sides of the Heegaard
splitting of $S^3$ into its two solid tori and count
meridian discs, obtaining $p q$ signed intersections.
Here $p=c'$ and $q=-b'$.

The local intersection of the two coned cycles is
that linking number. This equality follows directly
by replacing the two cones inside the small ball by
transverse relative two-chains with the same boundary
links; their intersection is the defining linking
number in its boundary three-sphere. At $Q$ both
cycle orientations reverse relative to the outgoing
radial coordinate, so the two reversals leave the
local intersection sign unchanged. Consequently
\[
 [\sigma_\gamma]^2
   =[\sigma_\gamma]\cdot[\sigma_{\gamma'}]
   =-2b'c'\ne0.                                  \tag{A.J.33b}
\]
Both cycles and their homotopy have compact support.
The weighted homothety
$(x,y,z)\mapsto(s^{b+c}x,s^a y,s^a z)$ maps
$H$ biholomorphically onto $\{h=s^{a(b+c)}\}$.
For sufficiently small positive $s$, it puts those
compact supports into any prescribed Milnor
neighbourhood of the origin. Thus (A.J.33b) is a
spherical class of nonzero square in an actual
local Milnor fibre $M$.

For an independent check when $\gcd(a,b,c)=1$,
we compute the full homotopy type and the rank of
the intersection form from two primary results.
The join theorem
\cite[Theorem L1 and Example 2.4(a)]{Nemethi1991Join}
applies to arbitrary analytic germs, including
non-isolated ones. Set $d=\gcd(b,c)$. The Milnor
fibre of $y^bz^c$ is a disjoint union of $d$ annuli:
its equations in a Milnor polydisc are
$y^{b/d}z^{c/d}=\zeta\epsilon^{1/d}$, and each
component is parametrized by
$y=A w^{c/d}$, $z=B w^{-b/d}$. Coprimality makes
this parametrization injective, and the two radius
bounds give a nonempty annulus in $w$.
The Milnor fibre of $x^a$ consists of $a$ points.
The join theorem thus identifies $M$, by a homotopy
equivalence, with the union of $a$ cones on a
disjoint union of $d$ circles, all glued along their
common base. Contracting one cone and then choosing
arcs between the $d$ base components gives
\[
 M\simeq
 \bigvee^{(a-1)d}S^2\ \vee\
 \bigvee^{(a-1)(d-1)}S^1.                          \tag{A.J.35}
\]
In particular, every class in $H_2(M,\Z)$ is a sum
of images of spherical fundamental classes.

For the same irreducible germ, the boundary calculation
of \cite[Theorem 4.1 and Lemma 4.6]{MPW2005Hirzebruch}
gives a Seifert three-manifold with base genus
\[
 g_\partial=(a-1)(d-1),\qquad
 e_0=\frac{a}{(b/d)(c/d)}>0.                     \tag{A.J.36}
\]
Here $e_0$ is its rational Seifert Euler number in
the orientation convention of that paper. Its
nonzero value gives
$b_1(\partial M)=2g_\partial$. To make this last
deduction explicit, the rational abelianized Seifert
presentation has a regular-fibre generator $v$,
exceptional-fibre generators $q_i$, the relations
$\alpha_iq_i+\beta_iv=0$ and
$\sum_iq_i+ev=0$, and $2g_\partial$ free base
generators. Eliminating the $q_i$ gives $e_0v=0$.
Thus all fibre generators vanish over $\Q$, leaving
exactly the $2g_\partial$ base generators.

Let $I_M:H_2(M,\Q)\to H_2(M,\partial M;\Q)$ be
the intersection morphism. Poincare--Lefschetz duality
and the homology sequence of the pair give
\[
 \operatorname{rank} I_M
   =b_2(M)+b_1(M)-b_1(\partial M)=a-1>0.          \tag{A.J.37}
\]
Indeed $H_3(M)=0$ by (A.J.35), so
$H_3(M,\partial M)\simeq H^1(M)$ injects into
$H_2(\partial M)$. Its cokernel is the kernel of
$I_M$, which has dimension
$b_1(\partial M)-b_1(M)$. This proves the displayed
rank formula. For $a=2$, $b=1$, $c=\ell$, the
more precise boundary is $L(2\ell,1)$ by
\cite[Theorem 4.2(2)]{MPW2005Hirzebruch}; the same
pair sequence over $\Z$ makes the rank-one
intersection form have absolute determinant
$2\ell$. The special case $\ell=1$ is the ordinary
quadratic singularity and has determinant two.

Finally let $j:M\hookrightarrow Y_t$ be the actual
embedding obtained from (A.J.34). A complex two-torus
has contractible universal cover $\C^2$, so
$\pi_2(Y_t)=0$. In particular the spherical class
constructed in (A.J.33a) maps to zero in
$H_2(Y_t,\Z)$. Intersection numbers of compact
two-cycles in the interior of $M$ are preserved by
this oriented codimension-zero embedding: the same
small transverse perturbations compute them in
both manifolds. Its square in $Y_t$ is therefore
both zero and $-2b'c'$, contradicting (A.J.33b).
In the irreducible case, (A.J.35)--(A.J.37) provide
an additional contradiction on the whole intersection
space.
This proves the exclusion for the actual map,
including every multiplicity in (A.J.33).
\end{proof}

The homotopy and boundary theorems just used are
established results; our application is the displayed
embedding and intersection comparison. The explicit
cycle also treats the reducible monomials. The exact
Seifert arithmetic, homology dimensions and ranks
are checked in \path{code/E_Hirzebruch.py}. No
reduction of an arbitrary nonordinary surface germ
to (A.J.33) is asserted.


% A.5.16 from part_A7_rational_bounds.tex

\begin{statement}{A.5.16}{Proved}{Rational normalization bounds with the exact constant-sheaf map}
Let $f:Y\to B$ be a proper holomorphic map from a smooth compact
complex threefold with $b_2(Y)=0$, with connected fibres and smooth
general fibre a two-torus. Let $mS$ be a full fibre with irreducible
nonnormal reduced support. Write $\eta:T\to S$ for normalization,
$\pi:Z\to T$ for its minimal resolution, and $\Gamma\subset S$
for the conductor-image scheme. Suppose that $Z$ is rational.
Put $Q=\eta_*\OO_T/\OO_S$ and $e=h^0(S,\omega_S)$.
Then $e\in\{0,1\}$, and
\[
\begin{gathered}
 R^1\pi_*\OO_Z=0,\qquad
 H^1(T;\mathbb C)=H^1(T,\OO_T)=H^2(T,\OO_T)=0,\\
 H^0(S,Q)\xrightarrow{\ \simeq\ }H^1(S,\OO_S),\qquad
 h^0(S,Q)=1+e,\quad h^1(S,Q)=e,\quad
 \chi(\Gamma,\OO_\Gamma)=-1.
\end{gathered}
\tag{A.RT.1}
\]
In particular $T$ has rational singularities. The coefficient map
\[
 H^1(S;\mathbb C)\longrightarrow H^1(S,\OO_S)
\tag{A.RT.2}
\]
is injective, so $b_1(S)\le1+e\le2$.

There is an effective integral divisor $D$ on $Z$ with
$\OO_Z(D)\simeq K_Z^{-1}$, whose strict-transform part is the
normalization conductor. It satisfies
\[
 h^0(D,\OO_D)=h^1(D,\OO_D)=1,
 \quad D_{\rm red}\text{ connected},
 \quad p_a(D_{\rm red})\le1.
\tag{A.RT.3}
\]
The conductor image is connected,
$p_a(\Gamma_{\rm red})\le2$, and
$e(\Gamma_{\rm red})\ge0$. Set-theoretically
$D_{\rm red}=\pi^{-1}(\eta^{-1}(\Gamma_{\rm red}))$.
Exceptional curves over points outside the conductor have
coefficient zero in $D$; they are not to be counted as conductor
branches. With $E_p=(\pi^{-1}(p))_{\rm red}$ the exact Euler
formula is
\[
 e(S)=e(Z)-e(D_{\rm red})+e(\Gamma_{\rm red})
       -\sum_{p\in\operatorname{Sing}T\setminus\eta^{-1}(\Gamma)}
                  \bigl(e(E_p)-1\bigr).
\tag{A.RT.4}
\]
In the torus smoothing of the statement, A.5.17 excludes
every point in this sum. Hence $T$ is smooth outside the
normalization conductor, and the sum is zero.
\end{statement}
\begin{proof}
The canonical evaluation argument in A.5.12 gives
$\omega_S^m\simeq\OO_S$ and $\chi(S,\OO_S)=0$.
Thus $e=0$ or $1$, with $e=1$ precisely when
$\omega_S\simeq\OO_S$, and
$h^1(S,\OO_S)=1+e$.

We first prove a useful fact without the rationality assumption.
Finite duality identifies the normal dualizing sheaf with
\[
 \omega_T=\mathfrak c_T\otimes\eta^*\omega_S
       \ \subset\ \eta^*\omega_S,
\tag{A.RT.5}
\]
where $\mathfrak c_T$ is the conductor ideal on $T$.
A nonzero global section of $\omega_T$ is therefore a nonzero
section of the torsion line $\eta^*\omega_S$ vanishing on its
nonempty conductor. Its $m$th tensor power would be a nonzero
holomorphic function on compact connected $T$ vanishing there,
which is impossible. Consequently
$H^2(T,\OO_T)=H^0(T,\omega_T)^\vee=0$.

For the proper resolution $\pi$, normality gives
$\pi_*\OO_Z=\OO_T$, and its fibres have dimension at most one.
The low-degree Leray sequence is
\[
\begin{split}
0\longrightarrow H^1(T,\OO_T)\longrightarrow H^1(Z,\OO_Z)
 \longrightarrow H^0(T,R^1\pi_*\OO_Z)\\
 \longrightarrow H^2(T,\OO_T)\longrightarrow H^2(Z,\OO_Z)
 \longrightarrow0.
\end{split}
\tag{A.RT.6}
\]
Rationality of $Z$ makes both its positive-degree groups zero.
Together with (A.RT.5), this gives $H^1(T,\OO_T)=0$ and
$R^1\pi_*\OO_Z=0$, since this last sheaf has finite support.
The fibres of $\pi$ are connected. Proper base change for the
constant sheaf gives $R^0\pi_*\mathbb C_Z=\mathbb C_T$.
The degree-one topological Leray edge map is consequently an
injection $H^1(T;\mathbb C)\hookrightarrow H^1(Z;\mathbb C)=0$.

In the normalization exact sequence
\[
 0\longrightarrow\OO_S\longrightarrow\eta_*\OO_T
       \longrightarrow Q\longrightarrow0,
\tag{A.RT.7}
\]
the map on $H^0$ is the identity on constants. The preceding
vanishings therefore make its connecting map
$H^0(Q)\to H^1(\OO_S)$ an isomorphism and its next connecting
map $H^1(Q)\to H^2(\OO_S)$ an isomorphism. These are the
maps and dimensions in (A.RT.1). In particular $\chi(Q)=1$.
The exact conductor-duality isomorphism proved in A.5.9 is
\[
 Q\simeq\omega_\Gamma\otimes\omega_S^{-1}|_\Gamma.
\tag{A.RT.8}
\]
The conductor image is a proper Cohen--Macaulay curve, including
any nilpotent structure. Serre duality on it gives
$\chi(Q)=-\chi(\omega_S|_\Gamma)$. A torsion line bundle has
degree zero on every reduced component, and therefore also has
total degree zero on this possibly nonreduced curve. Curve
Riemann--Roch gives
$\chi(\omega_S|_\Gamma)=\chi(\OO_\Gamma)$.
This proves $\chi(\OO_\Gamma)=-1$.

Here is the asserted injectivity of the complex coefficient map.
In the category of sheaves of complex vector spaces set
\[
 \mathcal V=\operatorname{coker}
     (\mathbb C_S\longrightarrow\eta_*\mathbb C_T).
\]
The diagram from constants to holomorphic functions induces the
specific sheaf map
\[
 \mathcal V\longrightarrow Q,
 \quad[(a_1,\ldots,a_r)]\longmapsto[(a_1,\ldots,a_r)] .
\tag{A.RT.9}
\]
At a point $s\in S$, the $r$ entries refer to the connected
normalization germs above $s$. A tuple of branchwise constants
that comes from an element of $\OO_{S,s}$ has equal values on
all these germs, because evaluation at their common image point
is one scalar. It is therefore a diagonal tuple. This proves
that (A.RT.9) is injective on every stalk. No $\OO_S$-module
structure on $\mathcal V$ is asserted or needed.

A finite map has zero positive-degree direct images of the
constant sheaf, by proper base change on its finite fibres.
Since $H^1(T;\mathbb C)=0$ and the map on constants is an
isomorphism, the connecting map from
$0\to\mathbb C_S\to\eta_*\mathbb C_T\to\mathcal V\to0$
identifies $H^0(S,\mathcal V)$ with $H^1(S;\mathbb C)$.
Naturality of the connecting maps gives the commutative square
\[
\begin{array}{ccc}
 H^0(S,\mathcal V)&\xrightarrow{\ \simeq\ }&H^1(S;\mathbb C)\\
 \big\downarrow&&\big\downarrow\\
 H^0(S,Q)&\xrightarrow{\ \simeq\ }&H^1(S,\OO_S).
\end{array}
\tag{A.RT.10}
\]
Its left arrow is injective by (A.RT.9), proving (A.RT.2).

A.5.3 gives the effective integral divisor $D$ and the exact
line-bundle relation
$K_Z\otimes\OO_Z(D)\simeq(\eta\pi)^*\omega_S$.
The torsion line on the right is trivial on the smooth rational
surface: its Picard group is torsion-free and has zero identity
component. Hence $\OO_Z(-D)\simeq K_Z$. The divisor sequence
$0\to K_Z\to\OO_Z\to\OO_D\to0$, rationality of $Z$ and
Serre duality prove (A.RT.3). In detail the first positive
cohomology of $K_Z$ is zero, and $H^2(K_Z)=\mathbb C$;
the positive cohomology of $\OO_Z$ is zero. Thus $D$ is
connected. Restriction $\OO_D\to\OO_{D_{\rm red}}$ is
surjective on $H^1$, because the kernel is supported in
dimension one. Consequently $p_a(D_{\rm red})\le1$.

We identify its support, rather than discarding exceptional
components. On a minimal resolution every exceptional
irreducible curve $E$ has $K_Z\cdot E\ge0$. Indeed
$E^2<0$, and adjunction gives
$K_Z\cdot E=2p_a(E)-2-E^2$; a smooth rational
$(-1)$-curve is excluded by minimality. If $E$ has coefficient
zero in $D$, then
$0=(K_Z+D)\cdot E$ forces $K_Z\cdot E=D\cdot E=0$.
Thus $E$ is disjoint from $D_{\rm red}$.
Every exceptional fibre over the normalization conductor meets
its strict transform. Connectedness of that fibre then forces
all its exceptional curves to have positive coefficient in $D$.
Conversely, a positive-coefficient exceptional curve over a point
outside the conductor would lie in a component of $D_{\rm red}$
disjoint from every strict conductor curve. The connectedness
just proved, and the nonempty strict conductor, exclude this.
This proves the stated support equality and the zero-coefficient
assertion. It also proves connectedness of $\Gamma_{\rm red}$.

For completeness the bound on its arithmetic genus follows
directly from (A.RT.8). One has
$h^1(\Gamma,\omega_S|_\Gamma)=h^0(Q)=1+e$,
and restriction to $\Gamma_{\rm red}$ is surjective on $H^1$.
For $e=1$ the line is trivial, giving
$p_a(\Gamma_{\rm red})\le2$. For $e=0$ its $H^1$ has
dimension at most one, and degree-zero Riemann--Roch gives
$p_a(\Gamma_{\rm red})=1-h^0(\omega_S|_{\Gamma_{\rm red}})
+h^1(\omega_S|_{\Gamma_{\rm red}})\le2$.
Finally normalization and resolution are isomorphisms off the
conductor and the finitely many exceptional fibres.
Stratifying by these disjoint closed sets gives (A.RT.4) by
additivity of the compactly supported Euler characteristic.
At a point outside the conductor, $T\to S$ is an
isomorphism and $S$ is normal. A normal complex surface
has isolated singularities. The reduced local equation
therefore has an isolated critical point whenever that
point is singular. Since the full fibre has equation
$u h^m$, the isolated clause of A.5.17 excludes it.
This proves the final assertion and eliminates the displayed
sum by an actual local comparison with the nearby torus.

We prove the asserted Euler bound for the conductor image,
including its possible positive-genus component. Put
$C=\Gamma_{\rm red}$. The normalization exact sequence for
the connected reduced curve $D_{\rm red}$ shows that the sum
of its normalized component genera is at most
$p_a(D_{\rm red})\le1$. Any component with genus one is
smooth, is unique, and meets the remaining rational curves
only in trees; a second attachment path joining distinct
points of that elliptic component would increase the
arithmetic genus to at least two. Each component of $C$ is
the finite image of a strict conductor component of
$D_{\rm red}$. Riemann--Hurwitz therefore shows that its
normalization is rational or elliptic, with at most one
elliptic-normalized component.

Write $n_0$ for the number of rational-normalized components
of $C$, $n_1\in\{0,1\}$ for the number of its
elliptic-normalized components, and $r_p,\delta_p$ for the
number of branches and the delta invariant at a singular
point. The exact coherent and topological normalization maps
give
\[
 \chi(\OO_C)=n_0-\sum_p\delta_p,
 \qquad e(C)=2n_0-\sum_p(r_p-1).
\]
For each reduced curve germ, independent constants on its
normalized branches give $\delta_p\ge r_p-1$.
The already proved arithmetic-genus bound and connectedness
give $\chi(\OO_C)\ge-1$, hence $e(C)\ge n_0-1$.
If $n_0\ge1$ this is nonnegative. If $n_0=0$, there is
one elliptic-normalized component and
$\sum\delta_p\le1$. A negative value could therefore only
be $-1$, from one ordinary node. Indeed its curve ring has
colength one in the direct sum of its two normalization
rings, and is contained in the subring of pairs with equal
constant value; equality of colengths makes these rings
equal. The two points over that node are distinct on the
elliptic normalization.

Here the hypersurface embedding excludes this last case.
At least one point $t\in T$ over the node receives both
local conductor branches. Otherwise partition the local
irreducible hypersurface branches according to which of
the two conductor branches they receive. Both parts are
nonempty. Two surface branches from opposite parts meet
in a curve, since their distinct equations in the regular
three-dimensional ambient local ring generate an ideal
of height two. That intersection curve belongs to the
conductor image, so is one of its two branches. Both
surface branches receive it, contradicting the partition.
For a locally irreducible hypersurface there is only one
normalization point and the assertion is immediate.

The connected resolution fibre $\pi^{-1}(t)$ lies in
$D_{\rm red}$ by the support equality proved above.
The strict conductor curves over the two branches are
elliptic, because a rational curve cannot map
nonconstantly to an elliptic curve. There is at most one
such curve on $D_{\rm red}$. If the two strict curves
were different, they would already give total genus two.
If they coincide, they meet $\pi^{-1}(t)$ at distinct
points: the map of that normal elliptic curve factors
through the normalization of its image, so one point
cannot map to both points over the node. The connected
resolution fibre then supplies an additional path joining
those points, outside the elliptic curve. Its components
are rational, since $T$ has rational singularities.
This produces a graph cycle in addition to the elliptic
genus, again forcing $p_a(D_{\rm red})\ge2$.
Both alternatives are impossible. Thus $e(C)\ge0$.
\end{proof}

We record the exact minimal-model calculation that follows from
A.5.16. This is a reduction of the outstanding Euler calculation;
the configurations with negative complement Euler number are
retained in the calculation.
Choose any sequence of contractions of smooth rational
$(-1)$-curves from $Z$ to a minimal rational surface $Z_0$,
and let $D_0$ be the pushforward of $D$. At one contraction
$\sigma:Z_i\to Z_{i-1}$, let $E$ be the exceptional curve,
$a$ its coefficient in $D_i$, and $r$ the multiplicity of
$D_{i-1}$ at the centre. The actual divisor comparison is
\[
 K_{Z_i}+D_i
   =\sigma^*(K_{Z_{i-1}}+D_{i-1})+(1+a-r)E.
 \tag{A.RT.11}
\]
Both canonical-plus-boundary line bundles are trivial, so
intersection with $E$ gives $a=r-1$. In particular the centre
lies on $D_{i-1}$, and effectivity is retained at every stage.
If $r=1$, the centre is a smooth point of a reduced boundary
component, $a=0$, and the strict boundary is isomorphic to the
old boundary. The blowup increases its complement Euler number
by one. If $r\ge2$, the new reduced boundary is exactly the
full inverse image of the old reduced boundary. Both the
surface and the boundary replace the centre by $\mathbf P^1$,
so the complement Euler number is unchanged. Thus
\[
 e(Z\setminus D_{\rm red})
 =e(Z_0\setminus (D_0)_{\rm red})+N_0,
 \qquad N_0\ge0,
 \tag{A.RT.12}
\]
where $N_0$ is the number of multiplicity-one centres in this
specified sequence of reverse blowups. Formula (A.RT.12) concerns
the resolution complement. For the actual torus smoothing,
the isolated clause of A.5.17 has already removed the
off-conductor exceptional correction in (A.RT.4). This local
argument is needed: on an abstract minimal $\mathbb F_2$,
its exceptional $(-2)$-section can be disjoint from an
anticanonical divisor without having arisen from one of the
$N_0$ blowups.

On $\mathbf P^2$, the divisor $D_0$ is an effective cubic.
Its possible reduced supports and their complement Euler
numbers are
\[
\begin{array}{c|c}
 \text{reduced support}&e(\mathbf P^2\setminus(D_0)_{\rm red})\\ \hline
 \text{smooth cubic, nodal cubic, cuspidal cubic}&3,\ 2,\ 1\\
 \text{line and conic, transverse or tangent}&1,\ 0\\
 \text{three lines forming a triangle, or concurrent}&0,\ -1\\
 \text{two distinct lines, or one line}&0,\ 1.
\end{array}
\tag{A.RT.13}
\]
For example the union of three concurrent lines has Euler
number $6-3+1=4$, whereas $e(\mathbf P^2)=3$.
The last row retains respectively a double line plus a line,
and a triple line. These are all cubic possibilities: an
irreducible singular cubic has arithmetic genus one and
therefore exactly a node or a cusp; the other degrees partition
three as $2+1$ or $1+1+1$.

On $\mathbb F_n$, with $n=0$ or $n\ge2$, write $C_0^2=-n$
for a minimal section and $F$ for a ruling fibre. Then
\[
 D_0\sim 2C_0+(n+2)F.
 \tag{A.RT.14}
\]
Its horizontal part has degree two over the ruling base.
If that part consists of two distinct sections $C_1+C_2$,
write $V$ for the vertical part, $d$ for its total degree and
$k=C_1\cdot C_2$. Adjunction on either section gives
$k+d=2$. Let $q$ be the number of distinct intersection
points of $C_1$ and $C_2$, and let $r$ count those points
whose ruling fibre occurs in $V_{\rm red}$. A vertical fibre
meets the two sections in two points unless they coincide at
that fibre. Additivity therefore gives the exact formula
\[
 e((D_0)_{\rm red})=4-q+r,
 \qquad e(\mathbb F_n\setminus(D_0)_{\rm red})=q-r\ge0.
 \tag{A.RT.15}
\]
This includes two sections and a fibre through their common
point: that union has Euler number four and complement Euler
number zero.

If the horizontal part is an irreducible bisection $C$, the
same adjunction calculation gives $p_a(C)=1-d$, so $d=0$
or $1$. For $d=0$, its Euler number is at most two, and the
complement Euler number is at least two. For $d=1$, the
bisection is a smooth rational curve, the vertical part is one
reduced fibre, and their set-theoretic intersection has one or
two points. The complement Euler number is respectively one
or two.

The remaining horizontal part is $2C$, where $C$ is a section.
Write $C\sim C_0+kF$. The vertical divisor has total degree
$d=n+2-2k$, and, with $v$ its number of distinct fibres,
\[
 e((D_0)_{\rm red})=2+v,
 \qquad e(\mathbb F_n\setminus(D_0)_{\rm red})=2-v.
 \tag{A.RT.16}
\]
If $C\ne C_0$, positivity of intersection with $C_0$ gives
$k\ge n$; hence $d\le2-n$ and $v\le2$.
For $n=0$ one also has $d\le2$. Consequently the only
negative values on a minimal ruled surface occur in the
explicit configurations
\[
 Z_0=\mathbb F_n,\quad n\ge2,\qquad
 D_0=2C_0+\sum_{j=1}^{v}a_jF_j,\qquad
 a_j\ge1,\quad \sum_j a_j=n+2,\quad v\ge3.
 \tag{A.RT.17}
\]
The fibres $F_j$ here are distinct. Equations
(A.RT.11)--(A.RT.17), together with (A.RT.4), isolate the
remaining contributions exactly; an arbitrary identification
of their conductor branches has not been substituted for an
actual hypersurface normalization map.

There is a further fundamental-group restriction on precisely
these negative minimal-model configurations. Their reduced
supports are unions of rational curves with simply connected
underlying spaces: three concurrent lines on $\mathbf P^2$,
or the section with its attached ruling fibres in (A.RT.17).
Every reverse blowup in (A.RT.11) preserves this property.
A multiplicity-one blowup leaves the support isomorphic to
the preceding one. A positive-coefficient blowup at a smooth
point adds a rational leaf, at an ordinary double point
subdivides the incidence edge, and at the initial concurrent
point replaces that vertex by an exceptional $\mathbf P^1$
with three attached branches. Subsequent supports have only
the resulting ordinary double points. None of these explicit
operations introduces a topological loop. Consequently
$D_{\rm red}$ is simply connected throughout these cases.

The actual map $h=\eta\pi$ expresses $S$ as the topological
pushout of
\[
 D_{\rm red}\hookrightarrow Z,
 \qquad D_{\rm red}\xrightarrow{\ h|_{D_{\rm red}}\ }
                           \Gamma_{\rm red}.
 \tag{A.RT.18}
\]
Indeed $h$ is an isomorphism off these closed sets and the
whole preimage of $\Gamma_{\rm red}$ is $D_{\rm red}$.
The continuous map from this pushout to $S$ is a bijection
from a compact space to a Hausdorff space, and hence is a
homeomorphism. A triangulation of the pair
$(Z,D_{\rm red})$ makes its inclusion a cofibration.
The double mapping-cylinder model for (A.RT.18) therefore
has the same homotopy type as the pushout, and van Kampen
gives
\[
 \pi_1(S)\simeq
 \pi_1(Z)*_{\pi_1(D_{\rm red})}\pi_1(\Gamma_{\rm red})
 \simeq\pi_1(\Gamma_{\rm red}).
 \tag{A.RT.19}
\]
All normalizations of the image components are $\mathbf P^1$,
by Riemann--Hurwitz. Gluing their finitely many points over
singularities identifies the topological curve with a graph
of simply connected spheres. Its fundamental group is a
free group, of rank $b_1(\Gamma_{\rm red})$.

That free group is a quotient of a metabelian group. To see
the exact comparison, triangulate $(Y,S)$ and retract a small
regular neighbourhood $N$ onto $S$, as in A.5.11. Properness
puts a sufficiently small full fibre tube $U=f^{-1}(\Delta)$
inside $N$. Its retraction is surjective on fundamental
groups because it restricts to the identity on $S$.
General position makes
$\pi_1(U\setminus S)\to\pi_1(U)$ surjective, and the
torus bundle over $\Delta^*$ has
$\pi_1(U\setminus S)=\mathbb Z^4\rtimes\mathbb Z$.
This last group has abelian commutator subgroup, so its
quotient (A.RT.19) is metabelian. A free group of rank at
least two is not metabelian: it maps onto $S_4$, whose
commutator group $A_4$ is not abelian. We obtain
\[
 \pi_1(S)\simeq\pi_1(\Gamma_{\rm red})\in\{1,\mathbb Z\},
 \qquad b_1(\Gamma_{\rm red})\le1.
 \tag{A.RT.20}
\]
Under the full CDP Betti hypotheses, the nonzero class on
every support from D.1.9 excludes the trivial group in
(A.RT.20). This last use of the global transgression is
separate from the preceding $b_2(Y)=0$ calculations.

For these two boundary classes one can further exclude a
one-component image with first Betti number one. Such an image
has one two-branch singular point $q$; write $\alpha,\beta$
for its two normalization points there. Any two normalization
points of $S$ above $q$ receive a common local conductor branch,
because the corresponding two hypersurface branches meet in
a curve. Thus there cannot be both a point receiving only
$\alpha$ and a point receiving only $\beta$. After interchanging
the labels, every normalization point above $q$ receives
$\alpha$.

Form the incidence graph with one vertex for each strict
conductor component and one vertex for each connected
resolution fibre over a normalization point above $q$.
Edges record their intersection points, labelled by the
branch $\alpha$ or $\beta$ received there. Smooth resolution
fibres, consisting of one point, are included. The graph is
a forest: it is obtained from the tree of components and
intersection points of $D_{\rm red}$ by contracting disjoint
exceptional subtrees and deleting the other incidence edges.
Every strict conductor component maps finitely onto the sole
image component, so it has a preimage of $\beta$ and therefore
a $\beta$-edge. Every resolution-fibre vertex has an
$\alpha$-edge by the preceding paragraph. Starting at any
strict-component vertex, follow a $\beta$-edge and then an
$\alpha$-edge, and repeat. This walk never immediately reverses
an edge, because its consecutive labels differ. It is an
infinite nonbacktracking walk in a finite forest, which is
impossible. The contradiction proves the exclusion.

Since every image component is rational-normalized, write
$n_\Gamma$ for their number and $\beta_\Gamma\le1$ for the
first Betti number. Their exact topological normalization
formula now yields, on these negative boundary classes,
\[
 e(\Gamma_{\rm red})=n_\Gamma+1-\beta_\Gamma\ge2.
 \tag{A.RT.21}
\]
For one component the excluded case is the only value below
two; for at least two components the inequality follows
directly from $\beta_\Gamma\le1$.


% A.5.17 from part_A6_milnor.tex

\begin{statement}{A.5.17}{Proved}{Isolated singularities are excluded and suspension germs are transverse products}
Let $f:Y\to C_0$ have the hypotheses of A.5.15.
No local equation $f=u h^m$, with $u$ a holomorphic
unit and $m\ge1$, can have a reduced equation $h$
with an isolated critical point at the origin.
Suppose its local equation at a point of a fibre is
\[
             f=(x^a-g(y,z))^m,
             \qquad a\ge2,\quad m\ge1,             \tag{A.J.38}
\]
where $g\ne0$ and $g(0,0)=0$. Then the reduced
plane germ $\{g=0\}_{\rm red}$ is smooth. If
$b$ is its multiplicity in $\operatorname{div}(g)$,
there are local holomorphic coordinates in which
$f=(x^a-y^b)^m$.
Consequently, at every point where the full fibre
is a uniform multiple of its reduced support and
the reduced equation has a nonzero Hessian, its
critical surface germ is a constant
transverse $A_{b-1}$ product. A nonordinary change
of transverse type at a point with nonzero Hessian
is excluded.
\end{statement}
\begin{proof}
We first prove the isolated assertion, including its
comparison with the actual global fibre. Replace $h$
by $H=u^{1/m}h$, taking a holomorphic root on a small
coordinate neighbourhood. Then $f=H^m$ exactly.
This replacement preserves an isolated critical point
as a germ. Indeed every local irreducible component
of the critical locus of $H$ has constant $H$-value,
necessarily zero when it contains the origin. On
$H=0$ one has $dH=u^{1/m}dh$. Thus its critical
locus has the same zero-dimensional germ as that
of $h$.

Choose a small closed ball $K\subset\C^3$ such that
$H^{-1}(0)$ is smooth along its boundary intersection
and transverse to $\partial K$.
The standard isolated-singularity Milnor-ball
construction supplies this choice; see
\cite[Lemma 6.1 and Appendix B]{Milnor1968}.
The restriction $H|_{\partial K}$ is a submersion
at its zero set. Compactness therefore supplies
$\eta>0$ for which it is a submersion wherever
$|H|<2\eta$ on the boundary. For the linear
deformation
\[
 H_\lambda(w)=H(w)-\lambda_1w_1
                         -\lambda_2w_2-\lambda_3w_3,
                                                        \tag{A.J.38a}
\]
this boundary submersion persists for all parameters
in a sufficiently small polydisc $P$ and all values
of modulus less than $\eta$. The gradient map of
$H$ has an isolated zero and consequently is a finite
holomorphic map germ of positive local degree. Choose
a small regular value $\lambda$ of that map. Its
nonempty finite inverse image consists of critical
points of $H_\lambda$ with invertible Hessian. All
these points and their critical values tend to the
origin as $\lambda\to0$.

Remove the critical-value discriminant from
$P\times\Delta_\eta$. The relative critical locus
over this domain misses the boundary, so its
projection is proper and its image is a proper
complex analytic subset. Its complement is path
connected. The projection of the corresponding
smooth level sets in $K$ is a proper submersion
with boundary, and hence an Ehresmann bundle.
At any one of the nondegenerate critical points,
the holomorphic Morse lemma gives
$H_\lambda-H_\lambda(p)=X^2+Y^2+Z^2$.
A local Milnor fibre at a sufficiently small regular
value embeds in the corresponding level set in $K$.
The bundle transports it, by an oriented embedding,
into $K\cap H^{-1}(\epsilon_0)$ for a nonzero small
$\epsilon_0$, which is a submanifold of the actual
global torus $Y_{\epsilon_0^m}$.

For a positive real local Milnor value, the real
locus $X,Y,Z\in\mathbb R$ is a two-sphere. It is
Lagrangian for the standard K\"ahler form; multiplication
by $\mathrm i$ identifies its tangent bundle with
its normal bundle with the opposite oriented Euler
class. Its self-intersection is therefore
$-\chi(S^2)=-2$. Equivalently this is the case
$a=2$, $b=c=1$ of the explicit cycle in A.5.15,
after the linear change $Y^2+Z^2=(Y+\mathrm iZ)(Y-\mathrm iZ)$.
The transported sphere has the same square in the
actual torus, whose $\pi_2$ is zero. This is the
required contradiction. No suspension representation
of the original isolated singularity has been used.

We prove the comparison of Milnor fibres needed for
the argument. Factor the analytic plane function as
\[
                   g=u\prod_{i=1}^r h_i^{b_i},
\]
where the $h_i$ are distinct reduced irreducible
germs and $u$ is a unit. Deform each $h_i$
holomorphically to $h_{i,s}$ in a finite-dimensional
parameter polydisc $P$, and retain its exponent:
\[
 g_s=u\prod_i h_{i,s}^{b_i},\qquad
 H_s=x^a-g_s.                                    \tag{A.J.39}
\]
We use only sufficiently small parameters.

Choose a small closed ball $K\subset\C^2_{y,z}$
such that each $h_i=0$ meets $\partial K$
transversely in its smooth part, and their boundary
links are disjoint. Choose its representative so
that every critical value of $g$ on $K$ is zero.
This is possible because the critical locus of a
holomorphic function has finitely many local
irreducible components, the function is constant
on each of them, and each component of the germ
contains the origin. The boundary smoothness,
transversality and disjointness persist for
$h_{i,s}$ by compactness and the implicit function
theorem. Near any of those links, the other factors
are units, so there are parameter-dependent local
coordinates $(v,w)$ with $g_s=v^{b_i}$, obtained
by taking a root of that unit. The boundary
$\partial K$ is transverse to the $w$ direction.
Consequently the restriction of $g_s$ to
$\partial K$ is a submersion at every point where
$0<|g_s|<\eta_0$, for a uniform sufficiently small
$\eta_0>0$. Away from these link neighbourhoods
$|g_s|$ has a uniform positive lower bound on the
boundary. This proves the asserted uniform
transversality, including at parameters for which
the germs inside $K$ change their singularities.

Take a small disc $D_x$ of radius $\rho$ with
$\rho^a<\eta_0/2$, then shrink $P$ so that every
critical value of $g_s$ on $K$ has modulus less than
$\rho^a/4$. This shrinking follows by compactness:
a sequence of critical points with parameters
tending to zero would otherwise limit to a
nonzero critical value of $g$ in the chosen small
representative. The latter has no such value.
Fix $0<\eta<\rho^a/4$ and put
$W=\overline D_x\times K$, a compact manifold
with corners contained in the original coordinate
chart in $Y$.

Every $H_s^{-1}(\epsilon)$ with
$0<|\epsilon|<\eta$ is transverse to this boundary.
On the face $\partial D_x\times K$, the equation
makes $|g_s|>3\rho^a/4$, so $dg_s\ne0$; its plane
directions give the submersion. On the face
$\overline D_x\times\partial K$, the $x$
derivative gives it when $x\ne0$. At $x=0$ we
have $g_s=-\epsilon\ne0$, and the preceding
boundary calculation gives it. At the corner
$|x|=\rho$, that same calculation for $g_s$
supplies both real target directions. The corner
stratum is therefore also transverse. We use
Ehresmann's theorem with corners: local lifts of
base tangent vectors can be chosen tangent to
each boundary face by these submersion statements,
then patched by a partition of unity and
integrated on the compact fibres. Thus no change
of the chosen domain $W$ is required.

Inside $P\times\{0<|\epsilon|<\eta\}$ remove the
critical-value discriminant $\mathcal D$ of the
maps $H_s$ on $W$. This is a proper complex
analytic subset. The relative critical locus has
no point on the boundary at such a value, so its
projection is proper over every compact subset
of this parameter domain, and the proper mapping
theorem applies. For fixed $s$ its critical
values are discrete, since a holomorphic function
is constant on each connected critical component.
The complement of a proper complex analytic
subset of a connected complex manifold is path
connected. The transversality just proved and
the implicit function theorem give a proper
smooth submersion, including its boundary,
\[
 \{(s,p):p\in W,\ (s,H_s(p))\in
                    (P\times\Delta_\eta^*)\setminus\mathcal D\}
 \longrightarrow
       (P\times\Delta_\eta^*)\setminus\mathcal D,
 \quad (s,p)\longmapsto(s,H_s(p)).                 \tag{A.J.40}
\]
Ehresmann's theorem identifies all its compact
fibres by orientation-preserving diffeomorphisms
of manifolds with corners.
In particular a sufficiently small local Milnor
fibre of an interior singularity of some $H_s$,
at a value near zero, embeds into a smooth fibre
in (A.J.40), hence into
$W\cap H_0^{-1}(\epsilon_0)$ for a sufficiently
small nonzero $\epsilon_0$. This last space is
an actual submanifold of the global torus
$Y_{\epsilon_0^m}$ by (A.J.38); no identification
of it with a standard Milnor fibre is needed.
Choose the value to avoid every other
critical value, and choose its local Milnor ball
in the interior of $W$. This is the precise
comparison we use. Arbitrary deformations which
change the exponents along the boundary links
are not used in (A.J.40).

Suppose first that $r\ge2$. Choose two distinct
factors $h_i,h_j$. Their map
$(h_i,h_j):(\C^2,0)\to(\C^2,0)$ is finite,
since their common zero germ is zero dimensional.
Its positive local degree is their intersection
multiplicity. A sufficiently small regular value
$(s_i,s_j)$ consequently has nonempty finite
inverse image. At each inverse-image point,
$h_i-s_i$ and $h_j-s_j$ are transverse coordinates.
Choose the remaining small shifts to avoid those
finitely many points. The deformation (A.J.39)
then has an interior point at which
\[
            H_s=x^a-U v^{b_i}w^{b_j},
\]
with $U$ a unit independent of $x$. Replacing
$v$ by $vU^{1/b_i}$ makes this exactly
$x^a-v^{b_i}w^{b_j}$ by an invertible holomorphic
coordinate change. A.5.15 constructs a spherical
class of nonzero square in its local Milnor fibre.
The embedding (A.J.40) puts that class in the
actual smooth torus fibre of $f=H_0^m$.
Preservation of intersection
numbers and $\pi_2$-vanishing give the same
contradiction as in A.5.15. Thus $r=1$.

Suppose its reduced factor $h=h_1$ is singular.
Write $g=u h^b$ and absorb a local $b$th root of
$u$ into $h$. This leaves its reduced zero germ
unchanged. Since $h$ is reduced, its gradient
has an isolated zero at the origin: a critical
curve would have constant value zero and would
be a repeated component of $h$. The finite
gradient map has a positive local degree.
Choose a small regular value
$(\lambda_1,\lambda_2)$ of this map. Then
\[
        h_s=h-\lambda_1y-\lambda_2z-\lambda_0
\]
has nondegenerate critical points near zero;
regularity of the gradient value says exactly
that their Hessians are invertible. Choose
$\lambda_0$ so that their critical values are
nonzero. At one such point $p$ put
$v=h_s(p)\ne0$. The Hessian of $h_s^b$ is
$b v^{b-1}$ times the Hessian of $h_s$, and
is invertible. The holomorphic Morse lemma on
the two plane variables identifies the germ of
$H_s-H_s(p)$ with
\[
                         x^a-Y^2-Z^2.
\]
A linear change of $(Y,Z)$ gives the Hirzebruch
germ $x^a-YZ$, already excluded by A.5.15.
Its Milnor fibre embeds by (A.J.40), producing
the same contradiction. Hence $h$ is smooth.
Taking $y=h u^{1/b}$ as a coordinate proves the
asserted transverse product.

For the Hessian consequence, take a local $m$th
root of the actual map, absorbing its unit.
Its reduced equation has the same nonzero
Hessian rank as every other reduced equation
at the critical point. The holomorphic splitting
lemma splits off a nondegenerate square and
gives $f=(x^2-g(y,z))^m$ exactly. If $g=0$,
the equation would not be reduced, so $g\ne0$.
At a critical point $dg(0)=0$. The result just
proved gives $g=y^b$ with $b\ge2$ in suitable
holomorphic coordinates, which is the constant
transverse $A_{b-1}$ product.
\end{proof}


% A.5.18 from part_A6_tacnode_bounds.tex

\begin{statement}{A.5.18}{Proved}{Every horizontal tacnode image is elliptic}
Let $f:Y\to C_0$ be a proper holomorphic map from a smooth
compact complex threefold with $b_2(Y)=0$, with connected
fibres and smooth general fibre a complex two-torus.
Let $mS$ be a full fibre with irreducible reduced support,
$\eta:T\to S$ its normalization, and $\pi:Z\to T$ a
resolution with an irrational ruling $p:Z\to B$.
Suppose every exceptional curve of $\pi$ is vertical and
the conductor divisor on $Z$ is $2C+V$, where $C$ is a
section of $p$ and $V$ is vertical. Let $D=\eta\pi(C)$.
Suppose the generic plane-curve transversal to $D$ in $S$ is an
ordinary tacnode.

Then $D$ is a smooth elliptic curve, $B=C$ is elliptic,
and $C\to D$ is an unramified double cover. The transversal
type along $D$ is everywhere the same ordinary tacnode.
The normalization $T$ is a smooth minimal ruled surface
over this elliptic curve. Its conductor is twice a
section, and $e(S)=0$. In particular its minimal
resolution is the identity.
In particular a rational image of
the horizontal tacnodal conductor cannot occur. The proof
allows vertical conductor and rational singularities of
$T$; it does not assume that the normalization map is
immersive at the special points under examination.
\end{statement}
\begin{proof}
If the normalization of $D$ has positive genus, A.5.14
and A.5.5 make $D$ smooth elliptic and its transversal
type constant. There are two generic normalization
branches of a tacnode. A vertical component of $Z$
is rational, so it cannot dominate this elliptic curve.
Thus both branches lie on $C$, and $C\to D$ is an
unramified double cover: the constant local equation
$x^2-y^4=0$ has two separately smooth normalization
branches on every sufficiently small disc in $D$.
Riemann--Hurwitz gives $g(C)=g(B)=1$.
It remains to exclude the case $D^\nu=\Sigma\simeq\PP^1$.

The finite map $C\to\Sigma$ has degree at least two,
because $C\simeq B$ has positive genus. Since a generic
tacnode has exactly two normalization branches, this
degree is exactly two. Consequently no vertical
conductor component can also dominate $D$: its generic
points would supply an additional normalization branch.
Write $g=g(C)\ge1$ and let $\mathcal B\subset\Sigma$
be the reduced branch divisor of this double cover.
Riemann--Hurwitz gives
\[
                   \deg\mathcal B=2g+2.             \tag{A.T.1}
\]

We first bound the number of branches of the singular
image curve itself. A.5.8 supplies the descended ruling
$\widetilde p:T\to B$. Identify $C$ with $B$ using the
section, let $\phi:B\to D$ be its actual map to the
image curve, and put
\[
                    \psi=\phi\widetilde p:T\to D.
\]
Every vertical conductor component has constant $\psi$
value. Its image under $\eta$ differs from $D$, as just
proved, and meets $D$ in finitely many points. Thus
there is a finite set $A\subset D$ such that, for each
$s\in S$, either $\psi(\eta^{-1}(s))$ is a singleton
or all its values belong to $A$. Include in $A$ the
constant values on the vertical conductor components
and the values of $\phi$ over their finitely many
intersections with $D$. Fibres over the remaining
points of $D$ have their preimages on $C$, where all
the $\psi$ values equal that same point of $D$.

Choose an embedded finite tree $J\subset D$ containing
$A$, in a triangulation of the compact analytic curve.
It can be constructed by joining these finitely many
points by arcs, subdividing their intersections and
retaining a spanning tree. The quotient
$q:D\to D/J$ is a homotopy equivalence: a contractible
CW subcomplex is a cofibration, and the homotopy
extension property extends its contraction. Since
$q\psi$ is constant on every fibre of the finite
proper quotient $\eta$, there is a continuous map
\[
 a:S\to D/J,\qquad a\eta=q\psi,\qquad a|_D=q.     \tag{A.T.2}
\]
The last equality holds first on the dense regular
locus of $D$ and then everywhere by continuity.

Use the explicit triangulated neighbourhood retraction
$r:N\to S$ of (A.Pi.3). For a sufficiently small disc
$\Delta$ around the fibre value, properness gives
$U=f^{-1}(\Delta)\subset N$. The map $(ar)|_U$ is
surjective on fundamental groups because its restriction
to $D$ is the homotopy equivalence $q$. Perturbing loops
off the real-codimension-two analytic support gives
$\pi_1(U\setminus S)\twoheadrightarrow\pi_1(U)$.
The torus-bundle exact sequence over $\Delta^*$ therefore
gives the actual quotient maps
\[
 \Z^4\rtimes\Z\twoheadrightarrow\pi_1(U)
       \twoheadrightarrow\pi_1(D/J)\simeq\pi_1(D).
                                                        \tag{A.T.3}
\]
The first group is metabelian, so the last is metabelian.

For each singular point $x$ of $D$, let $r_x$ be the
number of its normalization branches and put
$\beta=\sum_x(r_x-1)$. The finite normalization map
identifies exactly those finite sets of points on
$\PP^1$; since it is a compact Hausdorff quotient,
this is also the actual underlying topological quotient.
Adjoining arcs between the points to be identified
gives a CW presentation homotopy equivalent to
$S^2\vee\bigvee^\beta S^1$. Hence
$\pi_1(D)$ is the free group $F_\beta$.
For $\beta\ge2$ this group surjects onto $S_4$ by
sending two free generators to $(1234)$ and $(12)$;
the commutator $A_4$ of that quotient is nonabelian.
Equation (A.T.3) therefore forces $\beta\le1$.
By A.5.14 the total delta invariant $\delta=p_a(D)$
is at most two. If $s$ denotes the number of points
of $\Sigma$ over $\operatorname{Sing}D$, then
\[
 s=\#\operatorname{Sing}D+\beta\le\delta+\beta\le3.
                                                        \tag{A.T.4}
\]

First suppose that $D$ is a local complete intersection.
We retain the exact bundles and jets of A.5.5.
Let $\nu:\Sigma\to D$ be normalization, $I=\mathcal I_D$,
\[
 N=(\nu^*(I/I^2))^*,\qquad
 M=\nu^*(\mathcal O_Y(S)|_D).
\]
Thus $M$ is a degree-zero line bundle on $\Sigma$.
Let $\Xi$ be the curve-normalization
conductor divisor, $c=\deg\Xi=2\delta$, and let
$R\subset N$ and $Q=N/R$ be the saturated radical and
quotient of the generic rank-one Hessian. Write
$q_0=\deg Q$. The adjunction and conductor comparisons
(A.J.1) give
\[
 \deg N=-2+c,\quad \deg R=-2+c-q_0,\quad
 0\ne a\in H^0(\Sigma,M\otimes Q^{-2}),\quad q_0\le0.
                                                        \tag{A.T.5}
\]
Here $a$ denotes the intrinsic quadratic coefficient,
whose zero divisor has degree $-2q_0$.
The actual quartic invariant $\Delta=ae-b^2/4$ of
A.5.5 is meromorphic in
$L^2$, where $L=M\otimes Q^{-1}\otimes R^{-2}$.
Its pole bound (A.J.9) gives a nonzero holomorphic
section $\sigma$ of the actual square line bundle
$L^2(4\Xi)=(L(2\Xi))^2$. Its zero divisor has degree
\[
               d_\Delta=8+2q_0.                    \tag{A.T.6}
\]

The double cover $C\to\Sigma$ is obtained by adjoining
$\sqrt{-\Delta}$ to the function field, with the line
trivializations specified by $L$. Indeed over the
generic point a tacnode has normal expansion
$a\ell^2+b\ell r^2+er^4$, and the two normalization
branches are distinguished by the two roots of
$a(\ell/r^2)^2+b(\ell/r^2)+e$. Their discriminant
is $-4\Delta$. Changing a local frame of $L$ multiplies
the representing function of $\Delta$ by a square,
so this describes the same degree-two extension on
every overlap. At a point with local parameter $z$,
write that representing function as $z^n u(z)$,
$u(0)\ne0$. A root of the unit and removal of the
even power of $z$ show that its normalization is
ramified exactly when $n$ is odd. Adding $4\Xi$ changes
each order by an even number. Thus $\mathcal B$ is
exactly the odd-multiplicity support of
$\operatorname{div}(\sigma)$, and
\[
                  2g+2\le 8+2q_0.                 \tag{A.T.7}
\]

At a point of $D_{\rm reg}$ where $a\ne0$, A.5.17
applies to the actual local $m$th root of the full-fibre
map. Its nonzero Hessian makes the local equation a
constant transverse product, necessarily of type
$A_3$ because this is the generic type along $D$.
The conductor cover is therefore unramified there.
Every branch point outside $\nu^{-1}(\operatorname{Sing}D)$
is consequently a zero of $a$. Equations (A.T.4)--(A.T.5)
give the second bound
\[
                    2g+2\le -2q_0+s.              \tag{A.T.8}
\]
Adding (A.T.7) and (A.T.8) yields
$4g+4\le8+s\le11$. Since $g\ge1$, we have
\[
           g=1,\qquad \deg\mathcal B=4,
                    q_0\in\{-2,-1,0\}.            \tag{A.T.9}
\]

We next prove the exact local obstruction required in
the equality cases. Let a smooth point of $D$ be a
branch point at which $\Delta$ has a simple zero.
The preceding nonzero-Hessian argument gives $a(0)=0$.
In coordinates $D=(x,y)=0$ with curve parameter $z$,
choose the holomorphic saturated Hessian factor as the
first normal coordinate. The quadratic jet is then
$a(z)x^2$. The radical cubic vanishes identically,
so a local reduced equation has the exact expansion
\[
 f_0=a(z)x^2+x^2(xA+yB)+xy^2b(z)+xy^3C+y^4E,
                                                        \tag{A.T.10}
\]
with holomorphic coefficients. The identity
$\Delta=ae-b^2/4$, $a(0)=0$ and the simple zero of
$\Delta$ imply $b(0)=0$, $a'(0)\ne0$ and $e(0)\ne0$.
Indeed an order at least two for $a$, together with
$b(0)=0$, would give an order at least two for $\Delta$;
and then $\Delta'(0)=a'(0)e(0)$ is nonzero.
Thus $b/a$ is holomorphic and $E(0)\ne0$.

Set $X=x+b(z)y^2/(2a(z))$. In this actual analytic
coordinate change the coefficient of $Xy^2$ cancels.
All terms created from $x^2(xA+yB)$ lie in
$(X^2)+(Xy^3)+(y^4)$. Therefore, collecting the
terms of $X$-degree zero and one, gives exactly
\[
 f_0=X^2A_1(X,y,z)+Xy^3B_1(y,z)+y^4C_1(y,z),
 \qquad C_1(0)\ne0.
\]
Choose a fourth root of $C_1$ and define
\[
 Y=yC_1^{1/4}+
                  \frac{XB_1}{4C_1^{3/4}}.
                                                        \tag{A.T.11}
\]
The map $(X,y,z)\mapsto(X,Y,z)$ is locally
biholomorphic, and direct expansion gives
$f_0-Y^4\in(X^2)$. Hence $f_0=Y^4+X^2H(X,Y,z)$.
Its quadratic jet along $D$ is still $a(z)X^2$,
so $H(0,0,z)=a(z)$ and $\partial H/\partial z(0)\ne0$.
The further coordinate $Z=H$ gives exactly
\[
                         f_0=Y^4+X^2Z.             \tag{A.T.12}
\]
To apply this to the actual map, first take a local
$m$th root of its unit. This multiplies the reduced
equation by a unit, which changes neither zero order
of $a$ nor of $\Delta$. The same coordinate construction
then makes the full map $(Y^4+X^2Z)^m$.
Statement A.5.15 excludes it. Consequently every
smooth branch point has order at least three in
$\operatorname{div}(\sigma)$.

We also record precisely when the conductor denominator
can be removed at a singular point $x\in D$.
If the intrinsic section $a$ is nonzero at every
normalization point above $x$, choose an actual
constant linear combination of the two lci conormal
generators which is nonzero under the normalized
quotient at all those finitely many points. The
corresponding diagonal entry $a_0$ of the symmetric
Hessian matrix belongs to $\mathcal O_{D,x}$ and
has nonzero value on every normalization branch.
It is therefore a unit in $\mathcal O_{D,x}$.
The other normalized factor coefficient is the
actual quotient $h=q_{XY}/a_0\in\mathcal O_{D,x}$.
Thus square completion and monic polynomial division
in the proof of (A.J.7)--(A.J.8) both take place
over $\mathcal O_{D,x}$ itself. The regular-sequence
graded pieces inject into their normalization
pullbacks, so their cancellations already descend
on $D$. We may repeat that proof with $d=1$.
It gives a holomorphic section $\Delta$ of $L^2$
above $x$, without a conductor pole. At a branch
point above $x$, this proves the sharper inequality
\[
 \operatorname{ord}\sigma\ge
                        4\operatorname{ord}\Xi+1.
                                                        \tag{A.T.13}
\]

We now exhaust the three integer cases in (A.T.9).
If $q_0=-2$, the divisor of $\sigma$ has degree four
and has four odd-multiplicity points, so each is
simple. At most three lie above singularities of
$D$, by (A.T.4). The remaining smooth branch point
contradicts (A.T.12).
If $q_0=0$, the section $a$ has no zeros and all
four branch points would have to lie above the
singular locus; this contradicts $s\le3$.

Finally let $q_0=-1$, so that $\deg\operatorname{div}(a)=2$
and $d_\Delta=6$. If $k$ of the four branch points
lie above $\operatorname{Sing}D$, then the local
obstruction and parity give
\[
               6\ge k+3(4-k)=12-2k.
\]
Hence $k=s=3$, and equality holds: the three singular
preimages have order one in $\operatorname{div}(\sigma)$
and the one smooth branch point has order three.
The relations $s\le\delta+\beta\le3$ now force
$\delta=2$, $\beta=1$ and two singular points of
delta invariant one. The exact delta-one
classification in A.5.14 identifies them as one
ordinary node and one cusp. Their normalization
conductor orders are respectively $(1,1)$ and $2$;
these follow directly from the displayed node and
cusp rings in that proof.

At the cusp preimage, if $a$ were nonzero,
(A.T.13) would make the order of $\sigma$ at least
$4\cdot2+1=9$, exceeding its total degree six.
So $a$ vanishes there. It also vanishes at the
smooth branch point. These two distinct points
exhaust the degree-two divisor of $a$. The section
is consequently nonzero at both node preimages.
Apply (A.T.13) at the node: each of its two branch
points has order at least $4\cdot1+1=5$ in
$\operatorname{div}(\sigma)$, again exceeding the
total degree six. This excludes every rational lci image.

We finish by treating all non-lci image germs permitted
by $\delta\le2$ and $\beta\le1$, with explicit analytic
rings and transverse coordinate maps. For a unibranch
reduced curve ring $A\subset\C\{t\}$ of colength at
most two, an element of order one makes $A=\C\{t\}$.
Otherwise its least positive valuation is two or three:
a least valuation at least four leaves three linearly
independent missing coefficients. If it is three and
the colength is two, the inclusion
$A\subset\C+t^3\C\{t\}$ is equality, since the latter
already has colength two. If it is two, choose a
parameter with $t^2\in A$ by taking a square root of
the leading unit. The even part of every convergent
series belongs to $\C\{t^2\}\subset A$. Subtracting
it shows that the least odd element can be written
$t^{2k+1}v(t^2)$; division by its unit puts
$t^{2k+1}$ in $A$. Consequently
$A=\C\{t^2\}\oplus t^{2k+1}\C\{t^2\}$, of colength
$k\le2$. These are the planar cusp and the planar
$A_4$ curve, besides the smooth case.

For two branches, write their image rings as $A_1,A_2$.
The two projections of $A\subset A_1\oplus A_2$ are
surjective. Their kernels define a common finite
quotient $Q_0$ and an exact fibre-product description
$A=A_1\times_{Q_0}A_2$. Thus
$\delta=\delta_1+\delta_2+\operatorname{length}Q_0$.
If both branches are smooth, the quotient is
$\C\{t\}/(t^k)$ with $k\le2$; after changing the
second parameter, the two maps identify the same
$k$-jet. The map
\[
 x\longmapsto(t,s),\qquad
 y\longmapsto(t^k,-s^k)
\]
identifies this fibre product with
$\C\{x,y\}/(y^2-x^{2k})$: write a matching pair as
the sum of its average and $y$ times its difference
divided by $2x^k$. These are the node and tacnode.
If one branch is singular, the delta budget forces
one cusp, one smooth branch, and $Q_0=\C$.
Three branches are excluded by $\beta\le1$.

An abstract plane curve ring in the given smooth
threefold is lci by the actual transverse-coordinate
construction in A.5.14. Thus the only remaining
non-lci analytic rings, necessarily consuming the
whole delta budget at a single point, are
\[
\begin{array}{ll}
 A=\C\{t^3,t^4,t^5\},
 &I=(y^2-xz,\ yz-x^3,\ z^2-x^2y),\\[2pt]
 A=\C\{t^2,t^3\}\times_{\C}\C\{s\},
 &I=(xz,\ yz,\ y^2-x^3).
\end{array}                                             \tag{A.T.14}
\]
These are normal forms of the actual embeddings.
Their three displayed coordinate functions form a
basis of the cotangent space of each ring. Lift them
to $\mathcal O_{Y,x}$; their differentials are linearly
independent because the embedding dimension is three.
They are therefore ambient coordinates by the inverse
function theorem. The displayed ideals are their
kernels, obtained respectively by the parametrizations
$(t^3,t^4,t^5)$ and the union of $(t^2,t^3,0)$ with
$(0,0,s)$.

We show that the local reduced equation $f_0$ actually
lies in $I^2$ in each case. In the first case put
\[
 g_1=y^2-xz,\quad g_2=yz-x^3,\quad
 g_3=z^2-x^2y,\quad
 F=x^5+xy^3+z^3-3x^2yz.
\]
The exact identities
\[
 xF=g_2^2-g_1g_3,\qquad
 yF=xg_1^2+g_2g_3,\qquad
 zF=g_3^2+xg_1g_2                              \tag{A.T.15}
\]
show that the maximal ideal annihilates the class of
$F$ modulo $I^2$. Moreover
$I^{(2)}=I^2+(F)$ in the convergent local ring.
Here is a finite-module verification of that equality.
Let $J=I^2+(F)$. The quotient $\C\{x,y,z\}/J$ is
finite over the discrete valuation ring $\C\{x\}$,
since its reduction is the finite parametrized curve.
Its generic rank is nine: the curve has degree three
over the $x$ coordinate, and the square of its smooth
codimension-two ideal has length three along each
generic normal slice. Modulo $x$, the quotient has
the nine-element basis
\[
 1,y,y^2,y^3,\ z,yz,y^2z,\ z^2,yz^2,
\]
because its ideal is $(y^4,y^3z,y^2z^2,z^3)$.
A finite module over a discrete valuation ring is
the sum of its free and torsion parts; equality of
generic rank and residue-space dimension here makes
the torsion part zero. Thus $J$ is $x$-saturated.
By (A.T.15), $J$ and $I^2$ agree after inverting $x$.
The curve is lci away from the origin, so its symbolic
square differs from its ordinary square only by a
module supported there. Hence
$I^{(2)}=I^2:x^\infty=J$.

On its normalization let
$a_1=dg_1$ and $b_1=dg_2+t\,dg_1$.
They form a basis of the normalized conormal image;
$dg_3=-2t^2a_1+t b_1$.
The quadratic term of $F$ in this basis is
\[
 q_F=\frac3t a_1^2-\frac3{t^2}a_1b_1
                           +\frac1{t^3}b_1^2,
 \qquad \det(q_F)=\frac3{4t^4}.                  \tag{A.T.16}
\]
Every equation in $I^2$ has a holomorphic quadratic
form in this same basis. If $f_0=cF+I^2$ had
$c(0)\ne0$, the determinant of its quadratic form
would have the noncancellable leading pole
$3c(0)^2/(4t^4)$: the cross terms with a holomorphic
form have pole at most three. It could not have
generic rank one. Thus $c(0)=0$, and (A.T.15) gives
$f_0\in I^2$.

For the second ring put $G=y^2-x^3$ and
$P=(z,G)$, $Q_1=(x,y)$. Then $I=P\cap Q_1$ and
\[
 I^{(2)}=P^2\cap Q_1^2=I^2+(zG).                \tag{A.T.17}
\]
Indeed $P^2=(z^2,zG,G^2)$, while $zG,G^2$ already
belong to $Q_1^2$. Multiplication by $z^2$ is
injective modulo $Q_1^2$, so membership of
$z^2A+zGB+G^2C$ in $Q_1^2$ forces $A\in Q_1^2$.
This gives the asserted generating ideal directly.
Furthermore
$x(zG)=(xz)G$, $y(zG)=(yz)G$, and
$z(zG)=(yz)^2-x(xz)^2$ all lie in $I^2$.
On the cusp branch use $a_1=d(xz)=t^2dz$ and
$b_1=dG=-3t^4dx+2t^3dy$.
The quadratic form of $zG$ is $t^{-2}a_1b_1$,
of determinant $-1/(4t^4)$. The same rank-one
argument forces its coefficient to vanish at the
origin, and the three identities again give
$f_0\in I^2$.

For clarity, the bundle used at these non-lci points
is now defined by its actual morphism. Put
\[
 K=\operatorname{im}\bigl(\nu^*(I/I^2)
                         \longrightarrow\nu^*\Omega_Y^1\bigr),
 \qquad N=K^*.
\]
The coherent image $K$ is locally free of rank two
on the smooth curve $\Sigma$. Let $R_{\rm ram}$ be
the zero divisor of $d\nu$, and let $K_{\rm sat}$
be the kernel of the surjection
$\nu^*\Omega_Y^1\to\Omega_\Sigma^1(-R_{\rm ram})$.
The torsion quotient $K_{\rm sat}/K$ defines a divisor
$J_{\rm jac}$. Taking determinants in these actual
maps gives
\[
 \deg N=-2+\deg(J_{\rm jac}-R_{\rm ram}).         \tag{A.T.18}
\]
For $(t^3,t^4,t^5)$, the wedge $a_1\wedge b_1$
has coefficients $t^8(3,4t,5t^2)$, while $d\nu$
has common factor $t^2$. Thus the correction is
$8-2=6$. For the cusp branch in the second case,
the wedge is $-t^5(2,3t,0)$ and the tangent common
factor is $t$, giving correction four. On its
smooth branch the wedge is $s^2(dx\wedge dy)$,
giving correction two. In both cases (A.T.18)
therefore gives $\deg N=4$.

Because $f_0\in I^2$, its quadratic form is a
holomorphic section of $\operatorname{Sym}^2K\otimes M$,
independent of its expression as a sum of products
of ideal generators. This follows either from its
actual Hessian or from injectivity of
$\operatorname{Sym}^2K\to\nu^*\operatorname{Sym}^2\Omega_Y^1$.
Its saturated radical and quotient again give line
bundles $R,Q$, with $q_0=\deg Q\le0$ and
$\deg R=4-q_0$. The generic quartic invariant
is meromorphic in $L^2$, where
$L=M\otimes Q^{-1}\otimes R^{-2}$; now
\[
                        \deg L^2=-16+2q_0.        \tag{A.T.19}
\]

We bound its poles using actual transverse affine
coordinate maps, with the curve parameter held
fixed during every normal differentiation. For the
space cusp use, for $t\ne0$,
\[
 (u,v,t)\longmapsto
 \left(t^3,\ t^4+\frac{v}{3t^5},\
                   t^5-\frac{u}{t^3}+\frac{2v}{3t^4}\right).
                                                        \tag{A.T.20}
\]
This is locally biholomorphic and its two transverse
derivatives are the dual frame to $a_1,b_1$.
Each coefficient of a derivative has pole at most
five. The first normal derivatives of all three
$g_i$ are holomorphic, since
$dg_1=a_1$, $dg_2=b_1-ta_1$ and
$dg_3=-2t^2a_1+t b_1$.
Second normal derivatives have pole at most ten.
For $f_0=\sum c_{ij}g_i g_j$, the cubic Taylor
coefficients consequently have pole at most ten:
the only surviving terms use a first and a second
derivative of the $g_i$, or a first derivative of
$c_{ij}$ and two first derivatives of the $g_i$.
Every fourth derivative of a holomorphic ambient
function in (A.T.20) has pole at most twenty.

Choose a holomorphic change of the $N$ frame adapted
to its saturated radical and quotient. It does not
increase these bounds. The radical cubic is zero
because the generic transversal is $A_3$.
Writing the relevant Taylor terms as
$a\ell^2+b\ell r^2+er^4$ therefore gives a
holomorphic $a$, a mixed cubic $b$ of pole at most
ten, and $e$ of pole at most twenty. The invariant
$\Delta=ae-b^2/4$ has pole at most twenty.
These are the same intrinsic jets on overlaps:
to weighted order four, where $\ell$ has weight
two and $r$ has weight one, an adapted change is
$\ell'=\alpha\ell+u r^2$, $r'=\beta r$.
Direct substitution multiplies $\Delta$ by
$\alpha^{-2}\beta^{-4}$; changes of the defining
section supply the factor $M^2$. Higher weighted
terms and changes of the curve parameter do not
alter these coefficients. This verifies the
transition law for the line $L^2$ in (A.T.19).

There is only one singular preimage $p$ in this case.
Thus $\Delta$ is a holomorphic section of the square
line $L^2(20p)$, whose degree is $4+2q_0$.
As in (A.T.7), the even pole correction preserves
branch parity, so
\[
 \deg\mathcal B\le4+2q_0,
 \qquad \deg\mathcal B\le-2q_0+1.
\]
The second bound again uses A.5.17 away from that
single singular image preimage. Adding them gives
$2\deg\mathcal B\le5$, contradicting
$\deg\mathcal B=2g+2\ge4$.

For the cusp-plus-line germ, the corresponding
actual affine coordinate maps are
\[
 (u,v,t)\longmapsto
       \left(t^2,t^3+\frac{v}{2t^3},\frac{u}{t^2}\right),
 \qquad
 (u,v,s)\longmapsto\left(\frac us,\frac vs,s\right).
                                                        \tag{A.T.21}
\]
Their first generator derivatives are holomorphic
in the respective normalized conormal frames.
The same product differentiation gives cubic pole
bounds six and two, and the fourth derivative bounds
are twelve and four. Thus $\Delta$ has poles bounded
by $12p+4q$ at the two normalization preimages.
The line $L^2(12p+4q)$ is a square and has degree
$2q_0\le0$ by (A.T.19). A nonzero section has no
positive-degree odd divisor, whereas this double
cover has at least four branch points. This is the
last contradiction. All non-lci cases have now
been excluded, and the positive-genus case at the
start gives the asserted elliptic curves and constant
transversal type.

We finally prove the normalization and Euler assertions,
keeping track of any unnecessary blowups in the original
choice of resolution. Factor $\pi$ through the minimal
resolution $\pi_{\min}:Z_{\min}\to T$, with the resulting
sequence of exceptional-curve blowdowns
$q:Z\to Z_{\min}$. Every curve contracted by $q$ is
vertical for $p$, by the stated assumption on $\pi$.
The equality of structure-sheaf direct images for each
blowdown, in local coordinates on $B$, gives a descended
ruling $p_{\min}:Z_{\min}\to B$ with
$p=p_{\min}q$. The image $C_{\min}$ is still a section,
and the pushed-forward effective conductor is
$D_{\min}=2C_{\min}+V_{\min}$ with $V_{\min}$ vertical.
The normalization-canonical comparison gives
$K_{Z_{\min}}+D_{\min}=\pi_{\min}^*\eta^*\omega_S$
as line bundles, and its rational Chern class is zero.

Statement C2d gives a smooth normalization near the
inverse image of the elliptic curve $D$, with normal
line of degree zero on its unramified double cover.
The minimal resolution is an isomorphism there, so
$C_{\min}^2=0$. No other conductor or exceptional
curve meets it. Hence $V_{\min}\cdot C_{\min}=0$.
No ruling $(-1)$-curve can meet $C_{\min}$: one in
$V_{\min}$ is disjoint from it, while one outside
$V_{\min}$ would have
\[
 (K_{Z_{\min}}+2C_{\min}+V_{\min})\cdot E
       =-1+2C_{\min}\cdot E+V_{\min}\cdot E\ge1,
\]
contrary to the vanishing Chern class.

Contract ruling $(-1)$-curves to a minimal ruled bundle.
Every contraction therefore avoids the section. For
one blowdown $\sigma:W\to W'$ let $D'=\sigma_*D_W$,
let $a$ be the coefficient of its exceptional curve,
and let $d$ be the multiplicity of $D'$ at the centre.
The actual divisor comparison is
\[
 K_W+D_W=\sigma^*(K_{W'}+D')+(1+a-d)E.
                                                        \tag{A.T.22}
\]
Its intersection with $E$ gives $a=d-1$, so the
log-canonical line and its zero rational Chern class
descend. The section remains disjoint from the
effective vertical part at every step. On a minimal
ruled bundle every nonzero effective vertical divisor
is a positive sum of fibres and meets the section.
The final vertical part is consequently zero. If
there had been a blowup in the reverse sequence,
its first centre would be away from this section,
and (A.T.22) would give coefficient $-1$ on its
exceptional curve. Its strict transform retains that
negative coefficient under subsequent blowups, a
contradiction to effectiveness. Thus no ruling
blowdowns were needed: $Z_{\min}$ is already a
minimal ruled bundle and $D_{\min}=2C_{\min}$.

All exceptional curves of $\pi_{\min}$ would be
vertical. A minimal ruled bundle has irreducible
fibres of square zero, which cannot form the
negative-definite exceptional divisor of a surface
resolution. Hence $\pi_{\min}$ has no exceptional
curve and is an isomorphism. This proves the assertion
about $T$. Finally $\eta$ is an isomorphism away from
its conductor section $C_T$ and the elliptic image
$D$. Additivity of the topological Euler number on
these actual decompositions gives
\[
 e(S)=e(T)-e(C_T)+e(D)
      =e(\PP^1)e(B)-0+0=0.                       \tag{A.T.23}
\]
The full multiple fibre has the same underlying
topological space and therefore the same Euler number.
\end{proof}

\paragraph{Connected double covers in the two non-lci cases.}
The non-lci calculation also excludes a connected
degree-two conductor cover when its covering curve
is rational, for either of the two image rings in
(A.T.14). In the cusp-plus-line case the cleared
quartic has degree $2q_0\le0$, so it has no branch
divisor at all. In the space-cusp case a rational
double cover has branch degree two. The already
proved inequalities give
\[
                2\le4+2q_0,\qquad 2\le-2q_0+1,
\]
so the integer $q_0$ is $-1$. The cleared quartic
then has degree two and exactly two simple branch
points. Only one point lies above the singular
image germ. The other is a smooth simple branch
point, excluded by the actual coordinate reduction
(A.T.10)--(A.T.12). Thus these two non-lci images
cannot support a connected double cover of any
nonnegative genus in a torus degeneration.

The integer branch budgets, the discriminant square
weights and the coordinate divisibility used in
(A.T.10)--(A.T.12) are checked in
\path{code/E_tacnode_bounds.py}; the non-lci charts,
symbolic generators and cubic/quartic pole orders
are checked in \path{code/E_low_delta_tacnode.py}.
The topological
quotient (A.T.2) targets the actual singular curve
$D/J$; replacing it by its normalization would
lose the free-group obstruction in (A.T.4).


% A.5.19 from part_A6_triple_images.tex

\begin{statement}{A.5.19}{Proved}{Ordinary-triple images have arithmetic genus at most one}
Let $f:Y\to C_0$ be a proper holomorphic map from a smooth
compact complex threefold with $b_2(Y)=0$, with connected
fibres and smooth general fibre a complex two-torus.
Let $mS$ be a full fibre with irreducible reduced support,
and let $D\subset\operatorname{Sing}S$ be an irreducible
curve with generic ordinary-triple plane transversal.
Then $D$ is a local complete intersection in $Y$ and
$p_a(D)\le1$. The possibilities for its underlying curve
are a smooth rational curve, a rational curve with one
ordinary node or cusp, and a smooth elliptic curve.
In the elliptic case the transversal type is constant.

When $D$ is singular rational, every component of the
normalized conductor curve dominating $D$ is rational,
and its map to $D^\nu\simeq\PP^1$ has degree one.
Consequently, in the irrational ruled setup of A.5.18
with the generic transversal replaced by an ordinary
triple point, the image of the noncontracted horizontal
section is either smooth rational or smooth elliptic.
In the elliptic case the horizontal cover is unramified
of degree three, the normalization is smooth minimal
elliptic ruled, and the fibre has Euler number zero.
In the rational case the horizontal cover has degree
two or three; these cases are not excluded by this
statement.
\end{statement}
\begin{proof}
By A.5.14, $p_a(D)\le2$. If its normalization has
positive genus, that statement and A.5.5 already give
a smooth elliptic curve of constant transverse type.
Suppose its normalization $\nu:\Sigma\to D$ is
$\PP^1$, with total delta invariant $\delta\le2$.

First let $D$ be lci. Put
$N=(\nu^*(I_D/I_D^2))^*$ and
$M=\nu^*(\mathcal O_Y(S)|_D)$.
The quadratic normal jet vanishes generically, so
its vanishing descends on the reduced curve through
the locally free graded piece $I_D^2/I_D^3$.
Thus the initial cubic is an actual section of
$\operatorname{Sym}^3N^*\otimes M$.
Adjunction and the curve conductor give
$\deg N=-2+2\delta$, while $\deg M=0$.
Its nonzero binary-cubic discriminant is a section
of $M^4\otimes(\det N)^{-6}$, of degree
\[
                         12-12\delta.             \tag{A.U.1}
\]
Hence $\delta\le1$.

We exclude every non-lci germ of delta at most two
by its symbolic cube. The one- and two-branch analytic
classification is proved in (A.T.14). There is one
additional possibility when three branches occur:
delta at least two is forced by identifying their
constants, so equality makes the ring exactly
\[
 \{(a,b,c)\in\C\{t_1\}\oplus\C\{t_2\}\oplus\C\{t_3\}:
                      a(0)=b(0)=c(0)\}.
\]
Its actual embedding in the smooth threefold is the
union of the three coordinate axes, with ideal
$I=(xy,xz,yz)$. Indeed the three branch parameters
give a basis of the cotangent space, so their lifts
are ambient coordinates as in A.5.18.
Thus there are exactly three non-lci cases to consider;
each consumes the whole delta budget at one point.

For the space cusp retain $I=(g_1,g_2,g_3)$ and
$F$ from (A.T.15). Its exact symbolic-cube identity is
\[
                         I^{(3)}=I^3+FI.           \tag{A.U.2}
\]
Here is a verification in the convergent local ring.
The right-hand side agrees with $I^3$ after inverting
$x$, by (A.T.15). Its quotient is finite over
$\C\{x\}$, of generic rank $3\cdot6=18$: a square
normal plane has six monomials of total degree less
than three. Modulo $x$, its ideal in $\C\{y,z\}$ is
\[
 (y,z)^6+(y^2z^3,yz^4,z^5).
\]
The fifteen monomials of total degree at most four,
together with $y^5,y^4z,y^3z^2$, give residue dimension
eighteen. The same discrete-valuation-module argument
used for (A.T.15) makes the quotient $x$-torsion-free.
It is therefore the saturation of $I^3$, which equals
the symbolic cube since the curve is lci off the origin.
This proves (A.U.2).

For the cusp-plus-line ideal $I=(xz,yz,G)$ with
$G=y^2-x^3$, one has
\[
                   I^{(3)}=I^3+(zG)I.             \tag{A.U.3}
\]
To verify this directly, write $P=(z,G)$ and
$Q=(x,y)$. Their powers are primary because these
prime ideals are complete intersections. The
symbolic cube is $P^3\cap Q^3$.
Write an element of $P^3$ as
$z^3A+z^2GB+zG^2C+G^3E$.
The last two terms are already in $Q^3$.
Reducing the other terms modulo $Q^3$ gives
$zA+GB=0$ there. The constant-in-$(x,y)$ coefficient
of $B$ must therefore be divisible by $z$; write
$B=zB'+B_Q$, $B_Q\in Q$.
It follows that $A+GB'\in Q^3$.
Consequently
\[
 P^3\cap Q^3=(z^3Q^3,z^2GQ,zG^2,G^3)
                        =I^3+(zG)I,
\]
as claimed.

For the three axes, membership of
$x^a y^b z^c$ in the symbolic cube means
$a+b,a+c,b+c\ge3$.
If one exponent is zero, the other two are at
least three and the monomial belongs to $I^3$.
If all are positive, either at least two are at
least two, in which case it belongs to $(xyz)I$,
or two are one, which violates one of the three
inequalities. Hence
\[
                         I^{(3)}=I^3+(xyz)I.       \tag{A.U.4}
\]

In all three cases the generic ordinary-triple
equation $f_0$ belongs to $I^{(3)}$: on the smooth
locus of each curve branch its order in the two
normal variables is at least three. Define the
bundle by its actual conormal image,
\[
 K=\operatorname{im}\bigl(\nu^*(I/I^2)
                           \to\nu^*\Omega_Y^1\bigr),
 \qquad N=K^*.
\]
For the first two cases (A.T.18) gives $\deg N=4$.
For the three axes, on the branch $(t,0,0)$ the
frame is $t\,dy,t\,dz$, with conormal determinant
order two and no tangent ramification. Each of the
three branches has the same correction, for total
six, so again $\deg N=-2+6=4$.

Use the actual affine maps (A.T.20)--(A.T.21).
For the three axes add
\[
              (u,v,t)\longmapsto(t,u/t,v/t)
                                                        \tag{A.U.5}
\]
on the first branch, and the two analogous permuted
maps on the other branches. In these charts the
first normal derivatives of the ideal generators
are holomorphic in the specified conormal frame.
An $I^3$ term consequently has holomorphic initial
cubic. The extra term in each of (A.U.2)--(A.U.4)
has initial cubic equal to the quadratic term of
its displayed symbolic-square generator, multiplied
by the first derivative of an ideal generator.

For the space cusp that quadratic term is (A.T.16),
so the cubic pole is at most three. For the
cusp-plus-line germ it is $t^{-2}uv$ on the cusp
and $s^{-1}v^2$ on the smooth branch, so the pole
bounds are two and one. For the three axes it is
$t^{-1}uv$ in (A.U.5), giving pole bound one on
each branch. Coefficients multiplying these terms
are holomorphic; their derivatives do not enter
the initial cubic because the lower normal jets
of the products vanish.

The initial cubic is intrinsic on the dense smooth
part of $D$, as all lower jets are zero. Its
transition under a normal-coordinate change is
therefore the actual cubic tensor in
$\operatorname{Sym}^3N^*\otimes M$; higher-order
coordinate terms make no contribution. The binary
discriminant is homogeneous of degree four in
these coefficients. It is a meromorphic section
of $M^4\otimes(\det N)^{-6}$, of degree $-24$,
with pole divisor of total degree at most twelve
in each of the three cases. After this exact pole
correction its line has negative degree, at most
$-12$. It cannot have the required nonzero
holomorphic section. This excludes all non-lci
cases and proves that $D$ is lci with $p_a(D)\le1$.
The list of underlying curves follows from the
delta-one classification in A.5.14.

We give the finite-cover comparison at a singular
rational image. Here $\delta=1$, so (A.U.1) has
degree zero and its nonzero section is nowhere zero
on $\Sigma=\PP^1$. The cubic defines a subscheme
$R_3\subset\PP(N)$ of tangent-root lines. Every
fibre consists of three distinct points. In a
projective coordinate chart, its derivative with
respect to the root coordinate is nonzero, so the
implicit function theorem makes
$R_3\to\Sigma$ a finite etale map of degree three.
There is no component in a whole projective fibre,
since its nonzero discriminant forbids an identically
zero cubic there.

Over the generic point of $D$, normalization of an
ordinary triple transversal separates exactly its
three tangent lines. This supplies an isomorphism
between the function-field algebra of $R_3$ and
that of the normalized conductor curves dominating
$D$, after their maps are factored through $\Sigma$.
Both are normal finite curves over $\Sigma$.
Taking the integral closure of $\mathcal O_\Sigma$
in this same algebra therefore identifies the two
finite covers everywhere. Every connected finite
etale cover of the simply connected Riemann sphere
has degree one. Thus each of these conductor
components is rational and maps isomorphically to
$\Sigma$, proving the second assertion.

In the irrational ruled setting the horizontal
section has positive genus. Its finite map to a
rational image normalization must have degree at
least two. The preceding comparison excludes a
singular rational image, and the generic triple
transversal bounds its degree above by three.
The remaining rational image is therefore smooth,
with degree two or three. If the image is elliptic,
no vertical rational curve can dominate it, so all
three branches belong to the horizontal section.
Statement C2d gives an unramified degree-three
cover, normal degree zero and constant type.
The exact minimal-resolution factorization and
blowdown argument (A.T.22) applies without change
to this coefficient-two section: its proof uses
only those properties. It makes the normalization
smooth minimal elliptic ruled with conductor twice
the section. The actual normalization Euler formula
(A.T.23) then gives zero.
\end{proof}

The singular rational possibilities cannot be
removed by a blanket local spherical-cycle argument.
For example
\[
                     f=x(yz)(x-yz)
                                                        \tag{A.U.6}
\]
has ordinary-triple transversals along the nodal curve
$D=(x,yz)=0$ and an additional ordinary-double
critical line $y=z=0$. For $t\ne0$, put
$B_t=\{(x,v):xv(x-v)=t\}$. The exact maps
\[
 B_t\times\C^*\longrightarrow f^{-1}(t),
 \quad (x,v,y)\longmapsto(x,y,v/y),\qquad
 (x,y,z)\longmapsto(x,yz,y)
                                                        \tag{A.U.7}
\]
are inverse holomorphic maps. Nonzero $t$ forces
$v,y$ and $z$ to be nonzero. The smooth affine
curve $B_t$ and $\C^*$ are aspherical, so this
regular affine fibre has no spherical homotopy
class. This is a local model, not a compact
$b_2=0$ example. A deformation smoothing $yz$
changes the additional critical line's transverse
type on the boundary and does not satisfy the
comparison used in A.5.17.

Exact ideal, normal-coordinate and discriminant
calculations are retained in
\path{code/E_triple_images.py} and
\path{code/E_triple_images.out}; the independently
computed symbolic saturations are also in
\path{code/A_low_delta_cusp.py}.


% A.5.20 from part_A6_cusp_completion.tex

\begin{statement}{A.5.20}{Proved}{A noncontracted horizontal cusp has zero Euler number}
Let $f:Y\to C_0$ be a proper holomorphic map from a smooth
compact complex threefold with $b_2(Y)=0$, with connected
fibres and smooth general fibre a complex two-torus.
Let $mS$ be a full fibre with irreducible reduced support,
$\eta:T\to S$ its normalization and $\pi:Z\to T$ an
irrational ruled resolution with ruling $p:Z\to B$.
Suppose every exceptional curve of $\pi$ is vertical
and the resolution conductor is $2C+V$, where $C$
is a noncontracted section and $V$ is vertical.
If the generic transversal at $D=\eta\pi(C)$ is an
ordinary cusp, then $D$ and $C$ are smooth elliptic
curves and $C\to D$ is an isomorphism. The cusp type
is constant along $D$. The normalization $T$ is a
smooth minimal ruled surface over this elliptic curve,
its conductor is twice a section, and $e(S)=0$.
This permits an initially singular normalization;
smoothness is part of the conclusion.
\end{statement}
\begin{proof}
A transverse cusp has exactly one normalization
branch. The map from the horizontal section to the
normalization of its image consequently has degree
one. Both curves are smooth, so this finite map
identifies $C\simeq B$ with $D^\nu$. Its genus is
positive. Statement A.5.14, including its actual
delta-one lci comparison, and A.5.5 therefore make
$D$ itself smooth elliptic and its cusp transversal
constant at every point. Thus $C\to D$ is an
isomorphism, without any assumption about the
ambient Picard pullback.

For the exact local normalization $x=u^3,y=u^2$,
let $C_T\subset T$ be the image of $C$. The
normalization is smooth near this curve. Its normal
line $P=N_{C_T/T}$ has $P^2=R$ and $P^3=Q$
in the degree-zero radical and Hessian quotient
lines of C2d. Hence $\deg P=0$. The normalization
is smooth near $C$, and no other conductor curve
meets it. Factor $\pi$ through its minimal
resolution, descending the ruling across its
vertical exceptional blowdowns as in A.5.18.
On that minimal resolution the section is isomorphic
to $C_T$ in its unchanged smooth neighbourhood. It
therefore has self-intersection zero and is disjoint from
the effective vertical conductor. The canonical
comparison is the actual line-bundle identity
$K+2C+V=(\eta\pi)^*\omega_S$, of zero rational
Chern class.

Apply (A.T.22) to each ruling blowdown. Its
exceptional coefficient is exactly the multiplicity
at its centre minus one. No ruling $(-1)$-curve
can meet $C$, since the intersection with $K+2C+V$
would then be positive. On the minimal ruled
bundle the vertical conductor is zero because it
is disjoint from a section. A first reverse blowup
away from that section would have coefficient
$-1$, so there were no ruling blowups. All
resolution-exceptional curves are vertical, and
the square-zero fibres of the resulting bundle
cannot be exceptional. This identifies the minimal
resolution with $T$ itself and proves the asserted
smoothness and minimality.

Finally the normalization is an isomorphism off
$C$ and $D$. The exact Euler comparison is
$e(S)=e(T)-e(C)+e(D)$; all three terms vanish,
since $T$ is ruled over an elliptic curve and
$C,D$ are elliptic. Thus $e(S)=0$, also for the
underlying topological space of the full fibre $mS$.
\end{proof}


% A.5.21 from part_A6_binary_products.tex

\begin{statement}{A.5.21}{Proved}{A homogeneous binary-product grid is excluded}
Let $f:Y\to C_0$ be a proper holomorphic map from a smooth
complex threefold with smooth nearby fibres complex
two-tori. Let $P(X,V)$ be a homogeneous polynomial of
degree $d\ge2$, a product of distinct complex linear
forms, and suppose $P(1,0)\ne0$.
There is no point at which the actual fibre map has
the local equation
\[
             f=u(x,y,z)\,P(x,yz)^m,
             \qquad m\ge1,                         \tag{A.BP.1}
\]
with $u$ a holomorphic unit. In particular this excludes
$\prod_{i=1}^d(x-a_i yz)$ for distinct constants $a_i$,
and every positive full-fibre multiplicity of that germ.
No Betti-number assumption on $Y$ is required.

The same exclusion holds for the variable-coefficient
equation $f=u h^m$ with
\[
 h=\sum_{i=0}^d a_i(x,y,z)x^{d-i}(yz)^i,
 \qquad P_0(X,V)=\sum_{i=0}^d a_i(0)X^{d-i}V^i,
                                                        \tag{A.BP.4}
\]
provided $P_0$ has $d$ distinct projective roots and
$P_0(1,0)\ne0$. The comparison with its constant
initial form is the explicit analytic family below.
\end{statement}
\begin{proof}
Put $h=P(x,yz)$. It is weighted homogeneous of degree
$2d$ for the positive weights $(2,1,1)$.
On a small simply connected coordinate neighbourhood
take a holomorphic logarithm of the unit and set
\[
 \Phi(x,y,z)=\bigl(xu^{1/(md)},
                  yu^{1/(2md)},zu^{1/(2md)}\bigr).
                                                        \tag{A.BP.2}
\]
Its derivative at the origin is invertible, and direct
homogeneity gives $h\Phi=u^{1/m}h$. Thus
$f=(h\Phi)^m$ exactly. It suffices to construct a
compact spherical cycle of nonzero square in $h=1$,
because weighted homotheties then put it inside an
arbitrarily small actual fibre of (A.BP.1).

The curve $B=\{P(x,v)=1\}\subset\C^2$ is smooth:
Euler's identity $xP_x+vP_v=dP$ precludes a common
zero of its two partial derivatives on that level.
It is also connected. The map
\[
 B\longrightarrow\PP^1\setminus\{P=0\},
                       \quad(x,v)\longmapsto[x:v]
\]
is an etale cover of degree $d$. Locally a projective
representative has nonzero $P$-value, and its $d$
rescalings with $P=1$ give the sheets. Around any
one of the simple projective zeros of $P$, that
value has a simple zero, so its inverse $d$th root
has monodromy a full $d$-cycle. This monodromy is
transitive on the sheets, proving connectedness.

The function $v$ on $B$ has exactly $d$ distinct
simple zeros, since
$P(x,0)=P(1,0)x^d$ and $P_x(x,0)=d/x$ at those
points. Choose an embedded arc $\gamma$ in $B$
joining two of its zeros and avoiding all other
zeros. On the interior choose a continuous argument
of $v(\gamma(t))$ and two functions $A(t),B_0(t)$
with $A B_0=v(\gamma(t))$ and
$|A|=|B_0|=|v(\gamma(t))|^{1/2}$.
The map
\[
 (t,\theta)\longmapsto
   \bigl(x(\gamma(t)),A(t)e^{\mathrm i\theta},
                        B_0(t)e^{-\mathrm i\theta}\bigr)
                                                        \tag{A.BP.3}
\]
lies in $h=1$ and collapses both boundary circles
of $[0,1]\times S^1$ to points. It is therefore a
spherical cycle, with embedded interior circles.

A parallel-arc isotopy with the same endpoints gives
a homologous spherical cycle whose support meets
the first only at those endpoints. At each endpoint,
$P_x\ne0$, so $(y,z)$ are actual holomorphic
coordinates on $h=1$. The links of the two cycles
are parallel curves of slope $(1,-1)$ on a Clifford
torus in the local three-sphere. Their linking
number is $-1$. As in the explicit relative-chain
comparison in A.5.15, that number is the local
intersection of the two coned cycles. Both
orientations reverse at the second endpoint,
so its intersection sign is unchanged. The
spherical class consequently has square $-2$.

The weighted homothety
$(x,y,z)\mapsto(\lambda^2x,\lambda y,\lambda z)$
maps $h=1$ biholomorphically to $h=\lambda^{2d}$
and shrinks the compact cycles and their isotopy
to the origin. Together with (A.BP.2), this puts
them inside an actual smooth global fibre of $f$.
That fibre is a complex two-torus and has zero
$\pi_2$, so the spherical class has zero homology
image. The oriented codimension-zero embedding
preserves the intersection of the compact cycles,
giving both zero and $-2$ for its square. This
contradiction proves the statement.

We now prove the variable-coefficient assertion,
including the boundary comparison. Absorb a local
$m$th root of the unit into $h$; this preserves the
form (A.BP.4) and its two initial-form conditions.
For a complex parameter $s$ set
\[
 h_s=\sum_{i=0}^d a_i(s^2x,sy,sz)x^{d-i}(yz)^i.
                                                        \tag{A.BP.5}
\]
This is holomorphic also at $s=0$, with
$h_0=P_0(x,yz)$. For $s\ne0$ there is the exact
identity
\[
       h(s^2x,sy,sz)=s^{2d}h_s(x,y,z).           \tag{A.BP.6}
\]
The leading coefficient $a_0$ is a unit. The implicit
function theorem applied to its degree-$d$ root
polynomial gives distinct holomorphic functions
$b_i(s,x,y,z)$ and the exact factorization
\[
       h_s=a_0(s^2x,sy,sz)
                     \prod_{i=1}^d(x-b_i(s,x,y,z)yz)
                                                        \tag{A.BP.7}
\]
on one small parameter and coordinate neighbourhood.
Each factor defines a smooth surface, since its
$x$ derivative is a unit. For $i\ne j$ the ideal
generated by the two factors equals $(x,yz)$:
their difference is the unit $b_j-b_i$ times $yz$.
Thus the zero level has precisely the same two
critical curves $D=(x,yz)=0$ for every $s$.

Choose a small closed ball $K$ in the coordinate
chart such that the two axes of $D$ meet its boundary
transversely and the zero-level surfaces away from
those axes are transverse to the boundary. At
$s=0$ this follows directly for the smooth graph
surfaces $x=b_i(0)yz$ and sufficiently small radii.
It persists for small $s$ by compactness away from
neighbourhoods of the two boundary links.
Near either link, use normal coordinates $(U,V)$
for the smooth critical curve and a parameter $w$
along it. The normal initial term of $h_s$ is a
degree-$d$ binary form with distinct roots.
On the compact parameterized link its normal
gradient therefore satisfies a uniform bound
\[
 \|\partial_{U,V}h_s\|\ge c\|(U,V)\|^{d-1},
 \qquad |\partial_w h_s|\le C\|(U,V)\|^d
                                                        \tag{A.BP.8}
\]
after the neighbourhood is shrunk. The first bound
follows because the gradient of a homogeneous
reduced binary form has no zero on the unit sphere;
higher normal terms have one additional order.
The second follows from $h_s\in(U,V)^d$.

Transversality of the critical curve to $\partial K$
allows a boundary coordinate with one real component
of $w$ solved as a smooth function of $(U,V)$ and
the remaining component. The correction to normal
differentiation tangent to this boundary is a
bounded multiple of $\partial_w h_s$.
By (A.BP.8) it is smaller by one normal order than
the surjective complex normal derivative. Thus
$h_s|_{\partial K}$ is a real submersion at every
nearby point with $h_s\ne0$, uniformly in $s$.
Away from the two link neighbourhoods, compactness
and the preceding zero-level transversality give
the same conclusion for all sufficiently small
nonzero values.

Remove the proper critical-value discriminant from
the product of the parameter disc and a sufficiently
small punctured value disc. The critical locus over
that domain cannot meet the boundary, so its
projection is proper and its image is complex
analytic. Its complement is connected. The level
sets in $K$ form a proper submersion with boundary
there, and Ehresmann transport identifies their
compact fibres by oriented diffeomorphisms.
Choose a small sphere cycle of the already proved
model $h_0=P_0(x,yz)$, and transport it to a fibre
of $h_{s_0}$ for $s_0\ne0$. Equation (A.BP.6)
then embeds that fibre into the actual fibre
of $h^m$ with value $s_0^{2dm}\epsilon^m$.
It carries the spherical class of square $-2$
into a global torus. The same intersection and
asphericity contradiction proves (A.BP.4).
\end{proof}

The condition $P(1,0)\ne0$ is used to produce the
distinct collapsing circles. The exact product
model (A.U.6), for which $P(X,V)=XV(X-V)$, has
$P(1,0)=0$ and is outside this statement.
The variable-coefficient extension uses only the
specific comparison (A.BP.5)--(A.BP.8), with its
unchanged boundary critical curves. The weighted scaling,
simple-zero and endpoint-square identities are
checked in \path{code/E_binary_products.py}.

\paragraph{Two exact scalar-zero comparisons.}
The homogeneous germ
\[
                  h=xy\{z(x-y)+xy\}
\tag{A.BP.9}
\]
has a cubic normal tensor along $(x,y)=0$ which
vanishes identically at $z=0$. On $xy\ne0$ put
\[
 v=x-y,\qquad Z=y(y+v)z+y^2(2y+v).
\]
These are actual coordinates with $y$, with inverse
$x=y+v$, $z=(Z-y^2(2y+v))/(y(y+v))$, and
\[
                              h=vZ+y^4.
\tag{A.BP.10}
\]
Choose an arc between two roots of $y^4=1$ avoiding
$y=0$ and the other roots. The matching circles
$v=\rho A e^{\mathrm i\theta}$,
$Z=\rho^{-1}B e^{-\mathrm i\theta}$, where
$AB=1-y^4$ and $|A|=|B|=|1-y^4|^{1/2}$,
collapse at its endpoints. Choose the positive
constant $\rho$ small enough that $|v|<|y|/2$
on the arc. The entire spherical cycle then lies
in the indicated coordinate chart. Its two endpoint
linkings are $-1$, as in (A.BP.3), so its square
is $-2$. Ordinary homothety in $(x,y,z)$ shrinks
it because (A.BP.9) has degree four. The same
weighted unit map as (A.BP.2), with all three
weights equal to one, includes every local unit
and every positive full-fibre multiplicity.
Thus this particular scalar-zero germ is excluded.

The following different scalar-zero model records
the limit of that argument:
\[
                     h=xz(y^2-x^2z).
\tag{A.BP.11}
\]
For $h=t\ne0$, the substitution $u=y/x$ is
invertible and gives $h=x^3 z(u^2-z)$.
The map forgetting $x$ is a finite etale triple
cover of $E=\{z(u^2-z)\ne0\}$.
There is an explicit finite etale double cover
\[
 \C^*\times(\C\setminus\{1,-1\})\longrightarrow E,
                  \qquad(s,v)\longmapsto(s^2,sv)
\tag{A.BP.12}
\]
in coordinates $(z,u)$, with deck involution
$(s,v)\mapsto(-s,-v)$. Its source is aspherical,
so both $E$ and the regular affine fibre of
(A.BP.11) are aspherical by these actual covering
maps. The quadratic surface branch has normalization
$(v,u)\mapsto(v,vu,u^2)$ and conductor
$2(v=0)+2(u=0)$ in the full reduced surface.
The other two normalization branches have conductors
$2(y=0)+(z=0)$ on $x=0$, and
$(x=0)+2(y=0)$ on $z=0$.
Thus their reduced conductor supports have only
normal-crossing pairs. This is an exact local
model, not a construction of a compact $b_2=0$
fibre. It shows why a vanishing cubic coefficient
alone does not justify a local spherical-cycle
exclusion. The identities and the two finite-cover
formulas are checked in \path{code/E_cubic_double.py}.

